\documentclass[11pt]{article}

    \usepackage[breakable]{tcolorbox}
    \usepackage{parskip} % Stop auto-indenting (to mimic markdown behaviour)
    

    % Basic figure setup, for now with no caption control since it's done
    % automatically by Pandoc (which extracts ![](path) syntax from Markdown).
    \usepackage{graphicx}
    % Keep aspect ratio if custom image width or height is specified
    \setkeys{Gin}{keepaspectratio}
    % Maintain compatibility with old templates. Remove in nbconvert 6.0
    \let\Oldincludegraphics\includegraphics
    % Ensure that by default, figures have no caption (until we provide a
    % proper Figure object with a Caption API and a way to capture that
    % in the conversion process - todo).
    \usepackage{caption}
    \DeclareCaptionFormat{nocaption}{}
    \captionsetup{format=nocaption,aboveskip=0pt,belowskip=0pt}

    \usepackage{float}
    \floatplacement{figure}{H} % forces figures to be placed at the correct location
    \usepackage{xcolor} % Allow colors to be defined
    \usepackage{enumerate} % Needed for markdown enumerations to work
    \usepackage{geometry} % Used to adjust the document margins
    \usepackage{amsmath} % Equations
    \usepackage{amssymb} % Equations
    \usepackage{textcomp} % defines textquotesingle
    % Hack from http://tex.stackexchange.com/a/47451/13684:
    \AtBeginDocument{%
        \def\PYZsq{\textquotesingle}% Upright quotes in Pygmentized code
    }
    \usepackage{upquote} % Upright quotes for verbatim code
    \usepackage{eurosym} % defines \euro

    \usepackage{iftex}
    \ifPDFTeX
        \usepackage[T1]{fontenc}
        \IfFileExists{alphabeta.sty}{
              \usepackage{alphabeta}
          }{
              \usepackage[mathletters]{ucs}
              \usepackage[utf8x]{inputenc}
          }
    \else
        \usepackage{fontspec}
        \usepackage{unicode-math}
    \fi

    \usepackage{fancyvrb} % verbatim replacement that allows latex
    \usepackage{grffile} % extends the file name processing of package graphics
                         % to support a larger range
    \makeatletter % fix for old versions of grffile with XeLaTeX
    \@ifpackagelater{grffile}{2019/11/01}
    {
      % Do nothing on new versions
    }
    {
      \def\Gread@@xetex#1{%
        \IfFileExists{"\Gin@base".bb}%
        {\Gread@eps{\Gin@base.bb}}%
        {\Gread@@xetex@aux#1}%
      }
    }
    \makeatother
    \usepackage[Export]{adjustbox} % Used to constrain images to a maximum size
    \adjustboxset{max size={0.9\linewidth}{0.9\paperheight}}

    % The hyperref package gives us a pdf with properly built
    % internal navigation ('pdf bookmarks' for the table of contents,
    % internal cross-reference links, web links for URLs, etc.)
    \usepackage{hyperref}
    % The default LaTeX title has an obnoxious amount of whitespace. By default,
    % titling removes some of it. It also provides customization options.
    \usepackage{titling}
    \usepackage{longtable} % longtable support required by pandoc >1.10
    \usepackage{booktabs}  % table support for pandoc > 1.12.2
    \usepackage{array}     % table support for pandoc >= 2.11.3
    \usepackage{calc}      % table minipage width calculation for pandoc >= 2.11.1
    \usepackage[inline]{enumitem} % IRkernel/repr support (it uses the enumerate* environment)
    \usepackage[normalem]{ulem} % ulem is needed to support strikethroughs (\sout)
                                % normalem makes italics be italics, not underlines
    \usepackage{soul}      % strikethrough (\st) support for pandoc >= 3.0.0
    \usepackage{mathrsfs}
    

    
    % Colors for the hyperref package
    \definecolor{urlcolor}{rgb}{0,.145,.698}
    \definecolor{linkcolor}{rgb}{.71,0.21,0.01}
    \definecolor{citecolor}{rgb}{.12,.54,.11}

    % ANSI colors
    \definecolor{ansi-black}{HTML}{3E424D}
    \definecolor{ansi-black-intense}{HTML}{282C36}
    \definecolor{ansi-red}{HTML}{E75C58}
    \definecolor{ansi-red-intense}{HTML}{B22B31}
    \definecolor{ansi-green}{HTML}{00A250}
    \definecolor{ansi-green-intense}{HTML}{007427}
    \definecolor{ansi-yellow}{HTML}{DDB62B}
    \definecolor{ansi-yellow-intense}{HTML}{B27D12}
    \definecolor{ansi-blue}{HTML}{208FFB}
    \definecolor{ansi-blue-intense}{HTML}{0065CA}
    \definecolor{ansi-magenta}{HTML}{D160C4}
    \definecolor{ansi-magenta-intense}{HTML}{A03196}
    \definecolor{ansi-cyan}{HTML}{60C6C8}
    \definecolor{ansi-cyan-intense}{HTML}{258F8F}
    \definecolor{ansi-white}{HTML}{C5C1B4}
    \definecolor{ansi-white-intense}{HTML}{A1A6B2}
    \definecolor{ansi-default-inverse-fg}{HTML}{FFFFFF}
    \definecolor{ansi-default-inverse-bg}{HTML}{000000}

    % common color for the border for error outputs.
    \definecolor{outerrorbackground}{HTML}{FFDFDF}

    % commands and environments needed by pandoc snippets
    % extracted from the output of `pandoc -s`
    \providecommand{\tightlist}{%
      \setlength{\itemsep}{0pt}\setlength{\parskip}{0pt}}
    \DefineVerbatimEnvironment{Highlighting}{Verbatim}{commandchars=\\\{\}}
    % Add ',fontsize=\small' for more characters per line
    \newenvironment{Shaded}{}{}
    \newcommand{\KeywordTok}[1]{\textcolor[rgb]{0.00,0.44,0.13}{\textbf{{#1}}}}
    \newcommand{\DataTypeTok}[1]{\textcolor[rgb]{0.56,0.13,0.00}{{#1}}}
    \newcommand{\DecValTok}[1]{\textcolor[rgb]{0.25,0.63,0.44}{{#1}}}
    \newcommand{\BaseNTok}[1]{\textcolor[rgb]{0.25,0.63,0.44}{{#1}}}
    \newcommand{\FloatTok}[1]{\textcolor[rgb]{0.25,0.63,0.44}{{#1}}}
    \newcommand{\CharTok}[1]{\textcolor[rgb]{0.25,0.44,0.63}{{#1}}}
    \newcommand{\StringTok}[1]{\textcolor[rgb]{0.25,0.44,0.63}{{#1}}}
    \newcommand{\CommentTok}[1]{\textcolor[rgb]{0.38,0.63,0.69}{\textit{{#1}}}}
    \newcommand{\OtherTok}[1]{\textcolor[rgb]{0.00,0.44,0.13}{{#1}}}
    \newcommand{\AlertTok}[1]{\textcolor[rgb]{1.00,0.00,0.00}{\textbf{{#1}}}}
    \newcommand{\FunctionTok}[1]{\textcolor[rgb]{0.02,0.16,0.49}{{#1}}}
    \newcommand{\RegionMarkerTok}[1]{{#1}}
    \newcommand{\ErrorTok}[1]{\textcolor[rgb]{1.00,0.00,0.00}{\textbf{{#1}}}}
    \newcommand{\NormalTok}[1]{{#1}}

    % Additional commands for more recent versions of Pandoc
    \newcommand{\ConstantTok}[1]{\textcolor[rgb]{0.53,0.00,0.00}{{#1}}}
    \newcommand{\SpecialCharTok}[1]{\textcolor[rgb]{0.25,0.44,0.63}{{#1}}}
    \newcommand{\VerbatimStringTok}[1]{\textcolor[rgb]{0.25,0.44,0.63}{{#1}}}
    \newcommand{\SpecialStringTok}[1]{\textcolor[rgb]{0.73,0.40,0.53}{{#1}}}
    \newcommand{\ImportTok}[1]{{#1}}
    \newcommand{\DocumentationTok}[1]{\textcolor[rgb]{0.73,0.13,0.13}{\textit{{#1}}}}
    \newcommand{\AnnotationTok}[1]{\textcolor[rgb]{0.38,0.63,0.69}{\textbf{\textit{{#1}}}}}
    \newcommand{\CommentVarTok}[1]{\textcolor[rgb]{0.38,0.63,0.69}{\textbf{\textit{{#1}}}}}
    \newcommand{\VariableTok}[1]{\textcolor[rgb]{0.10,0.09,0.49}{{#1}}}
    \newcommand{\ControlFlowTok}[1]{\textcolor[rgb]{0.00,0.44,0.13}{\textbf{{#1}}}}
    \newcommand{\OperatorTok}[1]{\textcolor[rgb]{0.40,0.40,0.40}{{#1}}}
    \newcommand{\BuiltInTok}[1]{{#1}}
    \newcommand{\ExtensionTok}[1]{{#1}}
    \newcommand{\PreprocessorTok}[1]{\textcolor[rgb]{0.74,0.48,0.00}{{#1}}}
    \newcommand{\AttributeTok}[1]{\textcolor[rgb]{0.49,0.56,0.16}{{#1}}}
    \newcommand{\InformationTok}[1]{\textcolor[rgb]{0.38,0.63,0.69}{\textbf{\textit{{#1}}}}}
    \newcommand{\WarningTok}[1]{\textcolor[rgb]{0.38,0.63,0.69}{\textbf{\textit{{#1}}}}}
    \makeatletter
    \newsavebox\pandoc@box
    \newcommand*\pandocbounded[1]{%
      \sbox\pandoc@box{#1}%
      % scaling factors for width and height
      \Gscale@div\@tempa\textheight{\dimexpr\ht\pandoc@box+\dp\pandoc@box\relax}%
      \Gscale@div\@tempb\linewidth{\wd\pandoc@box}%
      % select the smaller of both
      \ifdim\@tempb\p@<\@tempa\p@
        \let\@tempa\@tempb
      \fi
      % scaling accordingly (\@tempa < 1)
      \ifdim\@tempa\p@<\p@
        \scalebox{\@tempa}{\usebox\pandoc@box}%
      % scaling not needed, use as it is
      \else
        \usebox{\pandoc@box}%
      \fi
    }
    \makeatother

    % Define a nice break command that doesn't care if a line doesn't already
    % exist.
    \def\br{\hspace*{\fill} \\* }
    % Math Jax compatibility definitions
    \def\gt{>}
    \def\lt{<}
    \let\Oldtex\TeX
    \let\Oldlatex\LaTeX
    \renewcommand{\TeX}{\textrm{\Oldtex}}
    \renewcommand{\LaTeX}{\textrm{\Oldlatex}}
    % Document parameters
    % Document title
    \title{practica\_01\_jordi\_blasco\_lozano}
    
    
    
    
    
    
    
% Pygments definitions
\makeatletter
\def\PY@reset{\let\PY@it=\relax \let\PY@bf=\relax%
    \let\PY@ul=\relax \let\PY@tc=\relax%
    \let\PY@bc=\relax \let\PY@ff=\relax}
\def\PY@tok#1{\csname PY@tok@#1\endcsname}
\def\PY@toks#1+{\ifx\relax#1\empty\else%
    \PY@tok{#1}\expandafter\PY@toks\fi}
\def\PY@do#1{\PY@bc{\PY@tc{\PY@ul{%
    \PY@it{\PY@bf{\PY@ff{#1}}}}}}}
\def\PY#1#2{\PY@reset\PY@toks#1+\relax+\PY@do{#2}}

\@namedef{PY@tok@w}{\def\PY@tc##1{\textcolor[rgb]{0.73,0.73,0.73}{##1}}}
\@namedef{PY@tok@c}{\let\PY@it=\textit\def\PY@tc##1{\textcolor[rgb]{0.24,0.48,0.48}{##1}}}
\@namedef{PY@tok@cp}{\def\PY@tc##1{\textcolor[rgb]{0.61,0.40,0.00}{##1}}}
\@namedef{PY@tok@k}{\let\PY@bf=\textbf\def\PY@tc##1{\textcolor[rgb]{0.00,0.50,0.00}{##1}}}
\@namedef{PY@tok@kp}{\def\PY@tc##1{\textcolor[rgb]{0.00,0.50,0.00}{##1}}}
\@namedef{PY@tok@kt}{\def\PY@tc##1{\textcolor[rgb]{0.69,0.00,0.25}{##1}}}
\@namedef{PY@tok@o}{\def\PY@tc##1{\textcolor[rgb]{0.40,0.40,0.40}{##1}}}
\@namedef{PY@tok@ow}{\let\PY@bf=\textbf\def\PY@tc##1{\textcolor[rgb]{0.67,0.13,1.00}{##1}}}
\@namedef{PY@tok@nb}{\def\PY@tc##1{\textcolor[rgb]{0.00,0.50,0.00}{##1}}}
\@namedef{PY@tok@nf}{\def\PY@tc##1{\textcolor[rgb]{0.00,0.00,1.00}{##1}}}
\@namedef{PY@tok@nc}{\let\PY@bf=\textbf\def\PY@tc##1{\textcolor[rgb]{0.00,0.00,1.00}{##1}}}
\@namedef{PY@tok@nn}{\let\PY@bf=\textbf\def\PY@tc##1{\textcolor[rgb]{0.00,0.00,1.00}{##1}}}
\@namedef{PY@tok@ne}{\let\PY@bf=\textbf\def\PY@tc##1{\textcolor[rgb]{0.80,0.25,0.22}{##1}}}
\@namedef{PY@tok@nv}{\def\PY@tc##1{\textcolor[rgb]{0.10,0.09,0.49}{##1}}}
\@namedef{PY@tok@no}{\def\PY@tc##1{\textcolor[rgb]{0.53,0.00,0.00}{##1}}}
\@namedef{PY@tok@nl}{\def\PY@tc##1{\textcolor[rgb]{0.46,0.46,0.00}{##1}}}
\@namedef{PY@tok@ni}{\let\PY@bf=\textbf\def\PY@tc##1{\textcolor[rgb]{0.44,0.44,0.44}{##1}}}
\@namedef{PY@tok@na}{\def\PY@tc##1{\textcolor[rgb]{0.41,0.47,0.13}{##1}}}
\@namedef{PY@tok@nt}{\let\PY@bf=\textbf\def\PY@tc##1{\textcolor[rgb]{0.00,0.50,0.00}{##1}}}
\@namedef{PY@tok@nd}{\def\PY@tc##1{\textcolor[rgb]{0.67,0.13,1.00}{##1}}}
\@namedef{PY@tok@s}{\def\PY@tc##1{\textcolor[rgb]{0.73,0.13,0.13}{##1}}}
\@namedef{PY@tok@sd}{\let\PY@it=\textit\def\PY@tc##1{\textcolor[rgb]{0.73,0.13,0.13}{##1}}}
\@namedef{PY@tok@si}{\let\PY@bf=\textbf\def\PY@tc##1{\textcolor[rgb]{0.64,0.35,0.47}{##1}}}
\@namedef{PY@tok@se}{\let\PY@bf=\textbf\def\PY@tc##1{\textcolor[rgb]{0.67,0.36,0.12}{##1}}}
\@namedef{PY@tok@sr}{\def\PY@tc##1{\textcolor[rgb]{0.64,0.35,0.47}{##1}}}
\@namedef{PY@tok@ss}{\def\PY@tc##1{\textcolor[rgb]{0.10,0.09,0.49}{##1}}}
\@namedef{PY@tok@sx}{\def\PY@tc##1{\textcolor[rgb]{0.00,0.50,0.00}{##1}}}
\@namedef{PY@tok@m}{\def\PY@tc##1{\textcolor[rgb]{0.40,0.40,0.40}{##1}}}
\@namedef{PY@tok@gh}{\let\PY@bf=\textbf\def\PY@tc##1{\textcolor[rgb]{0.00,0.00,0.50}{##1}}}
\@namedef{PY@tok@gu}{\let\PY@bf=\textbf\def\PY@tc##1{\textcolor[rgb]{0.50,0.00,0.50}{##1}}}
\@namedef{PY@tok@gd}{\def\PY@tc##1{\textcolor[rgb]{0.63,0.00,0.00}{##1}}}
\@namedef{PY@tok@gi}{\def\PY@tc##1{\textcolor[rgb]{0.00,0.52,0.00}{##1}}}
\@namedef{PY@tok@gr}{\def\PY@tc##1{\textcolor[rgb]{0.89,0.00,0.00}{##1}}}
\@namedef{PY@tok@ge}{\let\PY@it=\textit}
\@namedef{PY@tok@gs}{\let\PY@bf=\textbf}
\@namedef{PY@tok@ges}{\let\PY@bf=\textbf\let\PY@it=\textit}
\@namedef{PY@tok@gp}{\let\PY@bf=\textbf\def\PY@tc##1{\textcolor[rgb]{0.00,0.00,0.50}{##1}}}
\@namedef{PY@tok@go}{\def\PY@tc##1{\textcolor[rgb]{0.44,0.44,0.44}{##1}}}
\@namedef{PY@tok@gt}{\def\PY@tc##1{\textcolor[rgb]{0.00,0.27,0.87}{##1}}}
\@namedef{PY@tok@err}{\def\PY@bc##1{{\setlength{\fboxsep}{\string -\fboxrule}\fcolorbox[rgb]{1.00,0.00,0.00}{1,1,1}{\strut ##1}}}}
\@namedef{PY@tok@kc}{\let\PY@bf=\textbf\def\PY@tc##1{\textcolor[rgb]{0.00,0.50,0.00}{##1}}}
\@namedef{PY@tok@kd}{\let\PY@bf=\textbf\def\PY@tc##1{\textcolor[rgb]{0.00,0.50,0.00}{##1}}}
\@namedef{PY@tok@kn}{\let\PY@bf=\textbf\def\PY@tc##1{\textcolor[rgb]{0.00,0.50,0.00}{##1}}}
\@namedef{PY@tok@kr}{\let\PY@bf=\textbf\def\PY@tc##1{\textcolor[rgb]{0.00,0.50,0.00}{##1}}}
\@namedef{PY@tok@bp}{\def\PY@tc##1{\textcolor[rgb]{0.00,0.50,0.00}{##1}}}
\@namedef{PY@tok@fm}{\def\PY@tc##1{\textcolor[rgb]{0.00,0.00,1.00}{##1}}}
\@namedef{PY@tok@vc}{\def\PY@tc##1{\textcolor[rgb]{0.10,0.09,0.49}{##1}}}
\@namedef{PY@tok@vg}{\def\PY@tc##1{\textcolor[rgb]{0.10,0.09,0.49}{##1}}}
\@namedef{PY@tok@vi}{\def\PY@tc##1{\textcolor[rgb]{0.10,0.09,0.49}{##1}}}
\@namedef{PY@tok@vm}{\def\PY@tc##1{\textcolor[rgb]{0.10,0.09,0.49}{##1}}}
\@namedef{PY@tok@sa}{\def\PY@tc##1{\textcolor[rgb]{0.73,0.13,0.13}{##1}}}
\@namedef{PY@tok@sb}{\def\PY@tc##1{\textcolor[rgb]{0.73,0.13,0.13}{##1}}}
\@namedef{PY@tok@sc}{\def\PY@tc##1{\textcolor[rgb]{0.73,0.13,0.13}{##1}}}
\@namedef{PY@tok@dl}{\def\PY@tc##1{\textcolor[rgb]{0.73,0.13,0.13}{##1}}}
\@namedef{PY@tok@s2}{\def\PY@tc##1{\textcolor[rgb]{0.73,0.13,0.13}{##1}}}
\@namedef{PY@tok@sh}{\def\PY@tc##1{\textcolor[rgb]{0.73,0.13,0.13}{##1}}}
\@namedef{PY@tok@s1}{\def\PY@tc##1{\textcolor[rgb]{0.73,0.13,0.13}{##1}}}
\@namedef{PY@tok@mb}{\def\PY@tc##1{\textcolor[rgb]{0.40,0.40,0.40}{##1}}}
\@namedef{PY@tok@mf}{\def\PY@tc##1{\textcolor[rgb]{0.40,0.40,0.40}{##1}}}
\@namedef{PY@tok@mh}{\def\PY@tc##1{\textcolor[rgb]{0.40,0.40,0.40}{##1}}}
\@namedef{PY@tok@mi}{\def\PY@tc##1{\textcolor[rgb]{0.40,0.40,0.40}{##1}}}
\@namedef{PY@tok@il}{\def\PY@tc##1{\textcolor[rgb]{0.40,0.40,0.40}{##1}}}
\@namedef{PY@tok@mo}{\def\PY@tc##1{\textcolor[rgb]{0.40,0.40,0.40}{##1}}}
\@namedef{PY@tok@ch}{\let\PY@it=\textit\def\PY@tc##1{\textcolor[rgb]{0.24,0.48,0.48}{##1}}}
\@namedef{PY@tok@cm}{\let\PY@it=\textit\def\PY@tc##1{\textcolor[rgb]{0.24,0.48,0.48}{##1}}}
\@namedef{PY@tok@cpf}{\let\PY@it=\textit\def\PY@tc##1{\textcolor[rgb]{0.24,0.48,0.48}{##1}}}
\@namedef{PY@tok@c1}{\let\PY@it=\textit\def\PY@tc##1{\textcolor[rgb]{0.24,0.48,0.48}{##1}}}
\@namedef{PY@tok@cs}{\let\PY@it=\textit\def\PY@tc##1{\textcolor[rgb]{0.24,0.48,0.48}{##1}}}

\def\PYZbs{\char`\\}
\def\PYZus{\char`\_}
\def\PYZob{\char`\{}
\def\PYZcb{\char`\}}
\def\PYZca{\char`\^}
\def\PYZam{\char`\&}
\def\PYZlt{\char`\<}
\def\PYZgt{\char`\>}
\def\PYZsh{\char`\#}
\def\PYZpc{\char`\%}
\def\PYZdl{\char`\$}
\def\PYZhy{\char`\-}
\def\PYZsq{\char`\'}
\def\PYZdq{\char`\"}
\def\PYZti{\char`\~}
% for compatibility with earlier versions
\def\PYZat{@}
\def\PYZlb{[}
\def\PYZrb{]}
\makeatother


    % For linebreaks inside Verbatim environment from package fancyvrb.
    \makeatletter
        \newbox\Wrappedcontinuationbox
        \newbox\Wrappedvisiblespacebox
        \newcommand*\Wrappedvisiblespace {\textcolor{red}{\textvisiblespace}}
        \newcommand*\Wrappedcontinuationsymbol {\textcolor{red}{\llap{\tiny$\m@th\hookrightarrow$}}}
        \newcommand*\Wrappedcontinuationindent {3ex }
        \newcommand*\Wrappedafterbreak {\kern\Wrappedcontinuationindent\copy\Wrappedcontinuationbox}
        % Take advantage of the already applied Pygments mark-up to insert
        % potential linebreaks for TeX processing.
        %        {, <, #, %, $, ' and ": go to next line.
        %        _, }, ^, &, >, - and ~: stay at end of broken line.
        % Use of \textquotesingle for straight quote.
        \newcommand*\Wrappedbreaksatspecials {%
            \def\PYGZus{\discretionary{\char`\_}{\Wrappedafterbreak}{\char`\_}}%
            \def\PYGZob{\discretionary{}{\Wrappedafterbreak\char`\{}{\char`\{}}%
            \def\PYGZcb{\discretionary{\char`\}}{\Wrappedafterbreak}{\char`\}}}%
            \def\PYGZca{\discretionary{\char`\^}{\Wrappedafterbreak}{\char`\^}}%
            \def\PYGZam{\discretionary{\char`\&}{\Wrappedafterbreak}{\char`\&}}%
            \def\PYGZlt{\discretionary{}{\Wrappedafterbreak\char`\<}{\char`\<}}%
            \def\PYGZgt{\discretionary{\char`\>}{\Wrappedafterbreak}{\char`\>}}%
            \def\PYGZsh{\discretionary{}{\Wrappedafterbreak\char`\#}{\char`\#}}%
            \def\PYGZpc{\discretionary{}{\Wrappedafterbreak\char`\%}{\char`\%}}%
            \def\PYGZdl{\discretionary{}{\Wrappedafterbreak\char`\$}{\char`\$}}%
            \def\PYGZhy{\discretionary{\char`\-}{\Wrappedafterbreak}{\char`\-}}%
            \def\PYGZsq{\discretionary{}{\Wrappedafterbreak\textquotesingle}{\textquotesingle}}%
            \def\PYGZdq{\discretionary{}{\Wrappedafterbreak\char`\"}{\char`\"}}%
            \def\PYGZti{\discretionary{\char`\~}{\Wrappedafterbreak}{\char`\~}}%
        }
        % Some characters . , ; ? ! / are not pygmentized.
        % This macro makes them "active" and they will insert potential linebreaks
        \newcommand*\Wrappedbreaksatpunct {%
            \lccode`\~`\.\lowercase{\def~}{\discretionary{\hbox{\char`\.}}{\Wrappedafterbreak}{\hbox{\char`\.}}}%
            \lccode`\~`\,\lowercase{\def~}{\discretionary{\hbox{\char`\,}}{\Wrappedafterbreak}{\hbox{\char`\,}}}%
            \lccode`\~`\;\lowercase{\def~}{\discretionary{\hbox{\char`\;}}{\Wrappedafterbreak}{\hbox{\char`\;}}}%
            \lccode`\~`\:\lowercase{\def~}{\discretionary{\hbox{\char`\:}}{\Wrappedafterbreak}{\hbox{\char`\:}}}%
            \lccode`\~`\?\lowercase{\def~}{\discretionary{\hbox{\char`\?}}{\Wrappedafterbreak}{\hbox{\char`\?}}}%
            \lccode`\~`\!\lowercase{\def~}{\discretionary{\hbox{\char`\!}}{\Wrappedafterbreak}{\hbox{\char`\!}}}%
            \lccode`\~`\/\lowercase{\def~}{\discretionary{\hbox{\char`\/}}{\Wrappedafterbreak}{\hbox{\char`\/}}}%
            \catcode`\.\active
            \catcode`\,\active
            \catcode`\;\active
            \catcode`\:\active
            \catcode`\?\active
            \catcode`\!\active
            \catcode`\/\active
            \lccode`\~`\~
        }
    \makeatother
    \setcounter{secnumdepth}{0}

    \let\OriginalVerbatim=\Verbatim
    \makeatletter
    \renewcommand{\Verbatim}[1][1]{%
        %\parskip\z@skip
        \sbox\Wrappedcontinuationbox {\Wrappedcontinuationsymbol}%
        \sbox\Wrappedvisiblespacebox {\FV@SetupFont\Wrappedvisiblespace}%
        \def\FancyVerbFormatLine ##1{\hsize\linewidth
            \vtop{\raggedright\hyphenpenalty\z@\exhyphenpenalty\z@
                \doublehyphendemerits\z@\finalhyphendemerits\z@
                \strut ##1\strut}%
        }%
        % If the linebreak is at a space, the latter will be displayed as visible
        % space at end of first line, and a continuation symbol starts next line.
        % Stretch/shrink are however usually zero for typewriter font.
        \def\FV@Space {%
            \nobreak\hskip\z@ plus\fontdimen3\font minus\fontdimen4\font
            \discretionary{\copy\Wrappedvisiblespacebox}{\Wrappedafterbreak}
            {\kern\fontdimen2\font}%
        }%

        % Allow breaks at special characters using \PYG... macros.
        \Wrappedbreaksatspecials
        % Breaks at punctuation characters . , ; ? ! and / need catcode=\active
        \OriginalVerbatim[#1,codes*=\Wrappedbreaksatpunct]%
    }
    \makeatother

    % Exact colors from NB
    \definecolor{incolor}{HTML}{303F9F}
    \definecolor{outcolor}{HTML}{D84315}
    \definecolor{cellborder}{HTML}{CFCFCF}
    \definecolor{cellbackground}{HTML}{F7F7F7}

    % prompt
    \makeatletter
    \newcommand{\boxspacing}{\kern\kvtcb@left@rule\kern\kvtcb@boxsep}
    \makeatother
    \newcommand{\prompt}[4]{
        {\ttfamily\llap{{\color{#2}[#3]:\hspace{3pt}#4}}\vspace{-\baselineskip}}
    }
    

    
    % Prevent overflowing lines due to hard-to-break entities
    \sloppy
    % Setup hyperref package
    \hypersetup{
      breaklinks=true,  % so long urls are correctly broken across lines
      colorlinks=true,
      urlcolor=urlcolor,
      linkcolor=linkcolor,
      citecolor=citecolor,
      }
    % Slightly bigger margins than the latex defaults
    
    \geometry{verbose,tmargin=1in,bmargin=1in,lmargin=1in,rmargin=1in}
    
    

\begin{document}
    
    \maketitle
    
    

    
    \section{Práctica 1 --- Problema de Clique
Máximo}\label{pruxe1ctica-1-problema-de-clique-muxe1ximo}

\textbf{Autor:} Jordi Blasco Lozano

    \paragraph{Índice}\label{uxedndice}

\begin{enumerate}
\def\labelenumi{\arabic{enumi}.}
\tightlist
\item
  Resumen
\item
  Introducción
\item
  Carga de imports y conjuntos de datos
\item
  Tareas de implementación

  \begin{itemize}
  \tightlist
  \item
    4.1 Utilidades generales
  \item
    4.2 Fuerza bruta (parte 0)

    \begin{itemize}
    \tightlist
    \item
      Tarea 0.1: \texttt{is\_clique()}
    \item
      Tarea 0.2: \texttt{brute\_force\_max\_clique()}
    \item
      Tarea 0.3: \texttt{test\_brute\_force()}
    \item
      Tarea 0.4: evaluación sobre dataset
    \end{itemize}
  \item
    4.3 Formulación de Motzkin-Straus (parte 1)

    \begin{itemize}
    \tightlist
    \item
      Tarea 1.1: funciones objetivo
    \item
      Tarea 1.2: verificación del teorema
    \end{itemize}
  \item
    4.4 Dinámicas replicadoras (parte 2)

    \begin{itemize}
    \tightlist
    \item
      Tarea 2.1: \texttt{replicator\_dynamics()}
    \item
      Tarea 2.2: \texttt{decode\_clique()}
    \item
      Tarea 2.3: visualización de convergencia
    \item
      Tarea 2.4: evaluación sobre dataset
    \end{itemize}
  \item
    4.5 GNN + Dinámicas replicadoras (parte 3)

    \begin{itemize}
    \tightlist
    \item
      Tarea 3.1: \texttt{compute\_node\_features()}
    \item
      Tarea 3.2: \texttt{MaxCliqueGNN}
    \item
      Tarea 3.3: RD diferenciable
    \item
      Tarea 3.4: \texttt{maxclique\_loss()}
    \item
      Tarea 3.5: entrenamiento
    \item
      Tarea 3.6: evaluación
    \end{itemize}
  \item
    4.6 Presentación de resultados
  \end{itemize}
\item
  Ejercicio 1

  \begin{itemize}
  \tightlist
  \item
    5.1 Implementar \texttt{is\_clique()} y
    \texttt{brute\_force\_max\_clique()}
  \item
    5.2 Prueba en un grafo sintético pequeño
  \item
    5.3 Aplicación a los primeros 20 grafos de IMDB y contraste
    adicional con COLLAB
  \item
    5.4 Informe del ejercicio
  \end{itemize}
\item
  Ejercicio 2

  \begin{itemize}
  \tightlist
  \item
    6.1 Implementación de funciones objetivo
  \item
    6.2 Crear 3 grafos de test con máximas cliques conocidas
  \item
    6.3 Verificar que el teorema se cumple en cada caso
  \item
    6.4 Demostrar el problema de la solución espuria en el grafo cherry
  \item
    6.5 Demostrar que la regularización lo corrige
  \item
    6.6 Informe del ejercicio
  \end{itemize}
\item
  Ejercicio 3

  \begin{itemize}
  \tightlist
  \item
    7.1 Implementar RD con el decodificador
  \item
    7.2 Visualización de la convergencia en un grafo pequeño
  \item
    7.3 Aplicación de RD en IMDB y COLLAB
  \item
    7.4 Comparación con fuerza bruta en los grafos pequeños de IMDB
  \item
    7.5 Informe del ejercicio
  \end{itemize}
\item
  Ejercicio 4

  \begin{itemize}
  \tightlist
  \item
    8.1 Implementar el pipeline completo
  \item
    8.2 Entrenar en IMDB y evaluar en grafos hold-out
  \item
    8.3 Comparación de los tres enfoques

    \begin{itemize}
    \tightlist
    \item
      Fuerza Bruta
    \item
      RD puro desde el baricentro
    \item
      GNN + RD
    \end{itemize}
  \item
    8.4 Ablación: variar el número de iteraciones de RD (train: 0,1,3,5;
    test: 10,50,100)
  \item
    8.5 Informe del ejercicio
  \end{itemize}
\item
  Ejercicio 5

  \begin{itemize}
  \tightlist
  \item
    9.1 Tabla resumen final
  \item
    9.2 Informe comparativo final
  \end{itemize}
\item
  Conclusiones
\end{enumerate}

    \subsection{1. Resumen}\label{resumen}

    En esta práctica se ha abordado el problema de clique máximo desde tres
niveles de complejidad. En primer lugar se ha implementado una
referencia exacta por fuerza bruta, útil para fijar un ground truth en
grafos pequeños. Después se ha trabajado con la relajación continua de
Motzkin-Straus y con Replicator Dynamics como procedimiento de
optimización sobre el simplex. Por último se ha añadido un modelo
híbrido GNN + RD para estudiar si una inicialización aprendida puede
mejorar el comportamiento de RD puro.

El notebook se ha reorganizado para que la sección de implementación
quede separada de los ejercicios, pero respetando los apartados del
enunciado.

    \subsection{2. Introducción}\label{introducciuxf3n}

    Una clique en un grafo no dirigido es un subconjunto de nodos en el que
todos los pares posibles están conectados. La clique máxima es la de
mayor tamaño, y encontrarla es un problema NP-hard. Esa dificultad
explica por qué una estrategia exacta solo es útil en instancias
pequeñas y por qué tiene sentido estudiar formulaciones relajadas y
métodos iterativos.

En esta práctica se ha seguido la secuencia natural del enunciado.
Primero se ha programado fuerza bruta como referencia exacta. Después se
ha comprobado numéricamente el teorema de Motzkin-Straus y el papel de
la regularización frente a soluciones espurias. A continuación se ha
implementado Replicator Dynamics y se ha evaluado sobre IMDB y COLLAB.
Finalmente se ha entrenado un modelo híbrido GNN + RD para comparar su
comportamiento frente a RD puro y analizar una pequeña ablación sobre el
número de iteraciones de refinamiento.

    \subsection{3. Carga de imports y conjuntos de
datos}\label{carga-de-imports-y-conjuntos-de-datos}

    \begin{tcolorbox}[breakable, size=fbox, boxrule=1pt, pad at break*=1mm,colback=cellbackground, colframe=cellborder]
\prompt{In}{incolor}{253}{\boxspacing}
\begin{Verbatim}[commandchars=\\\{\}]
\PY{k+kn}{import}\PY{+w}{ }\PY{n+nn}{random}
\PY{k+kn}{import}\PY{+w}{ }\PY{n+nn}{time}
\PY{k+kn}{from}\PY{+w}{ }\PY{n+nn}{itertools}\PY{+w}{ }\PY{k+kn}{import} \PY{n}{combinations}
\PY{k+kn}{from}\PY{+w}{ }\PY{n+nn}{pathlib}\PY{+w}{ }\PY{k+kn}{import} \PY{n}{Path}

\PY{k+kn}{import}\PY{+w}{ }\PY{n+nn}{matplotlib}\PY{n+nn}{.}\PY{n+nn}{pyplot}\PY{+w}{ }\PY{k}{as}\PY{+w}{ }\PY{n+nn}{plt}
\PY{k+kn}{import}\PY{+w}{ }\PY{n+nn}{networkx}\PY{+w}{ }\PY{k}{as}\PY{+w}{ }\PY{n+nn}{nx}
\PY{k+kn}{import}\PY{+w}{ }\PY{n+nn}{numpy}\PY{+w}{ }\PY{k}{as}\PY{+w}{ }\PY{n+nn}{np}
\PY{k+kn}{import}\PY{+w}{ }\PY{n+nn}{pandas}\PY{+w}{ }\PY{k}{as}\PY{+w}{ }\PY{n+nn}{pd}
\PY{k+kn}{import}\PY{+w}{ }\PY{n+nn}{torch}
\PY{k+kn}{import}\PY{+w}{ }\PY{n+nn}{torch}\PY{n+nn}{.}\PY{n+nn}{nn}\PY{+w}{ }\PY{k}{as}\PY{+w}{ }\PY{n+nn}{nn}
\PY{k+kn}{import}\PY{+w}{ }\PY{n+nn}{torch}\PY{n+nn}{.}\PY{n+nn}{nn}\PY{n+nn}{.}\PY{n+nn}{functional}\PY{+w}{ }\PY{k}{as}\PY{+w}{ }\PY{n+nn}{F}
\PY{k+kn}{from}\PY{+w}{ }\PY{n+nn}{IPython}\PY{n+nn}{.}\PY{n+nn}{display}\PY{+w}{ }\PY{k+kn}{import} \PY{n}{Markdown}\PY{p}{,} \PY{n}{display}
\PY{k+kn}{from}\PY{+w}{ }\PY{n+nn}{torch\PYZus{}geometric}\PY{n+nn}{.}\PY{n+nn}{datasets}\PY{+w}{ }\PY{k+kn}{import} \PY{n}{TUDataset}
\PY{k+kn}{from}\PY{+w}{ }\PY{n+nn}{torch\PYZus{}geometric}\PY{n+nn}{.}\PY{n+nn}{nn}\PY{+w}{ }\PY{k+kn}{import} \PY{n}{GCNConv}
\PY{k+kn}{from}\PY{+w}{ }\PY{n+nn}{torch\PYZus{}geometric}\PY{n+nn}{.}\PY{n+nn}{utils}\PY{+w}{ }\PY{k+kn}{import} \PY{n}{to\PYZus{}networkx}

\PY{n}{plt}\PY{o}{.}\PY{n}{style}\PY{o}{.}\PY{n}{use}\PY{p}{(}\PY{l+s+s2}{\PYZdq{}}\PY{l+s+s2}{seaborn\PYZhy{}v0\PYZus{}8\PYZhy{}whitegrid}\PY{l+s+s2}{\PYZdq{}}\PY{p}{)}
\PY{n}{np}\PY{o}{.}\PY{n}{set\PYZus{}printoptions}\PY{p}{(}\PY{n}{precision}\PY{o}{=}\PY{l+m+mi}{4}\PY{p}{,} \PY{n}{suppress}\PY{o}{=}\PY{k+kc}{True}\PY{p}{)}
\PY{n}{pd}\PY{o}{.}\PY{n}{set\PYZus{}option}\PY{p}{(}\PY{l+s+s2}{\PYZdq{}}\PY{l+s+s2}{display.max\PYZus{}columns}\PY{l+s+s2}{\PYZdq{}}\PY{p}{,} \PY{k+kc}{None}\PY{p}{)}
\PY{n}{pd}\PY{o}{.}\PY{n}{set\PYZus{}option}\PY{p}{(}\PY{l+s+s2}{\PYZdq{}}\PY{l+s+s2}{display.precision}\PY{l+s+s2}{\PYZdq{}}\PY{p}{,} \PY{l+m+mi}{4}\PY{p}{)}

\PY{n}{SEED} \PY{o}{=} \PY{l+m+mi}{42}
\PY{n}{random}\PY{o}{.}\PY{n}{seed}\PY{p}{(}\PY{n}{SEED}\PY{p}{)}
\PY{n}{np}\PY{o}{.}\PY{n}{random}\PY{o}{.}\PY{n}{seed}\PY{p}{(}\PY{n}{SEED}\PY{p}{)}
\PY{n}{torch}\PY{o}{.}\PY{n}{manual\PYZus{}seed}\PY{p}{(}\PY{n}{SEED}\PY{p}{)}

\PY{n}{DEVICE} \PY{o}{=} \PY{n}{torch}\PY{o}{.}\PY{n}{device}\PY{p}{(}\PY{l+s+s2}{\PYZdq{}}\PY{l+s+s2}{cuda}\PY{l+s+s2}{\PYZdq{}} \PY{k}{if} \PY{n}{torch}\PY{o}{.}\PY{n}{cuda}\PY{o}{.}\PY{n}{is\PYZus{}available}\PY{p}{(}\PY{p}{)} \PY{k}{else} \PY{l+s+s2}{\PYZdq{}}\PY{l+s+s2}{cpu}\PY{l+s+s2}{\PYZdq{}}\PY{p}{)}
\PY{n}{BASE\PYZus{}DIR} \PY{o}{=} \PY{n}{Path}\PY{o}{.}\PY{n}{cwd}\PY{p}{(}\PY{p}{)}
\PY{n}{DATA\PYZus{}DIR} \PY{o}{=} \PY{n}{BASE\PYZus{}DIR} \PY{o}{/} \PY{l+s+s2}{\PYZdq{}}\PY{l+s+s2}{data}\PY{l+s+s2}{\PYZdq{}}
\PY{n}{IMAGES\PYZus{}DIR} \PY{o}{=} \PY{n}{BASE\PYZus{}DIR} \PY{o}{/} \PY{l+s+s2}{\PYZdq{}}\PY{l+s+s2}{images}\PY{l+s+s2}{\PYZdq{}}
\PY{n}{IMAGES\PYZus{}DIR}\PY{o}{.}\PY{n}{mkdir}\PY{p}{(}\PY{n}{parents}\PY{o}{=}\PY{k+kc}{True}\PY{p}{,} \PY{n}{exist\PYZus{}ok}\PY{o}{=}\PY{k+kc}{True}\PY{p}{)}

\PY{n}{imdb\PYZus{}dataset} \PY{o}{=} \PY{n}{TUDataset}\PY{p}{(}\PY{n}{root}\PY{o}{=}\PY{n+nb}{str}\PY{p}{(}\PY{n}{DATA\PYZus{}DIR}\PY{p}{)}\PY{p}{,} \PY{n}{name}\PY{o}{=}\PY{l+s+s2}{\PYZdq{}}\PY{l+s+s2}{IMDB\PYZhy{}BINARY}\PY{l+s+s2}{\PYZdq{}}\PY{p}{)}
\PY{n}{collab\PYZus{}dataset} \PY{o}{=} \PY{n}{TUDataset}\PY{p}{(}\PY{n}{root}\PY{o}{=}\PY{n+nb}{str}\PY{p}{(}\PY{n}{DATA\PYZus{}DIR}\PY{p}{)}\PY{p}{,} \PY{n}{name}\PY{o}{=}\PY{l+s+s2}{\PYZdq{}}\PY{l+s+s2}{COLLAB}\PY{l+s+s2}{\PYZdq{}}\PY{p}{)}

\PY{n+nb}{print}\PY{p}{(}\PY{l+s+sa}{f}\PY{l+s+s2}{\PYZdq{}}\PY{l+s+s2}{Dispositivo de trabajo: }\PY{l+s+si}{\PYZob{}}\PY{n}{DEVICE}\PY{l+s+si}{\PYZcb{}}\PY{l+s+s2}{\PYZdq{}}\PY{p}{)}
\PY{n+nb}{print}\PY{p}{(}\PY{l+s+sa}{f}\PY{l+s+s2}{\PYZdq{}}\PY{l+s+s2}{IMDB\PYZhy{}BINARY: }\PY{l+s+si}{\PYZob{}}\PY{n+nb}{len}\PY{p}{(}\PY{n}{imdb\PYZus{}dataset}\PY{p}{)}\PY{l+s+si}{\PYZcb{}}\PY{l+s+s2}{ grafos}\PY{l+s+s2}{\PYZdq{}}\PY{p}{)}
\PY{n+nb}{print}\PY{p}{(}\PY{l+s+sa}{f}\PY{l+s+s2}{\PYZdq{}}\PY{l+s+s2}{COLLAB: }\PY{l+s+si}{\PYZob{}}\PY{n+nb}{len}\PY{p}{(}\PY{n}{collab\PYZus{}dataset}\PY{p}{)}\PY{l+s+si}{\PYZcb{}}\PY{l+s+s2}{ grafos}\PY{l+s+s2}{\PYZdq{}}\PY{p}{)}

\PY{n}{sample\PYZus{}graph} \PY{o}{=} \PY{n}{imdb\PYZus{}dataset}\PY{p}{[}\PY{l+m+mi}{0}\PY{p}{]}
\PY{n+nb}{print}\PY{p}{(}\PY{l+s+s2}{\PYZdq{}}\PY{l+s+se}{\PYZbs{}n}\PY{l+s+s2}{Ejemplo de IMDB\PYZhy{}BINARY:}\PY{l+s+s2}{\PYZdq{}}\PY{p}{)}
\PY{n+nb}{print}\PY{p}{(}\PY{l+s+sa}{f}\PY{l+s+s2}{\PYZdq{}}\PY{l+s+s2}{  Nodos: }\PY{l+s+si}{\PYZob{}}\PY{n}{sample\PYZus{}graph}\PY{o}{.}\PY{n}{num\PYZus{}nodes}\PY{l+s+si}{\PYZcb{}}\PY{l+s+s2}{\PYZdq{}}\PY{p}{)}
\PY{n+nb}{print}\PY{p}{(}\PY{l+s+sa}{f}\PY{l+s+s2}{\PYZdq{}}\PY{l+s+s2}{  Aristas: }\PY{l+s+si}{\PYZob{}}\PY{n}{sample\PYZus{}graph}\PY{o}{.}\PY{n}{num\PYZus{}edges}\PY{l+s+si}{\PYZcb{}}\PY{l+s+s2}{\PYZdq{}}\PY{p}{)}
\PY{n+nb}{print}\PY{p}{(}\PY{l+s+sa}{f}\PY{l+s+s2}{\PYZdq{}}\PY{l+s+s2}{  Etiqueta del grafo: }\PY{l+s+si}{\PYZob{}}\PY{n+nb}{int}\PY{p}{(}\PY{n}{sample\PYZus{}graph}\PY{o}{.}\PY{n}{y}\PY{o}{.}\PY{n}{item}\PY{p}{(}\PY{p}{)}\PY{p}{)}\PY{+w}{ }\PY{k}{if}\PY{+w}{ }\PY{n+nb}{hasattr}\PY{p}{(}\PY{n}{sample\PYZus{}graph}\PY{p}{,}\PY{+w}{ }\PY{l+s+s1}{\PYZsq{}}\PY{l+s+s1}{y}\PY{l+s+s1}{\PYZsq{}}\PY{p}{)}\PY{+w}{ }\PY{k}{else}\PY{+w}{ }\PY{l+s+s1}{\PYZsq{}}\PY{l+s+s1}{N/D}\PY{l+s+s1}{\PYZsq{}}\PY{l+s+si}{\PYZcb{}}\PY{l+s+s2}{\PYZdq{}}\PY{p}{)}
\end{Verbatim}
\end{tcolorbox}

    \begin{Verbatim}[commandchars=\\\{\}]
Dispositivo de trabajo: cuda
IMDB-BINARY: 1000 grafos
COLLAB: 5000 grafos

Ejemplo de IMDB-BINARY:
  Nodos: 20
  Aristas: 146
  Etiqueta del grafo: 0
    \end{Verbatim}

    \subsection{4. Tareas de
implementación}\label{tareas-de-implementaciuxf3n}

    En esta sección se han agrupado las funciones y clases reutilizables.

    \paragraph{4.1 Utilidades generales}\label{utilidades-generales}

    Aquí se han reunido funciones auxiliares de conversión, medida de
tiempos, particionado y resumen de resultados. No pertenecen al
enunciado como tareas específicas, pero permiten que los ejercicios
posteriores queden más limpios.

    \begin{tcolorbox}[breakable, size=fbox, boxrule=1pt, pad at break*=1mm,colback=cellbackground, colframe=cellborder]
\prompt{In}{incolor}{254}{\boxspacing}
\begin{Verbatim}[commandchars=\\\{\}]
\PY{k}{def}\PY{+w}{ }\PY{n+nf}{pyg\PYZus{}to\PYZus{}adjacency}\PY{p}{(}\PY{n}{data}\PY{p}{)}\PY{p}{:}
\PY{+w}{    }\PY{l+s+sd}{\PYZdq{}\PYZdq{}\PYZdq{}Convierte un grafo de PyG a matriz de adyacencia densa de tipo float.\PYZdq{}\PYZdq{}\PYZdq{}}
    \PY{n}{G} \PY{o}{=} \PY{n}{to\PYZus{}networkx}\PY{p}{(}\PY{n}{data}\PY{p}{,} \PY{n}{to\PYZus{}undirected}\PY{o}{=}\PY{k+kc}{True}\PY{p}{)}
    \PY{n}{A} \PY{o}{=} \PY{n}{nx}\PY{o}{.}\PY{n}{to\PYZus{}numpy\PYZus{}array}\PY{p}{(}\PY{n}{G}\PY{p}{,} \PY{n}{dtype}\PY{o}{=}\PY{n+nb}{float}\PY{p}{)}
    \PY{n}{np}\PY{o}{.}\PY{n}{fill\PYZus{}diagonal}\PY{p}{(}\PY{n}{A}\PY{p}{,} \PY{l+m+mf}{0.0}\PY{p}{)}
    \PY{k}{return} \PY{n}{A}


\PY{k}{def}\PY{+w}{ }\PY{n+nf}{clique\PYZus{}to\PYZus{}simplex}\PY{p}{(}\PY{n}{clique}\PY{p}{,} \PY{n}{n}\PY{p}{)}\PY{p}{:}
\PY{+w}{    }\PY{l+s+sd}{\PYZdq{}\PYZdq{}\PYZdq{}Construye el punto del simplex asociado uniformemente a una clique.\PYZdq{}\PYZdq{}\PYZdq{}}
    \PY{n}{x} \PY{o}{=} \PY{n}{np}\PY{o}{.}\PY{n}{zeros}\PY{p}{(}\PY{n}{n}\PY{p}{,} \PY{n}{dtype}\PY{o}{=}\PY{n+nb}{float}\PY{p}{)}
    \PY{k}{if} \PY{n+nb}{len}\PY{p}{(}\PY{n}{clique}\PY{p}{)} \PY{o}{==} \PY{l+m+mi}{0}\PY{p}{:}
        \PY{k}{return} \PY{n}{x}
    \PY{n}{x}\PY{p}{[}\PY{n}{clique}\PY{p}{]} \PY{o}{=} \PY{l+m+mf}{1.0} \PY{o}{/} \PY{n+nb}{len}\PY{p}{(}\PY{n}{clique}\PY{p}{)}
    \PY{k}{return} \PY{n}{x}


\PY{k}{def}\PY{+w}{ }\PY{n+nf}{timed\PYZus{}call}\PY{p}{(}\PY{n}{fn}\PY{p}{,} \PY{o}{*}\PY{n}{args}\PY{p}{,} \PY{o}{*}\PY{o}{*}\PY{n}{kwargs}\PY{p}{)}\PY{p}{:}
\PY{+w}{    }\PY{l+s+sd}{\PYZdq{}\PYZdq{}\PYZdq{}Ejecuta una función y devuelve resultado y tiempo transcurrido.\PYZdq{}\PYZdq{}\PYZdq{}}
    \PY{n}{start} \PY{o}{=} \PY{n}{time}\PY{o}{.}\PY{n}{time}\PY{p}{(}\PY{p}{)}
    \PY{n}{out} \PY{o}{=} \PY{n}{fn}\PY{p}{(}\PY{o}{*}\PY{n}{args}\PY{p}{,} \PY{o}{*}\PY{o}{*}\PY{n}{kwargs}\PY{p}{)}
    \PY{n}{elapsed} \PY{o}{=} \PY{n}{time}\PY{o}{.}\PY{n}{time}\PY{p}{(}\PY{p}{)} \PY{o}{\PYZhy{}} \PY{n}{start}
    \PY{k}{return} \PY{n}{out}\PY{p}{,} \PY{n}{elapsed}


\PY{k}{def}\PY{+w}{ }\PY{n+nf}{summarize\PYZus{}results}\PY{p}{(}\PY{n}{results}\PY{p}{)}\PY{p}{:}
\PY{+w}{    }\PY{l+s+sd}{\PYZdq{}\PYZdq{}\PYZdq{}Resume una lista de diccionarios homogéneos en un DataFrame y estadísticas básicas.\PYZdq{}\PYZdq{}\PYZdq{}}
    \PY{k}{if} \PY{n+nb}{len}\PY{p}{(}\PY{n}{results}\PY{p}{)} \PY{o}{==} \PY{l+m+mi}{0}\PY{p}{:}
        \PY{k}{return} \PY{n}{pd}\PY{o}{.}\PY{n}{DataFrame}\PY{p}{(}\PY{p}{)}\PY{p}{,} \PY{p}{\PYZob{}}\PY{p}{\PYZcb{}}

    \PY{n}{df} \PY{o}{=} \PY{n}{pd}\PY{o}{.}\PY{n}{DataFrame}\PY{p}{(}\PY{n}{results}\PY{p}{)}
    \PY{n}{summary} \PY{o}{=} \PY{p}{\PYZob{}}\PY{p}{\PYZcb{}}
    \PY{k}{for} \PY{n}{column} \PY{o+ow}{in} \PY{p}{[}\PY{l+s+s2}{\PYZdq{}}\PY{l+s+s2}{clique\PYZus{}size}\PY{l+s+s2}{\PYZdq{}}\PY{p}{,} \PY{l+s+s2}{\PYZdq{}}\PY{l+s+s2}{time}\PY{l+s+s2}{\PYZdq{}}\PY{p}{,} \PY{l+s+s2}{\PYZdq{}}\PY{l+s+s2}{iterations}\PY{l+s+s2}{\PYZdq{}}\PY{p}{]}\PY{p}{:}
        \PY{k}{if} \PY{n}{column} \PY{o+ow}{in} \PY{n}{df}\PY{o}{.}\PY{n}{columns}\PY{p}{:}
            \PY{n}{summary}\PY{p}{[}\PY{n}{column}\PY{p}{]} \PY{o}{=} \PY{p}{\PYZob{}}
                \PY{l+s+s2}{\PYZdq{}}\PY{l+s+s2}{mean}\PY{l+s+s2}{\PYZdq{}}\PY{p}{:} \PY{n+nb}{float}\PY{p}{(}\PY{n}{df}\PY{p}{[}\PY{n}{column}\PY{p}{]}\PY{o}{.}\PY{n}{mean}\PY{p}{(}\PY{p}{)}\PY{p}{)}\PY{p}{,}
                \PY{l+s+s2}{\PYZdq{}}\PY{l+s+s2}{std}\PY{l+s+s2}{\PYZdq{}}\PY{p}{:} \PY{n+nb}{float}\PY{p}{(}\PY{n}{df}\PY{p}{[}\PY{n}{column}\PY{p}{]}\PY{o}{.}\PY{n}{std}\PY{p}{(}\PY{n}{ddof}\PY{o}{=}\PY{l+m+mi}{0}\PY{p}{)}\PY{p}{)}\PY{p}{,}
            \PY{p}{\PYZcb{}}
    \PY{k}{if} \PY{l+s+s2}{\PYZdq{}}\PY{l+s+s2}{valid}\PY{l+s+s2}{\PYZdq{}} \PY{o+ow}{in} \PY{n}{df}\PY{o}{.}\PY{n}{columns}\PY{p}{:}
        \PY{n}{summary}\PY{p}{[}\PY{l+s+s2}{\PYZdq{}}\PY{l+s+s2}{valid\PYZus{}rate}\PY{l+s+s2}{\PYZdq{}}\PY{p}{]} \PY{o}{=} \PY{n+nb}{float}\PY{p}{(}\PY{n}{df}\PY{p}{[}\PY{l+s+s2}{\PYZdq{}}\PY{l+s+s2}{valid}\PY{l+s+s2}{\PYZdq{}}\PY{p}{]}\PY{o}{.}\PY{n}{mean}\PY{p}{(}\PY{p}{)}\PY{p}{)}
    \PY{n}{summary}\PY{p}{[}\PY{l+s+s2}{\PYZdq{}}\PY{l+s+s2}{num\PYZus{}graphs}\PY{l+s+s2}{\PYZdq{}}\PY{p}{]} \PY{o}{=} \PY{n+nb}{int}\PY{p}{(}\PY{n+nb}{len}\PY{p}{(}\PY{n}{df}\PY{p}{)}\PY{p}{)}
    \PY{k}{return} \PY{n}{df}\PY{p}{,} \PY{n}{summary}


\PY{k}{def}\PY{+w}{ }\PY{n+nf}{format\PYZus{}metric}\PY{p}{(}\PY{n}{value}\PY{p}{,} \PY{n}{digits}\PY{o}{=}\PY{l+m+mi}{3}\PY{p}{)}\PY{p}{:}
    \PY{k}{if} \PY{n}{value} \PY{o+ow}{is} \PY{k+kc}{None} \PY{o+ow}{or} \PY{p}{(}\PY{n+nb}{isinstance}\PY{p}{(}\PY{n}{value}\PY{p}{,} \PY{n+nb}{float}\PY{p}{)} \PY{o+ow}{and} \PY{n}{np}\PY{o}{.}\PY{n}{isnan}\PY{p}{(}\PY{n}{value}\PY{p}{)}\PY{p}{)}\PY{p}{:}
        \PY{k}{return} \PY{l+s+s2}{\PYZdq{}}\PY{l+s+s2}{N/D}\PY{l+s+s2}{\PYZdq{}}
    \PY{k}{return} \PY{l+s+sa}{f}\PY{l+s+s2}{\PYZdq{}}\PY{l+s+si}{\PYZob{}}\PY{n}{value}\PY{l+s+si}{:}\PY{l+s+s2}{.}\PY{l+s+si}{\PYZob{}}\PY{n}{digits}\PY{l+s+si}{\PYZcb{}}\PY{l+s+s2}{f}\PY{l+s+si}{\PYZcb{}}\PY{l+s+s2}{\PYZdq{}}


\PY{k}{def}\PY{+w}{ }\PY{n+nf}{make\PYZus{}summary\PYZus{}dict}\PY{p}{(}\PY{n}{results}\PY{p}{)}\PY{p}{:}
    \PY{k}{if} \PY{n+nb}{len}\PY{p}{(}\PY{n}{results}\PY{p}{)} \PY{o}{==} \PY{l+m+mi}{0}\PY{p}{:}
        \PY{k}{return} \PY{p}{\PYZob{}}
            \PY{l+s+s2}{\PYZdq{}}\PY{l+s+s2}{num\PYZus{}graphs}\PY{l+s+s2}{\PYZdq{}}\PY{p}{:} \PY{l+m+mi}{0}\PY{p}{,}
            \PY{l+s+s2}{\PYZdq{}}\PY{l+s+s2}{mean\PYZus{}clique\PYZus{}size}\PY{l+s+s2}{\PYZdq{}}\PY{p}{:} \PY{n}{np}\PY{o}{.}\PY{n}{nan}\PY{p}{,}
            \PY{l+s+s2}{\PYZdq{}}\PY{l+s+s2}{mean\PYZus{}time}\PY{l+s+s2}{\PYZdq{}}\PY{p}{:} \PY{n}{np}\PY{o}{.}\PY{n}{nan}\PY{p}{,}
            \PY{l+s+s2}{\PYZdq{}}\PY{l+s+s2}{valid\PYZus{}rate}\PY{l+s+s2}{\PYZdq{}}\PY{p}{:} \PY{n}{np}\PY{o}{.}\PY{n}{nan}\PY{p}{,}
        \PY{p}{\PYZcb{}}
    \PY{n}{df} \PY{o}{=} \PY{n}{pd}\PY{o}{.}\PY{n}{DataFrame}\PY{p}{(}\PY{n}{results}\PY{p}{)}
    \PY{n}{clique\PYZus{}series} \PY{o}{=} \PY{n}{pd}\PY{o}{.}\PY{n}{to\PYZus{}numeric}\PY{p}{(}\PY{n}{df}\PY{p}{[}\PY{l+s+s2}{\PYZdq{}}\PY{l+s+s2}{clique\PYZus{}size}\PY{l+s+s2}{\PYZdq{}}\PY{p}{]}\PY{p}{,} \PY{n}{errors}\PY{o}{=}\PY{l+s+s2}{\PYZdq{}}\PY{l+s+s2}{coerce}\PY{l+s+s2}{\PYZdq{}}\PY{p}{)} \PY{k}{if} \PY{l+s+s2}{\PYZdq{}}\PY{l+s+s2}{clique\PYZus{}size}\PY{l+s+s2}{\PYZdq{}} \PY{o+ow}{in} \PY{n}{df} \PY{k}{else} \PY{n}{pd}\PY{o}{.}\PY{n}{Series}\PY{p}{(}\PY{n}{dtype}\PY{o}{=}\PY{n+nb}{float}\PY{p}{)}
    \PY{n}{time\PYZus{}series} \PY{o}{=} \PY{n}{pd}\PY{o}{.}\PY{n}{to\PYZus{}numeric}\PY{p}{(}\PY{n}{df}\PY{p}{[}\PY{l+s+s2}{\PYZdq{}}\PY{l+s+s2}{time}\PY{l+s+s2}{\PYZdq{}}\PY{p}{]}\PY{p}{,} \PY{n}{errors}\PY{o}{=}\PY{l+s+s2}{\PYZdq{}}\PY{l+s+s2}{coerce}\PY{l+s+s2}{\PYZdq{}}\PY{p}{)} \PY{k}{if} \PY{l+s+s2}{\PYZdq{}}\PY{l+s+s2}{time}\PY{l+s+s2}{\PYZdq{}} \PY{o+ow}{in} \PY{n}{df} \PY{k}{else} \PY{n}{pd}\PY{o}{.}\PY{n}{Series}\PY{p}{(}\PY{n}{dtype}\PY{o}{=}\PY{n+nb}{float}\PY{p}{)}
    \PY{n}{valid\PYZus{}series} \PY{o}{=} \PY{n}{pd}\PY{o}{.}\PY{n}{to\PYZus{}numeric}\PY{p}{(}\PY{n}{df}\PY{p}{[}\PY{l+s+s2}{\PYZdq{}}\PY{l+s+s2}{valid}\PY{l+s+s2}{\PYZdq{}}\PY{p}{]}\PY{p}{,} \PY{n}{errors}\PY{o}{=}\PY{l+s+s2}{\PYZdq{}}\PY{l+s+s2}{coerce}\PY{l+s+s2}{\PYZdq{}}\PY{p}{)} \PY{k}{if} \PY{l+s+s2}{\PYZdq{}}\PY{l+s+s2}{valid}\PY{l+s+s2}{\PYZdq{}} \PY{o+ow}{in} \PY{n}{df} \PY{k}{else} \PY{n}{pd}\PY{o}{.}\PY{n}{Series}\PY{p}{(}\PY{n}{dtype}\PY{o}{=}\PY{n+nb}{float}\PY{p}{)}
    \PY{k}{return} \PY{p}{\PYZob{}}
        \PY{l+s+s2}{\PYZdq{}}\PY{l+s+s2}{num\PYZus{}graphs}\PY{l+s+s2}{\PYZdq{}}\PY{p}{:} \PY{n+nb}{int}\PY{p}{(}\PY{n+nb}{len}\PY{p}{(}\PY{n}{df}\PY{p}{)}\PY{p}{)}\PY{p}{,}
        \PY{l+s+s2}{\PYZdq{}}\PY{l+s+s2}{mean\PYZus{}clique\PYZus{}size}\PY{l+s+s2}{\PYZdq{}}\PY{p}{:} \PY{n+nb}{float}\PY{p}{(}\PY{n}{clique\PYZus{}series}\PY{o}{.}\PY{n}{mean}\PY{p}{(}\PY{p}{)}\PY{p}{)} \PY{k}{if} \PY{o+ow}{not} \PY{n}{clique\PYZus{}series}\PY{o}{.}\PY{n}{empty} \PY{k}{else} \PY{n}{np}\PY{o}{.}\PY{n}{nan}\PY{p}{,}
        \PY{l+s+s2}{\PYZdq{}}\PY{l+s+s2}{mean\PYZus{}time}\PY{l+s+s2}{\PYZdq{}}\PY{p}{:} \PY{n+nb}{float}\PY{p}{(}\PY{n}{time\PYZus{}series}\PY{o}{.}\PY{n}{mean}\PY{p}{(}\PY{p}{)}\PY{p}{)} \PY{k}{if} \PY{o+ow}{not} \PY{n}{time\PYZus{}series}\PY{o}{.}\PY{n}{empty} \PY{k}{else} \PY{n}{np}\PY{o}{.}\PY{n}{nan}\PY{p}{,}
        \PY{l+s+s2}{\PYZdq{}}\PY{l+s+s2}{valid\PYZus{}rate}\PY{l+s+s2}{\PYZdq{}}\PY{p}{:} \PY{n+nb}{float}\PY{p}{(}\PY{n}{valid\PYZus{}series}\PY{o}{.}\PY{n}{mean}\PY{p}{(}\PY{p}{)}\PY{p}{)} \PY{k}{if} \PY{o+ow}{not} \PY{n}{valid\PYZus{}series}\PY{o}{.}\PY{n}{empty} \PY{k}{else} \PY{n}{np}\PY{o}{.}\PY{n}{nan}\PY{p}{,}
    \PY{p}{\PYZcb{}}


\PY{k}{def}\PY{+w}{ }\PY{n+nf}{plot\PYZus{}weighted\PYZus{}graph}\PY{p}{(}\PY{n}{A}\PY{p}{,} \PY{n}{weights}\PY{p}{,} \PY{n}{title}\PY{p}{)}\PY{p}{:}
\PY{+w}{    }\PY{l+s+sd}{\PYZdq{}\PYZdq{}\PYZdq{}Dibuja el grafo coloreando cada nodo según un peso.\PYZdq{}\PYZdq{}\PYZdq{}}
    \PY{n}{G} \PY{o}{=} \PY{n}{nx}\PY{o}{.}\PY{n}{from\PYZus{}numpy\PYZus{}array}\PY{p}{(}\PY{n}{A}\PY{p}{)}
    \PY{n}{pos} \PY{o}{=} \PY{n}{nx}\PY{o}{.}\PY{n}{spring\PYZus{}layout}\PY{p}{(}\PY{n}{G}\PY{p}{,} \PY{n}{seed}\PY{o}{=}\PY{n}{SEED}\PY{p}{)}
    \PY{n}{fig}\PY{p}{,} \PY{n}{ax} \PY{o}{=} \PY{n}{plt}\PY{o}{.}\PY{n}{subplots}\PY{p}{(}\PY{n}{figsize}\PY{o}{=}\PY{p}{(}\PY{l+m+mi}{5}\PY{p}{,} \PY{l+m+mi}{4}\PY{p}{)}\PY{p}{)}
    \PY{n}{nodes} \PY{o}{=} \PY{n}{nx}\PY{o}{.}\PY{n}{draw\PYZus{}networkx\PYZus{}nodes}\PY{p}{(}
        \PY{n}{G}\PY{p}{,}
        \PY{n}{pos}\PY{p}{,}
        \PY{n}{node\PYZus{}color}\PY{o}{=}\PY{n}{weights}\PY{p}{,}
        \PY{n}{cmap}\PY{o}{=}\PY{l+s+s2}{\PYZdq{}}\PY{l+s+s2}{viridis}\PY{l+s+s2}{\PYZdq{}}\PY{p}{,}
        \PY{n}{node\PYZus{}size}\PY{o}{=}\PY{l+m+mi}{350}\PY{p}{,}
        \PY{n}{ax}\PY{o}{=}\PY{n}{ax}\PY{p}{,}
    \PY{p}{)}
    \PY{n}{nx}\PY{o}{.}\PY{n}{draw\PYZus{}networkx\PYZus{}edges}\PY{p}{(}\PY{n}{G}\PY{p}{,} \PY{n}{pos}\PY{p}{,} \PY{n}{alpha}\PY{o}{=}\PY{l+m+mf}{0.5}\PY{p}{,} \PY{n}{width}\PY{o}{=}\PY{l+m+mf}{1.2}\PY{p}{,} \PY{n}{ax}\PY{o}{=}\PY{n}{ax}\PY{p}{)}
    \PY{n}{nx}\PY{o}{.}\PY{n}{draw\PYZus{}networkx\PYZus{}labels}\PY{p}{(}\PY{n}{G}\PY{p}{,} \PY{n}{pos}\PY{p}{,} \PY{n}{font\PYZus{}size}\PY{o}{=}\PY{l+m+mi}{9}\PY{p}{,} \PY{n}{ax}\PY{o}{=}\PY{n}{ax}\PY{p}{)}
    \PY{n}{plt}\PY{o}{.}\PY{n}{colorbar}\PY{p}{(}\PY{n}{nodes}\PY{p}{,} \PY{n}{ax}\PY{o}{=}\PY{n}{ax}\PY{p}{,} \PY{n}{label}\PY{o}{=}\PY{l+s+s2}{\PYZdq{}}\PY{l+s+s2}{Peso final}\PY{l+s+s2}{\PYZdq{}}\PY{p}{)}
    \PY{n}{ax}\PY{o}{.}\PY{n}{set\PYZus{}title}\PY{p}{(}\PY{n}{title}\PY{p}{)}
    \PY{n}{ax}\PY{o}{.}\PY{n}{axis}\PY{p}{(}\PY{l+s+s2}{\PYZdq{}}\PY{l+s+s2}{off}\PY{l+s+s2}{\PYZdq{}}\PY{p}{)}
    \PY{n}{plt}\PY{o}{.}\PY{n}{tight\PYZus{}layout}\PY{p}{(}\PY{p}{)}
    \PY{n}{plt}\PY{o}{.}\PY{n}{show}\PY{p}{(}\PY{p}{)}


\PY{k}{def}\PY{+w}{ }\PY{n+nf}{split\PYZus{}dataset}\PY{p}{(}\PY{n}{dataset}\PY{p}{,} \PY{n}{train\PYZus{}ratio}\PY{o}{=}\PY{l+m+mf}{0.8}\PY{p}{,} \PY{n}{max\PYZus{}graphs}\PY{o}{=}\PY{k+kc}{None}\PY{p}{)}\PY{p}{:}
\PY{+w}{    }\PY{l+s+sd}{\PYZdq{}\PYZdq{}\PYZdq{}Crea una partición secuencial sencilla train/test sobre un subconjunto del dataset.\PYZdq{}\PYZdq{}\PYZdq{}}
    \PY{n}{total} \PY{o}{=} \PY{n+nb}{len}\PY{p}{(}\PY{n}{dataset}\PY{p}{)} \PY{k}{if} \PY{n}{max\PYZus{}graphs} \PY{o+ow}{is} \PY{k+kc}{None} \PY{k}{else} \PY{n+nb}{min}\PY{p}{(}\PY{n+nb}{len}\PY{p}{(}\PY{n}{dataset}\PY{p}{)}\PY{p}{,} \PY{n}{max\PYZus{}graphs}\PY{p}{)}
    \PY{n}{subset} \PY{o}{=} \PY{p}{[}\PY{n}{dataset}\PY{p}{[}\PY{n}{i}\PY{p}{]} \PY{k}{for} \PY{n}{i} \PY{o+ow}{in} \PY{n+nb}{range}\PY{p}{(}\PY{n}{total}\PY{p}{)}\PY{p}{]}
    \PY{n}{split\PYZus{}idx} \PY{o}{=} \PY{n+nb}{max}\PY{p}{(}\PY{l+m+mi}{1}\PY{p}{,} \PY{n+nb}{int}\PY{p}{(}\PY{n}{total} \PY{o}{*} \PY{n}{train\PYZus{}ratio}\PY{p}{)}\PY{p}{)}
    \PY{n}{split\PYZus{}idx} \PY{o}{=} \PY{n+nb}{min}\PY{p}{(}\PY{n}{split\PYZus{}idx}\PY{p}{,} \PY{n}{total} \PY{o}{\PYZhy{}} \PY{l+m+mi}{1}\PY{p}{)}
    \PY{k}{return} \PY{n}{subset}\PY{p}{[}\PY{p}{:}\PY{n}{split\PYZus{}idx}\PY{p}{]}\PY{p}{,} \PY{n}{subset}\PY{p}{[}\PY{n}{split\PYZus{}idx}\PY{p}{:}\PY{p}{]}
\end{Verbatim}
\end{tcolorbox}

    \subsubsection{4.2 Fuerza bruta (tarea 0)}\label{fuerza-bruta-tarea-0}

    Tarea 0.1: \texttt{is\_clique()}

En esta tarea se ha implementado la función elemental que comprueba si
un conjunto de nodos forma realmente una clique. Es la base sobre la que
descansa toda la enumeración exhaustiva posterior.

    \begin{tcolorbox}[breakable, size=fbox, boxrule=1pt, pad at break*=1mm,colback=cellbackground, colframe=cellborder]
\prompt{In}{incolor}{255}{\boxspacing}
\begin{Verbatim}[commandchars=\\\{\}]
\PY{k}{def}\PY{+w}{ }\PY{n+nf}{is\PYZus{}clique}\PY{p}{(}\PY{n}{nodes}\PY{p}{,} \PY{n}{A}\PY{p}{)}\PY{p}{:}
\PY{+w}{    }\PY{l+s+sd}{\PYZdq{}\PYZdq{}\PYZdq{}}
\PY{l+s+sd}{    Check if a set of nodes forms a clique.}

\PY{l+s+sd}{    Args:}
\PY{l+s+sd}{        nodes: list of node indices}
\PY{l+s+sd}{        A: numpy array of shape (n, n) \PYZhy{} adjacency matrix}

\PY{l+s+sd}{    Returns:}
\PY{l+s+sd}{        bool: True if nodes form a clique, False otherwise}
\PY{l+s+sd}{    \PYZdq{}\PYZdq{}\PYZdq{}}
    \PY{k}{for} \PY{n}{i}\PY{p}{,} \PY{n}{j} \PY{o+ow}{in} \PY{n}{combinations}\PY{p}{(}\PY{n}{nodes}\PY{p}{,} \PY{l+m+mi}{2}\PY{p}{)}\PY{p}{:}
        \PY{k}{if} \PY{n}{A}\PY{p}{[}\PY{n}{i}\PY{p}{,} \PY{n}{j}\PY{p}{]} \PY{o}{==} \PY{l+m+mi}{0}\PY{p}{:}
            \PY{k}{return} \PY{k+kc}{False}
    \PY{k}{return} \PY{k+kc}{True}
\end{Verbatim}
\end{tcolorbox}

    Tarea 0.2: \texttt{brute\_force\_max\_clique()}

En esta tarea se ha implementado la búsqueda exhaustiva. El recorrido se
hace desde subconjuntos grandes a pequeños para poder detener la
búsqueda en cuanto aparece la primera clique máxima.

    \begin{tcolorbox}[breakable, size=fbox, boxrule=1pt, pad at break*=1mm,colback=cellbackground, colframe=cellborder]
\prompt{In}{incolor}{256}{\boxspacing}
\begin{Verbatim}[commandchars=\\\{\}]
\PY{k}{def}\PY{+w}{ }\PY{n+nf}{brute\PYZus{}force\PYZus{}max\PYZus{}clique}\PY{p}{(}\PY{n}{A}\PY{p}{)}\PY{p}{:}
\PY{+w}{    }\PY{l+s+sd}{\PYZdq{}\PYZdq{}\PYZdq{}}
\PY{l+s+sd}{    Find the maximum clique by brute force enumeration.}

\PY{l+s+sd}{    Args:}
\PY{l+s+sd}{        A: numpy array of shape (n, n) \PYZhy{} adjacency matrix}

\PY{l+s+sd}{    Returns:}
\PY{l+s+sd}{        list: nodes forming the maximum clique}
\PY{l+s+sd}{    \PYZdq{}\PYZdq{}\PYZdq{}}
    \PY{n}{n} \PY{o}{=} \PY{n}{A}\PY{o}{.}\PY{n}{shape}\PY{p}{[}\PY{l+m+mi}{0}\PY{p}{]}

    \PY{k}{for} \PY{n}{k} \PY{o+ow}{in} \PY{n+nb}{range}\PY{p}{(}\PY{n}{n}\PY{p}{,} \PY{l+m+mi}{0}\PY{p}{,} \PY{o}{\PYZhy{}}\PY{l+m+mi}{1}\PY{p}{)}\PY{p}{:}
        \PY{k}{for} \PY{n}{subset} \PY{o+ow}{in} \PY{n}{combinations}\PY{p}{(}\PY{n+nb}{range}\PY{p}{(}\PY{n}{n}\PY{p}{)}\PY{p}{,} \PY{n}{k}\PY{p}{)}\PY{p}{:}
            \PY{k}{if} \PY{n}{is\PYZus{}clique}\PY{p}{(}\PY{n}{subset}\PY{p}{,} \PY{n}{A}\PY{p}{)}\PY{p}{:}
                \PY{k}{return} \PY{n+nb}{list}\PY{p}{(}\PY{n}{subset}\PY{p}{)}
    \PY{k}{return} \PY{p}{[}\PY{p}{]}
\end{Verbatim}
\end{tcolorbox}

    Tarea 0.3: \texttt{test\_brute\_force()}

En esta tarea se ha preparado el ejemplo sintético del enunciado para
comprobar rápidamente que la implementación exacta devuelve una clique
válida y medir su tiempo de ejecución en un caso pequeño.

    \begin{tcolorbox}[breakable, size=fbox, boxrule=1pt, pad at break*=1mm,colback=cellbackground, colframe=cellborder]
\prompt{In}{incolor}{257}{\boxspacing}
\begin{Verbatim}[commandchars=\\\{\}]
\PY{k}{def}\PY{+w}{ }\PY{n+nf}{test\PYZus{}brute\PYZus{}force}\PY{p}{(}\PY{p}{)}\PY{p}{:}
\PY{+w}{    }\PY{l+s+sd}{\PYZdq{}\PYZdq{}\PYZdq{}}
\PY{l+s+sd}{    Test brute force on small graphs and measure computation time.}
\PY{l+s+sd}{    \PYZdq{}\PYZdq{}\PYZdq{}}
    \PY{n}{A} \PY{o}{=} \PY{n}{np}\PY{o}{.}\PY{n}{array}\PY{p}{(}\PY{p}{[}
        \PY{p}{[}\PY{l+m+mi}{0}\PY{p}{,} \PY{l+m+mi}{1}\PY{p}{,} \PY{l+m+mi}{1}\PY{p}{,} \PY{l+m+mi}{0}\PY{p}{]}\PY{p}{,}
        \PY{p}{[}\PY{l+m+mi}{1}\PY{p}{,} \PY{l+m+mi}{0}\PY{p}{,} \PY{l+m+mi}{1}\PY{p}{,} \PY{l+m+mi}{0}\PY{p}{]}\PY{p}{,}
        \PY{p}{[}\PY{l+m+mi}{1}\PY{p}{,} \PY{l+m+mi}{1}\PY{p}{,} \PY{l+m+mi}{0}\PY{p}{,} \PY{l+m+mi}{1}\PY{p}{]}\PY{p}{,}
        \PY{p}{[}\PY{l+m+mi}{0}\PY{p}{,} \PY{l+m+mi}{0}\PY{p}{,} \PY{l+m+mi}{1}\PY{p}{,} \PY{l+m+mi}{0}\PY{p}{]}
    \PY{p}{]}\PY{p}{,} \PY{n}{dtype}\PY{o}{=}\PY{n+nb}{float}\PY{p}{)}

    \PY{n}{start} \PY{o}{=} \PY{n}{time}\PY{o}{.}\PY{n}{time}\PY{p}{(}\PY{p}{)}
    \PY{n}{clique} \PY{o}{=} \PY{n}{brute\PYZus{}force\PYZus{}max\PYZus{}clique}\PY{p}{(}\PY{n}{A}\PY{p}{)}
    \PY{n}{elapsed} \PY{o}{=} \PY{n}{time}\PY{o}{.}\PY{n}{time}\PY{p}{(}\PY{p}{)} \PY{o}{\PYZhy{}} \PY{n}{start}

    \PY{n}{df} \PY{o}{=} \PY{n}{pd}\PY{o}{.}\PY{n}{DataFrame}\PY{p}{(}\PY{p}{\PYZob{}}
        \PY{l+s+s2}{\PYZdq{}}\PY{l+s+s2}{Maximum clique:}\PY{l+s+s2}{\PYZdq{}}\PY{p}{:} \PY{p}{[}\PY{n}{clique}\PY{p}{]}\PY{p}{,}
        \PY{l+s+s2}{\PYZdq{}}\PY{l+s+s2}{Clique size:}\PY{l+s+s2}{\PYZdq{}}\PY{p}{:} \PY{p}{[}\PY{n+nb}{len}\PY{p}{(}\PY{n}{clique}\PY{p}{)}\PY{p}{]}\PY{p}{,}
        \PY{l+s+s2}{\PYZdq{}}\PY{l+s+s2}{Time:}\PY{l+s+s2}{\PYZdq{}}\PY{p}{:} \PY{p}{[}\PY{n}{elapsed}\PY{p}{]}\PY{p}{,}
        \PY{l+s+s2}{\PYZdq{}}\PY{l+s+s2}{is valid clique:}\PY{l+s+s2}{\PYZdq{}}\PY{p}{:} \PY{p}{[}\PY{n}{is\PYZus{}clique}\PY{p}{(}\PY{n}{clique}\PY{p}{,} \PY{n}{A}\PY{p}{)}\PY{p}{]}\PY{p}{,}
    \PY{p}{\PYZcb{}}\PY{p}{)}
    

    \PY{k}{return} \PY{n}{df}
\end{Verbatim}
\end{tcolorbox}

    Tarea 0.4: evaluación sobre dataset

Aquí se ha preparado la evaluación por fuerza bruta sobre un subconjunto
pequeño del dataset, omitiendo los grafos con demasiados nodos para
evitar tiempos prohibitivos.

    \begin{tcolorbox}[breakable, size=fbox, boxrule=1pt, pad at break*=1mm,colback=cellbackground, colframe=cellborder]
\prompt{In}{incolor}{258}{\boxspacing}
\begin{Verbatim}[commandchars=\\\{\}]
\PY{k}{def}\PY{+w}{ }\PY{n+nf}{evaluate\PYZus{}brute\PYZus{}force\PYZus{}on\PYZus{}dataset}\PY{p}{(}\PY{n}{dataset}\PY{p}{,} \PY{n}{max\PYZus{}graphs}\PY{o}{=}\PY{l+m+mi}{10}\PY{p}{,} \PY{n}{max\PYZus{}nodes}\PY{o}{=}\PY{l+m+mi}{25}\PY{p}{)}\PY{p}{:}
\PY{+w}{    }\PY{l+s+sd}{\PYZdq{}\PYZdq{}\PYZdq{}}
\PY{l+s+sd}{    Evaluate brute force on dataset (only small graphs).}

\PY{l+s+sd}{    Args:}
\PY{l+s+sd}{        dataset: PyTorch Geometric dataset}
\PY{l+s+sd}{        max\PYZus{}graphs: maximum number of graphs to evaluate}
\PY{l+s+sd}{        max\PYZus{}nodes: skip graphs larger than this (too slow)}

\PY{l+s+sd}{    Returns:}
\PY{l+s+sd}{        results: list of (graph\PYZus{}idx, clique\PYZus{}size, time) tuples}
\PY{l+s+sd}{    \PYZdq{}\PYZdq{}\PYZdq{}}
    \PY{n}{results} \PY{o}{=} \PY{p}{[}\PY{p}{]}

    \PY{k}{for} \PY{n}{i} \PY{o+ow}{in} \PY{n+nb}{range}\PY{p}{(}\PY{n+nb}{min}\PY{p}{(}\PY{n}{max\PYZus{}graphs}\PY{p}{,} \PY{n+nb}{len}\PY{p}{(}\PY{n}{dataset}\PY{p}{)}\PY{p}{)}\PY{p}{)}\PY{p}{:}
        \PY{n}{data} \PY{o}{=} \PY{n}{dataset}\PY{p}{[}\PY{n}{i}\PY{p}{]}

        \PY{k}{if} \PY{n}{data}\PY{o}{.}\PY{n}{num\PYZus{}nodes} \PY{o}{\PYZgt{}} \PY{n}{max\PYZus{}nodes}\PY{p}{:}
            \PY{n+nb}{print}\PY{p}{(}\PY{l+s+sa}{f}\PY{l+s+s2}{\PYZdq{}}\PY{l+s+s2}{Graph }\PY{l+s+si}{\PYZob{}}\PY{n}{i}\PY{l+s+si}{\PYZcb{}}\PY{l+s+s2}{: Skipping (too large: }\PY{l+s+si}{\PYZob{}}\PY{n}{data}\PY{o}{.}\PY{n}{num\PYZus{}nodes}\PY{l+s+si}{\PYZcb{}}\PY{l+s+s2}{ nodes)}\PY{l+s+s2}{\PYZdq{}}\PY{p}{)}
            \PY{k}{continue}

        \PY{n}{A} \PY{o}{=} \PY{n}{pyg\PYZus{}to\PYZus{}adjacency}\PY{p}{(}\PY{n}{data}\PY{p}{)}
        \PY{n}{clique}\PY{p}{,} \PY{n}{elapsed} \PY{o}{=} \PY{n}{timed\PYZus{}call}\PY{p}{(}\PY{n}{brute\PYZus{}force\PYZus{}max\PYZus{}clique}\PY{p}{,} \PY{n}{A}\PY{p}{)}

        \PY{n}{results}\PY{o}{.}\PY{n}{append}\PY{p}{(}\PY{p}{\PYZob{}}
            \PY{l+s+s2}{\PYZdq{}}\PY{l+s+s2}{graph\PYZus{}idx}\PY{l+s+s2}{\PYZdq{}}\PY{p}{:} \PY{n}{i}\PY{p}{,}
            \PY{l+s+s2}{\PYZdq{}}\PY{l+s+s2}{num\PYZus{}nodes}\PY{l+s+s2}{\PYZdq{}}\PY{p}{:} \PY{n+nb}{int}\PY{p}{(}\PY{n}{data}\PY{o}{.}\PY{n}{num\PYZus{}nodes}\PY{p}{)}\PY{p}{,}
            \PY{l+s+s2}{\PYZdq{}}\PY{l+s+s2}{clique\PYZus{}size}\PY{l+s+s2}{\PYZdq{}}\PY{p}{:} \PY{n+nb}{int}\PY{p}{(}\PY{n+nb}{len}\PY{p}{(}\PY{n}{clique}\PY{p}{)}\PY{p}{)}\PY{p}{,}
            \PY{l+s+s2}{\PYZdq{}}\PY{l+s+s2}{clique}\PY{l+s+s2}{\PYZdq{}}\PY{p}{:} \PY{n}{clique}\PY{p}{,}
            \PY{l+s+s2}{\PYZdq{}}\PY{l+s+s2}{valid}\PY{l+s+s2}{\PYZdq{}}\PY{p}{:} \PY{n+nb}{bool}\PY{p}{(}\PY{n}{is\PYZus{}clique}\PY{p}{(}\PY{n}{clique}\PY{p}{,} \PY{n}{A}\PY{p}{)}\PY{p}{)}\PY{p}{,}
            \PY{l+s+s2}{\PYZdq{}}\PY{l+s+s2}{time}\PY{l+s+s2}{\PYZdq{}}\PY{p}{:} \PY{l+s+sa}{f}\PY{l+s+s2}{\PYZdq{}}\PY{l+s+si}{\PYZob{}}\PY{n}{elapsed}\PY{l+s+si}{:}\PY{l+s+s2}{.6f}\PY{l+s+si}{\PYZcb{}}\PY{l+s+s2}{\PYZdq{}}\PY{p}{,}
        \PY{p}{\PYZcb{}}\PY{p}{)}

    
    \PY{k}{return} \PY{n}{pd}\PY{o}{.}\PY{n}{DataFrame}\PY{p}{(}\PY{n}{results}\PY{p}{)}
\end{Verbatim}
\end{tcolorbox}

    \subsubsection{4.3 Formulación de
Motzkin-Straus}\label{formulaciuxf3n-de-motzkin-straus}

    Tarea 1.1: funciones objetivo

En esta tarea se han implementado la función objetivo original y su
versión regularizada. La segunda es la que permite separar soluciones
espurias en grafos como el cherry.

    \begin{tcolorbox}[breakable, size=fbox, boxrule=1pt, pad at break*=1mm,colback=cellbackground, colframe=cellborder]
\prompt{In}{incolor}{259}{\boxspacing}
\begin{Verbatim}[commandchars=\\\{\}]
\PY{k}{def}\PY{+w}{ }\PY{n+nf}{motzkin\PYZus{}straus\PYZus{}objective}\PY{p}{(}\PY{n}{x}\PY{p}{,} \PY{n}{A}\PY{p}{)}\PY{p}{:}
\PY{+w}{    }\PY{l+s+sd}{\PYZdq{}\PYZdq{}\PYZdq{}}
\PY{l+s+sd}{    Compute the Motzkin\PYZhy{}Straus objective f(x) = x\PYZca{}T A x}

\PY{l+s+sd}{    Args:}
\PY{l+s+sd}{        x: numpy array of shape (n,) \PYZhy{} point in the simplex}
\PY{l+s+sd}{        A: numpy array of shape (n, n) \PYZhy{} adjacency matrix}

\PY{l+s+sd}{    Returns:}
\PY{l+s+sd}{        float: objective value}
\PY{l+s+sd}{    \PYZdq{}\PYZdq{}\PYZdq{}}
    \PY{k}{return} \PY{n+nb}{float}\PY{p}{(}\PY{n}{x} \PY{o}{@} \PY{n}{A} \PY{o}{@} \PY{n}{x}\PY{p}{)}


\PY{k}{def}\PY{+w}{ }\PY{n+nf}{regularized\PYZus{}objective}\PY{p}{(}\PY{n}{x}\PY{p}{,} \PY{n}{A}\PY{p}{)}\PY{p}{:}
\PY{+w}{    }\PY{l+s+sd}{\PYZdq{}\PYZdq{}\PYZdq{}}
\PY{l+s+sd}{    Compute the regularized objective f\PYZus{}hat(x) = x\PYZca{}T (A + 0.5*I) x}

\PY{l+s+sd}{    Args:}
\PY{l+s+sd}{        x: numpy array of shape (n,) \PYZhy{} point in the simplex}
\PY{l+s+sd}{        A: numpy array of shape (n, n) \PYZhy{} adjacency matrix}

\PY{l+s+sd}{    Returns:}
\PY{l+s+sd}{        float: regularized objective value}
\PY{l+s+sd}{    \PYZdq{}\PYZdq{}\PYZdq{}}
    \PY{k}{return} \PY{n+nb}{float}\PY{p}{(}\PY{n}{x} \PY{o}{@} \PY{p}{(}\PY{n}{A} \PY{o}{+} \PY{l+m+mf}{0.5} \PY{o}{*} \PY{n}{np}\PY{o}{.}\PY{n}{eye}\PY{p}{(}\PY{n}{A}\PY{o}{.}\PY{n}{shape}\PY{p}{[}\PY{l+m+mi}{0}\PY{p}{]}\PY{p}{)}\PY{p}{)} \PY{o}{@} \PY{n}{x}\PY{p}{)}
\end{Verbatim}
\end{tcolorbox}

    Tarea 1.2: verificación del teorema

En esta tarea se ha reunido en una sola función la comprobación numérica
de los ejemplos del enunciado: triángulo, cuadrado con diagonal y
cherry.

    \begin{tcolorbox}[breakable, size=fbox, boxrule=1pt, pad at break*=1mm,colback=cellbackground, colframe=cellborder]
\prompt{In}{incolor}{260}{\boxspacing}
\begin{Verbatim}[commandchars=\\\{\}]
\PY{k}{def}\PY{+w}{ }\PY{n+nf}{verify\PYZus{}motzkin\PYZus{}straus}\PY{p}{(}\PY{p}{)}\PY{p}{:}
\PY{+w}{    }\PY{l+s+sd}{\PYZdq{}\PYZdq{}\PYZdq{}}
\PY{l+s+sd}{    Verify the Motzkin\PYZhy{}Straus theorem on test graphs.}
\PY{l+s+sd}{    \PYZdq{}\PYZdq{}\PYZdq{}}
    \PY{n}{A1} \PY{o}{=} \PY{n}{np}\PY{o}{.}\PY{n}{array}\PY{p}{(}\PY{p}{[}
        \PY{p}{[}\PY{l+m+mi}{0}\PY{p}{,} \PY{l+m+mi}{1}\PY{p}{,} \PY{l+m+mi}{1}\PY{p}{]}\PY{p}{,}
        \PY{p}{[}\PY{l+m+mi}{1}\PY{p}{,} \PY{l+m+mi}{0}\PY{p}{,} \PY{l+m+mi}{1}\PY{p}{]}\PY{p}{,}
        \PY{p}{[}\PY{l+m+mi}{1}\PY{p}{,} \PY{l+m+mi}{1}\PY{p}{,} \PY{l+m+mi}{0}\PY{p}{]}
    \PY{p}{]}\PY{p}{,} \PY{n}{dtype}\PY{o}{=}\PY{n+nb}{float}\PY{p}{)}

    \PY{n}{clique1} \PY{o}{=} \PY{n}{brute\PYZus{}force\PYZus{}max\PYZus{}clique}\PY{p}{(}\PY{n}{A1}\PY{p}{)}
    \PY{n}{x1} \PY{o}{=} \PY{n}{clique\PYZus{}to\PYZus{}simplex}\PY{p}{(}\PY{n}{clique1}\PY{p}{,} \PY{n}{A1}\PY{o}{.}\PY{n}{shape}\PY{p}{[}\PY{l+m+mi}{0}\PY{p}{]}\PY{p}{)}

    \PY{n}{A2} \PY{o}{=} \PY{n}{np}\PY{o}{.}\PY{n}{array}\PY{p}{(}\PY{p}{[}
        \PY{p}{[}\PY{l+m+mi}{0}\PY{p}{,} \PY{l+m+mi}{1}\PY{p}{,} \PY{l+m+mi}{1}\PY{p}{,} \PY{l+m+mi}{1}\PY{p}{]}\PY{p}{,}
        \PY{p}{[}\PY{l+m+mi}{1}\PY{p}{,} \PY{l+m+mi}{0}\PY{p}{,} \PY{l+m+mi}{1}\PY{p}{,} \PY{l+m+mi}{0}\PY{p}{]}\PY{p}{,}
        \PY{p}{[}\PY{l+m+mi}{1}\PY{p}{,} \PY{l+m+mi}{1}\PY{p}{,} \PY{l+m+mi}{0}\PY{p}{,} \PY{l+m+mi}{1}\PY{p}{]}\PY{p}{,}
        \PY{p}{[}\PY{l+m+mi}{1}\PY{p}{,} \PY{l+m+mi}{0}\PY{p}{,} \PY{l+m+mi}{1}\PY{p}{,} \PY{l+m+mi}{0}\PY{p}{]}
    \PY{p}{]}\PY{p}{,} \PY{n}{dtype}\PY{o}{=}\PY{n+nb}{float}\PY{p}{)}

    \PY{n}{clique2} \PY{o}{=} \PY{n}{brute\PYZus{}force\PYZus{}max\PYZus{}clique}\PY{p}{(}\PY{n}{A2}\PY{p}{)}
    \PY{n}{x2} \PY{o}{=} \PY{n}{clique\PYZus{}to\PYZus{}simplex}\PY{p}{(}\PY{n}{clique2}\PY{p}{,} \PY{n}{A2}\PY{o}{.}\PY{n}{shape}\PY{p}{[}\PY{l+m+mi}{0}\PY{p}{]}\PY{p}{)}

    \PY{n}{A3} \PY{o}{=} \PY{n}{np}\PY{o}{.}\PY{n}{array}\PY{p}{(}\PY{p}{[}
        \PY{p}{[}\PY{l+m+mi}{0}\PY{p}{,} \PY{l+m+mi}{0}\PY{p}{,} \PY{l+m+mi}{1}\PY{p}{]}\PY{p}{,}
        \PY{p}{[}\PY{l+m+mi}{0}\PY{p}{,} \PY{l+m+mi}{0}\PY{p}{,} \PY{l+m+mi}{1}\PY{p}{]}\PY{p}{,}
        \PY{p}{[}\PY{l+m+mi}{1}\PY{p}{,} \PY{l+m+mi}{1}\PY{p}{,} \PY{l+m+mi}{0}\PY{p}{]}
    \PY{p}{]}\PY{p}{,} \PY{n}{dtype}\PY{o}{=}\PY{n+nb}{float}\PY{p}{)}

    \PY{n}{true\PYZus{}x} \PY{o}{=} \PY{n}{np}\PY{o}{.}\PY{n}{array}\PY{p}{(}\PY{p}{[}\PY{l+m+mf}{0.5}\PY{p}{,} \PY{l+m+mf}{0.0}\PY{p}{,} \PY{l+m+mf}{0.5}\PY{p}{]}\PY{p}{,} \PY{n}{dtype}\PY{o}{=}\PY{n+nb}{float}\PY{p}{)}
    \PY{n}{spurious\PYZus{}x} \PY{o}{=} \PY{n}{np}\PY{o}{.}\PY{n}{array}\PY{p}{(}\PY{p}{[}\PY{l+m+mf}{0.25}\PY{p}{,} \PY{l+m+mf}{0.25}\PY{p}{,} \PY{l+m+mf}{0.5}\PY{p}{]}\PY{p}{,} \PY{n}{dtype}\PY{o}{=}\PY{n+nb}{float}\PY{p}{)}

    \PY{k}{return} \PY{n}{pd}\PY{o}{.}\PY{n}{DataFrame}\PY{p}{(}\PY{p}{[}
        \PY{p}{\PYZob{}}
            \PY{l+s+s2}{\PYZdq{}}\PY{l+s+s2}{graph}\PY{l+s+s2}{\PYZdq{}}\PY{p}{:} \PY{l+s+s2}{\PYZdq{}}\PY{l+s+s2}{Triángulo K3}\PY{l+s+s2}{\PYZdq{}}\PY{p}{,}
            \PY{l+s+s2}{\PYZdq{}}\PY{l+s+s2}{omega}\PY{l+s+s2}{\PYZdq{}}\PY{p}{:} \PY{n+nb}{len}\PY{p}{(}\PY{n}{clique1}\PY{p}{)}\PY{p}{,}
            \PY{l+s+s2}{\PYZdq{}}\PY{l+s+s2}{f\PYZus{}x}\PY{l+s+s2}{\PYZdq{}}\PY{p}{:} \PY{n}{motzkin\PYZus{}straus\PYZus{}objective}\PY{p}{(}\PY{n}{x1}\PY{p}{,} \PY{n}{A1}\PY{p}{)}\PY{p}{,}
            \PY{l+s+s2}{\PYZdq{}}\PY{l+s+s2}{theoretical}\PY{l+s+s2}{\PYZdq{}}\PY{p}{:} \PY{l+m+mi}{1} \PY{o}{\PYZhy{}} \PY{l+m+mi}{1} \PY{o}{/} \PY{n+nb}{len}\PY{p}{(}\PY{n}{clique1}\PY{p}{)}\PY{p}{,}
            \PY{l+s+s2}{\PYZdq{}}\PY{l+s+s2}{regularized\PYZus{}f\PYZus{}x}\PY{l+s+s2}{\PYZdq{}}\PY{p}{:} \PY{n}{regularized\PYZus{}objective}\PY{p}{(}\PY{n}{x1}\PY{p}{,} \PY{n}{A1}\PY{p}{)}\PY{p}{,}
        \PY{p}{\PYZcb{}}\PY{p}{,}
        \PY{p}{\PYZob{}}
            \PY{l+s+s2}{\PYZdq{}}\PY{l+s+s2}{graph}\PY{l+s+s2}{\PYZdq{}}\PY{p}{:} \PY{l+s+s2}{\PYZdq{}}\PY{l+s+s2}{Cuadrado con diagonal}\PY{l+s+s2}{\PYZdq{}}\PY{p}{,}
            \PY{l+s+s2}{\PYZdq{}}\PY{l+s+s2}{omega}\PY{l+s+s2}{\PYZdq{}}\PY{p}{:} \PY{n+nb}{len}\PY{p}{(}\PY{n}{clique2}\PY{p}{)}\PY{p}{,}
            \PY{l+s+s2}{\PYZdq{}}\PY{l+s+s2}{f\PYZus{}x}\PY{l+s+s2}{\PYZdq{}}\PY{p}{:} \PY{n}{motzkin\PYZus{}straus\PYZus{}objective}\PY{p}{(}\PY{n}{x2}\PY{p}{,} \PY{n}{A2}\PY{p}{)}\PY{p}{,}
            \PY{l+s+s2}{\PYZdq{}}\PY{l+s+s2}{theoretical}\PY{l+s+s2}{\PYZdq{}}\PY{p}{:} \PY{l+m+mi}{1} \PY{o}{\PYZhy{}} \PY{l+m+mi}{1} \PY{o}{/} \PY{n+nb}{len}\PY{p}{(}\PY{n}{clique2}\PY{p}{)}\PY{p}{,}
            \PY{l+s+s2}{\PYZdq{}}\PY{l+s+s2}{regularized\PYZus{}f\PYZus{}x}\PY{l+s+s2}{\PYZdq{}}\PY{p}{:} \PY{n}{regularized\PYZus{}objective}\PY{p}{(}\PY{n}{x2}\PY{p}{,} \PY{n}{A2}\PY{p}{)}\PY{p}{,}
        \PY{p}{\PYZcb{}}\PY{p}{,}
        \PY{p}{\PYZob{}}
            \PY{l+s+s2}{\PYZdq{}}\PY{l+s+s2}{graph}\PY{l+s+s2}{\PYZdq{}}\PY{p}{:} \PY{l+s+s2}{\PYZdq{}}\PY{l+s+s2}{Cherry \PYZhy{} solución verdadera}\PY{l+s+s2}{\PYZdq{}}\PY{p}{,}
            \PY{l+s+s2}{\PYZdq{}}\PY{l+s+s2}{omega}\PY{l+s+s2}{\PYZdq{}}\PY{p}{:} \PY{l+m+mi}{2}\PY{p}{,}
            \PY{l+s+s2}{\PYZdq{}}\PY{l+s+s2}{f\PYZus{}x}\PY{l+s+s2}{\PYZdq{}}\PY{p}{:} \PY{n}{motzkin\PYZus{}straus\PYZus{}objective}\PY{p}{(}\PY{n}{true\PYZus{}x}\PY{p}{,} \PY{n}{A3}\PY{p}{)}\PY{p}{,}
            \PY{l+s+s2}{\PYZdq{}}\PY{l+s+s2}{theoretical}\PY{l+s+s2}{\PYZdq{}}\PY{p}{:} \PY{l+m+mf}{0.5}\PY{p}{,}
            \PY{l+s+s2}{\PYZdq{}}\PY{l+s+s2}{regularized\PYZus{}f\PYZus{}x}\PY{l+s+s2}{\PYZdq{}}\PY{p}{:} \PY{n}{regularized\PYZus{}objective}\PY{p}{(}\PY{n}{true\PYZus{}x}\PY{p}{,} \PY{n}{A3}\PY{p}{)}\PY{p}{,}
        \PY{p}{\PYZcb{}}\PY{p}{,}
        \PY{p}{\PYZob{}}
            \PY{l+s+s2}{\PYZdq{}}\PY{l+s+s2}{graph}\PY{l+s+s2}{\PYZdq{}}\PY{p}{:} \PY{l+s+s2}{\PYZdq{}}\PY{l+s+s2}{Cherry \PYZhy{} solución espuria}\PY{l+s+s2}{\PYZdq{}}\PY{p}{,}
            \PY{l+s+s2}{\PYZdq{}}\PY{l+s+s2}{omega}\PY{l+s+s2}{\PYZdq{}}\PY{p}{:} \PY{l+m+mi}{2}\PY{p}{,}
            \PY{l+s+s2}{\PYZdq{}}\PY{l+s+s2}{f\PYZus{}x}\PY{l+s+s2}{\PYZdq{}}\PY{p}{:} \PY{n}{motzkin\PYZus{}straus\PYZus{}objective}\PY{p}{(}\PY{n}{spurious\PYZus{}x}\PY{p}{,} \PY{n}{A3}\PY{p}{)}\PY{p}{,}
            \PY{l+s+s2}{\PYZdq{}}\PY{l+s+s2}{theoretical}\PY{l+s+s2}{\PYZdq{}}\PY{p}{:} \PY{l+m+mf}{0.5}\PY{p}{,}
            \PY{l+s+s2}{\PYZdq{}}\PY{l+s+s2}{regularized\PYZus{}f\PYZus{}x}\PY{l+s+s2}{\PYZdq{}}\PY{p}{:} \PY{n}{regularized\PYZus{}objective}\PY{p}{(}\PY{n}{spurious\PYZus{}x}\PY{p}{,} \PY{n}{A3}\PY{p}{)}\PY{p}{,}
        \PY{p}{\PYZcb{}}\PY{p}{,}
    \PY{p}{]}\PY{p}{)}
\end{Verbatim}
\end{tcolorbox}

    \subsubsection{4.4 Dinámicas
replicadoras}\label{dinuxe1micas-replicadoras}

    Tarea 2.1: \texttt{replicator\_dynamics()}

En esta tarea se ha implementado la dinámica multiplicativa de Pelillo
sobre la matriz regularizada. La función devuelve tanto el punto final
del simplex como la historia del objetivo para estudiar la convergencia.

    \begin{tcolorbox}[breakable, size=fbox, boxrule=1pt, pad at break*=1mm,colback=cellbackground, colframe=cellborder]
\prompt{In}{incolor}{261}{\boxspacing}
\begin{Verbatim}[commandchars=\\\{\}]
\PY{k}{def}\PY{+w}{ }\PY{n+nf}{replicator\PYZus{}dynamics}\PY{p}{(}\PY{n}{A}\PY{p}{,} \PY{n}{max\PYZus{}iter}\PY{o}{=}\PY{l+m+mi}{1000}\PY{p}{,} \PY{n}{tol}\PY{o}{=}\PY{l+m+mf}{1e\PYZhy{}12}\PY{p}{,} \PY{n}{use\PYZus{}regularization}\PY{o}{=}\PY{k+kc}{True}\PY{p}{)}\PY{p}{:}
\PY{+w}{    }\PY{l+s+sd}{\PYZdq{}\PYZdq{}\PYZdq{}}
\PY{l+s+sd}{    Run Replicator Dynamics to find the maximum clique.}

\PY{l+s+sd}{    This follows the formulation from Pelillo et al.:}
\PY{l+s+sd}{    x\PYZus{}new = x * (A @ x)}
\PY{l+s+sd}{    x\PYZus{}new = x\PYZus{}new / x\PYZus{}new.sum()}

\PY{l+s+sd}{    Args:}
\PY{l+s+sd}{        A: numpy array of shape (n, n) \PYZhy{} adjacency matrix}
\PY{l+s+sd}{        max\PYZus{}iter: int \PYZhy{} maximum number of iterations}
\PY{l+s+sd}{        tol: float \PYZhy{} convergence tolerance}
\PY{l+s+sd}{        use\PYZus{}regularization: bool \PYZhy{} whether to use A + 0.5*I}

\PY{l+s+sd}{    Returns:}
\PY{l+s+sd}{        x: numpy array of shape (n,) \PYZhy{} final simplex point}
\PY{l+s+sd}{        history: dict with \PYZsq{}objectives\PYZsq{} and \PYZsq{}iterations\PYZsq{}}
\PY{l+s+sd}{    \PYZdq{}\PYZdq{}\PYZdq{}}
    \PY{n}{n} \PY{o}{=} \PY{n}{A}\PY{o}{.}\PY{n}{shape}\PY{p}{[}\PY{l+m+mi}{0}\PY{p}{]}

    \PY{k}{if} \PY{n}{use\PYZus{}regularization}\PY{p}{:}
        \PY{n}{A\PYZus{}reg} \PY{o}{=} \PY{n}{A} \PY{o}{+} \PY{l+m+mf}{0.5} \PY{o}{*} \PY{n}{np}\PY{o}{.}\PY{n}{eye}\PY{p}{(}\PY{n}{n}\PY{p}{)}
    \PY{k}{else}\PY{p}{:}
        \PY{n}{A\PYZus{}reg} \PY{o}{=} \PY{n}{A}\PY{o}{.}\PY{n}{copy}\PY{p}{(}\PY{p}{)}

    \PY{n}{x} \PY{o}{=} \PY{n}{np}\PY{o}{.}\PY{n}{ones}\PY{p}{(}\PY{n}{n}\PY{p}{)} \PY{o}{/} \PY{n}{n}

    \PY{n}{history} \PY{o}{=} \PY{p}{\PYZob{}}
        \PY{l+s+s1}{\PYZsq{}}\PY{l+s+s1}{objectives}\PY{l+s+s1}{\PYZsq{}}\PY{p}{:} \PY{p}{[}\PY{n}{x} \PY{o}{@} \PY{n}{A\PYZus{}reg} \PY{o}{@} \PY{n}{x}\PY{p}{]}\PY{p}{,}
        \PY{l+s+s1}{\PYZsq{}}\PY{l+s+s1}{iterations}\PY{l+s+s1}{\PYZsq{}}\PY{p}{:} \PY{l+m+mi}{0}
    \PY{p}{\PYZcb{}}

    \PY{k}{for} \PY{n}{t} \PY{o+ow}{in} \PY{n+nb}{range}\PY{p}{(}\PY{n}{max\PYZus{}iter}\PY{p}{)}\PY{p}{:}
        \PY{n}{x\PYZus{}old} \PY{o}{=} \PY{n}{x}\PY{o}{.}\PY{n}{copy}\PY{p}{(}\PY{p}{)}
        \PY{n}{x} \PY{o}{=} \PY{n}{x} \PY{o}{*} \PY{p}{(}\PY{n}{A\PYZus{}reg} \PY{o}{@} \PY{n}{x}\PY{p}{)}
        \PY{n}{x} \PY{o}{=} \PY{n}{x} \PY{o}{/} \PY{n+nb}{max}\PY{p}{(}\PY{n}{x}\PY{o}{.}\PY{n}{sum}\PY{p}{(}\PY{p}{)}\PY{p}{,} \PY{l+m+mf}{1e\PYZhy{}12}\PY{p}{)}
        \PY{n}{history}\PY{p}{[}\PY{l+s+s1}{\PYZsq{}}\PY{l+s+s1}{objectives}\PY{l+s+s1}{\PYZsq{}}\PY{p}{]}\PY{o}{.}\PY{n}{append}\PY{p}{(}\PY{n+nb}{float}\PY{p}{(}\PY{n}{x} \PY{o}{@} \PY{n}{A\PYZus{}reg} \PY{o}{@} \PY{n}{x}\PY{p}{)}\PY{p}{)}
        \PY{n}{history}\PY{p}{[}\PY{l+s+s1}{\PYZsq{}}\PY{l+s+s1}{iterations}\PY{l+s+s1}{\PYZsq{}}\PY{p}{]} \PY{o}{=} \PY{n}{t} \PY{o}{+} \PY{l+m+mi}{1}
        \PY{k}{if} \PY{n}{np}\PY{o}{.}\PY{n}{linalg}\PY{o}{.}\PY{n}{norm}\PY{p}{(}\PY{n}{x} \PY{o}{\PYZhy{}} \PY{n}{x\PYZus{}old}\PY{p}{,} \PY{n+nb}{ord}\PY{o}{=}\PY{l+m+mi}{1}\PY{p}{)} \PY{o}{\PYZlt{}} \PY{n}{tol}\PY{p}{:}
            \PY{k}{break}

    \PY{k}{return} \PY{n}{x}\PY{p}{,} \PY{n}{history}
\end{Verbatim}
\end{tcolorbox}

    Tarea 2.2: \texttt{decode\_clique()}

En esta tarea se ha implementado el decodificador voraz que transforma
la solución continua en una clique discreta válida.

    \begin{tcolorbox}[breakable, size=fbox, boxrule=1pt, pad at break*=1mm,colback=cellbackground, colframe=cellborder]
\prompt{In}{incolor}{262}{\boxspacing}
\begin{Verbatim}[commandchars=\\\{\}]
\PY{k}{def}\PY{+w}{ }\PY{n+nf}{decode\PYZus{}clique}\PY{p}{(}\PY{n}{x}\PY{p}{,} \PY{n}{A}\PY{p}{,} \PY{n}{threshold}\PY{o}{=}\PY{l+m+mf}{1e\PYZhy{}6}\PY{p}{)}\PY{p}{:}
\PY{+w}{    }\PY{l+s+sd}{\PYZdq{}\PYZdq{}\PYZdq{}}
\PY{l+s+sd}{    Extract the maximum clique from the RD solution.}

\PY{l+s+sd}{    Strategy: }
\PY{l+s+sd}{    1. Sort nodes by x value (descending)}
\PY{l+s+sd}{    2. Greedily add nodes that maintain clique property}

\PY{l+s+sd}{    Args:}
\PY{l+s+sd}{        x: numpy array of shape (n,) \PYZhy{} RD solution}
\PY{l+s+sd}{        A: numpy array of shape (n, n) \PYZhy{} adjacency matrix}
\PY{l+s+sd}{        threshold: float \PYZhy{} minimum x value to consider}

\PY{l+s+sd}{    Returns:}
\PY{l+s+sd}{        clique: list of node indices}
\PY{l+s+sd}{    \PYZdq{}\PYZdq{}\PYZdq{}}
    \PY{n}{indices} \PY{o}{=} \PY{n}{np}\PY{o}{.}\PY{n}{argsort}\PY{p}{(}\PY{o}{\PYZhy{}}\PY{n}{x}\PY{p}{)}
    \PY{n}{clique} \PY{o}{=} \PY{p}{[}\PY{p}{]}

    \PY{k}{for} \PY{n}{i} \PY{o+ow}{in} \PY{n}{indices}\PY{p}{:}
        \PY{k}{if} \PY{n}{x}\PY{p}{[}\PY{n}{i}\PY{p}{]} \PY{o}{\PYZlt{}} \PY{n}{threshold} \PY{o+ow}{and} \PY{n+nb}{len}\PY{p}{(}\PY{n}{clique}\PY{p}{)} \PY{o}{\PYZgt{}} \PY{l+m+mi}{0}\PY{p}{:}
            \PY{k}{break}
        \PY{k}{if} \PY{n+nb}{all}\PY{p}{(}\PY{n}{A}\PY{p}{[}\PY{n}{j}\PY{p}{,} \PY{n}{i}\PY{p}{]} \PY{o}{==} \PY{l+m+mi}{1} \PY{k}{for} \PY{n}{j} \PY{o+ow}{in} \PY{n}{clique}\PY{p}{)}\PY{p}{:}
            \PY{n}{clique}\PY{o}{.}\PY{n}{append}\PY{p}{(}\PY{n+nb}{int}\PY{p}{(}\PY{n}{i}\PY{p}{)}\PY{p}{)}

    \PY{k}{if} \PY{n+nb}{len}\PY{p}{(}\PY{n}{clique}\PY{p}{)} \PY{o}{==} \PY{l+m+mi}{0} \PY{o+ow}{and} \PY{n+nb}{len}\PY{p}{(}\PY{n}{indices}\PY{p}{)} \PY{o}{\PYZgt{}} \PY{l+m+mi}{0}\PY{p}{:}
        \PY{n}{clique} \PY{o}{=} \PY{p}{[}\PY{n+nb}{int}\PY{p}{(}\PY{n}{indices}\PY{p}{[}\PY{l+m+mi}{0}\PY{p}{]}\PY{p}{)}\PY{p}{]}

    \PY{k}{return} \PY{n}{clique}
\end{Verbatim}
\end{tcolorbox}

    Tarea 2.3: visualización de convergencia

En esta tarea se ha preparado una visualización del objetivo, de los
pesos finales y del grafo coloreado para interpretar mejor el
comportamiento de RD.

    \begin{tcolorbox}[breakable, size=fbox, boxrule=1pt, pad at break*=1mm,colback=cellbackground, colframe=cellborder]
\prompt{In}{incolor}{263}{\boxspacing}
\begin{Verbatim}[commandchars=\\\{\}]
\PY{k}{def}\PY{+w}{ }\PY{n+nf}{visualize\PYZus{}rd}\PY{p}{(}\PY{n}{A}\PY{p}{,} \PY{n}{graph\PYZus{}name}\PY{o}{=}\PY{l+s+s2}{\PYZdq{}}\PY{l+s+s2}{Test Graph}\PY{l+s+s2}{\PYZdq{}}\PY{p}{)}\PY{p}{:}
\PY{+w}{    }\PY{l+s+sd}{\PYZdq{}\PYZdq{}\PYZdq{}}
\PY{l+s+sd}{    Run RD and create visualization of convergence.}
\PY{l+s+sd}{    \PYZdq{}\PYZdq{}\PYZdq{}}
    \PY{n}{x\PYZus{}final}\PY{p}{,} \PY{n}{history} \PY{o}{=} \PY{n}{replicator\PYZus{}dynamics}\PY{p}{(}\PY{n}{A}\PY{p}{)}

    \PY{n}{fig}\PY{p}{,} \PY{n}{axes} \PY{o}{=} \PY{n}{plt}\PY{o}{.}\PY{n}{subplots}\PY{p}{(}\PY{l+m+mi}{1}\PY{p}{,} \PY{l+m+mi}{2}\PY{p}{,} \PY{n}{figsize}\PY{o}{=}\PY{p}{(}\PY{l+m+mi}{12}\PY{p}{,} \PY{l+m+mi}{4}\PY{p}{)}\PY{p}{)}

    \PY{n}{axes}\PY{p}{[}\PY{l+m+mi}{0}\PY{p}{]}\PY{o}{.}\PY{n}{plot}\PY{p}{(}\PY{n}{history}\PY{p}{[}\PY{l+s+s1}{\PYZsq{}}\PY{l+s+s1}{objectives}\PY{l+s+s1}{\PYZsq{}}\PY{p}{]}\PY{p}{,} \PY{n}{color}\PY{o}{=}\PY{l+s+s1}{\PYZsq{}}\PY{l+s+s1}{\PYZsh{}1f77b4}\PY{l+s+s1}{\PYZsq{}}\PY{p}{)}
    \PY{n}{axes}\PY{p}{[}\PY{l+m+mi}{0}\PY{p}{]}\PY{o}{.}\PY{n}{set\PYZus{}title}\PY{p}{(}\PY{l+s+s1}{\PYZsq{}}\PY{l+s+s1}{Objective value over iterations}\PY{l+s+s1}{\PYZsq{}}\PY{p}{)}
    \PY{n}{axes}\PY{p}{[}\PY{l+m+mi}{0}\PY{p}{]}\PY{o}{.}\PY{n}{set\PYZus{}xlabel}\PY{p}{(}\PY{l+s+s1}{\PYZsq{}}\PY{l+s+s1}{Iteration}\PY{l+s+s1}{\PYZsq{}}\PY{p}{)}
    \PY{n}{axes}\PY{p}{[}\PY{l+m+mi}{0}\PY{p}{]}\PY{o}{.}\PY{n}{set\PYZus{}ylabel}\PY{p}{(}\PY{l+s+s1}{\PYZsq{}}\PY{l+s+s1}{f(x)}\PY{l+s+s1}{\PYZsq{}}\PY{p}{)}

    \PY{n}{axes}\PY{p}{[}\PY{l+m+mi}{1}\PY{p}{]}\PY{o}{.}\PY{n}{bar}\PY{p}{(}\PY{n}{np}\PY{o}{.}\PY{n}{arange}\PY{p}{(}\PY{n+nb}{len}\PY{p}{(}\PY{n}{x\PYZus{}final}\PY{p}{)}\PY{p}{)}\PY{p}{,} \PY{n}{x\PYZus{}final}\PY{p}{,} \PY{n}{color}\PY{o}{=}\PY{l+s+s1}{\PYZsq{}}\PY{l+s+s1}{\PYZsh{}2ca02c}\PY{l+s+s1}{\PYZsq{}}\PY{p}{)}
    \PY{n}{axes}\PY{p}{[}\PY{l+m+mi}{1}\PY{p}{]}\PY{o}{.}\PY{n}{set\PYZus{}title}\PY{p}{(}\PY{l+s+s1}{\PYZsq{}}\PY{l+s+s1}{Final node weights}\PY{l+s+s1}{\PYZsq{}}\PY{p}{)}
    \PY{n}{axes}\PY{p}{[}\PY{l+m+mi}{1}\PY{p}{]}\PY{o}{.}\PY{n}{set\PYZus{}xlabel}\PY{p}{(}\PY{l+s+s1}{\PYZsq{}}\PY{l+s+s1}{Node}\PY{l+s+s1}{\PYZsq{}}\PY{p}{)}
    \PY{n}{axes}\PY{p}{[}\PY{l+m+mi}{1}\PY{p}{]}\PY{o}{.}\PY{n}{set\PYZus{}ylabel}\PY{p}{(}\PY{l+s+s1}{\PYZsq{}}\PY{l+s+s1}{x\PYZus{}i}\PY{l+s+s1}{\PYZsq{}}\PY{p}{)}

    \PY{n}{plt}\PY{o}{.}\PY{n}{suptitle}\PY{p}{(}\PY{l+s+sa}{f}\PY{l+s+s2}{\PYZdq{}}\PY{l+s+s2}{Replicator Dynamics: }\PY{l+s+si}{\PYZob{}}\PY{n}{graph\PYZus{}name}\PY{l+s+si}{\PYZcb{}}\PY{l+s+s2}{\PYZdq{}}\PY{p}{)}
    \PY{n}{plt}\PY{o}{.}\PY{n}{tight\PYZus{}layout}\PY{p}{(}\PY{p}{)}
    \PY{n}{plt}\PY{o}{.}\PY{n}{show}\PY{p}{(}\PY{p}{)}

    \PY{n}{plot\PYZus{}weighted\PYZus{}graph}\PY{p}{(}\PY{n}{A}\PY{p}{,} \PY{n}{x\PYZus{}final}\PY{p}{,} \PY{l+s+sa}{f}\PY{l+s+s2}{\PYZdq{}}\PY{l+s+si}{\PYZob{}}\PY{n}{graph\PYZus{}name}\PY{l+s+si}{\PYZcb{}}\PY{l+s+s2}{ coloreado por x final}\PY{l+s+s2}{\PYZdq{}}\PY{p}{)}
    \PY{k}{return} \PY{n}{x\PYZus{}final}\PY{p}{,} \PY{n}{history}
\end{Verbatim}
\end{tcolorbox}

    Tarea 2.4: evaluación sobre dataset

En esta tarea se ha implementado la evaluación de RD sobre un
subconjunto de grafos y una comparación adicional con la referencia
exacta allí donde existe ground truth.

    \begin{tcolorbox}[breakable, size=fbox, boxrule=1pt, pad at break*=1mm,colback=cellbackground, colframe=cellborder]
\prompt{In}{incolor}{264}{\boxspacing}
\begin{Verbatim}[commandchars=\\\{\}]
\PY{k}{def}\PY{+w}{ }\PY{n+nf}{evaluate\PYZus{}rd\PYZus{}on\PYZus{}dataset}\PY{p}{(}\PY{n}{dataset}\PY{p}{,} \PY{n}{dataset\PYZus{}name}\PY{p}{,} \PY{n}{num\PYZus{}graphs}\PY{o}{=}\PY{l+m+mi}{50}\PY{p}{)}\PY{p}{:}
\PY{+w}{    }\PY{l+s+sd}{\PYZdq{}\PYZdq{}\PYZdq{}}
\PY{l+s+sd}{    Evaluate Replicator Dynamics on a dataset.}
\PY{l+s+sd}{    \PYZdq{}\PYZdq{}\PYZdq{}}
    \PY{n}{results} \PY{o}{=} \PY{p}{[}\PY{p}{]}

    \PY{k}{for} \PY{n}{i} \PY{o+ow}{in} \PY{n+nb}{range}\PY{p}{(}\PY{n+nb}{min}\PY{p}{(}\PY{n}{num\PYZus{}graphs}\PY{p}{,} \PY{n+nb}{len}\PY{p}{(}\PY{n}{dataset}\PY{p}{)}\PY{p}{)}\PY{p}{)}\PY{p}{:}
        \PY{n}{data} \PY{o}{=} \PY{n}{dataset}\PY{p}{[}\PY{n}{i}\PY{p}{]}
        \PY{n}{A} \PY{o}{=} \PY{n}{pyg\PYZus{}to\PYZus{}adjacency}\PY{p}{(}\PY{n}{data}\PY{p}{)}

        \PY{n}{start\PYZus{}time} \PY{o}{=} \PY{n}{time}\PY{o}{.}\PY{n}{time}\PY{p}{(}\PY{p}{)}
        \PY{n}{x\PYZus{}final}\PY{p}{,} \PY{n}{history} \PY{o}{=} \PY{n}{replicator\PYZus{}dynamics}\PY{p}{(}\PY{n}{A}\PY{p}{,} \PY{n}{max\PYZus{}iter}\PY{o}{=}\PY{l+m+mi}{500}\PY{p}{)}
        \PY{n}{clique} \PY{o}{=} \PY{n}{decode\PYZus{}clique}\PY{p}{(}\PY{n}{x\PYZus{}final}\PY{p}{,} \PY{n}{A}\PY{p}{)}
        \PY{n}{elapsed} \PY{o}{=} \PY{n}{time}\PY{o}{.}\PY{n}{time}\PY{p}{(}\PY{p}{)} \PY{o}{\PYZhy{}} \PY{n}{start\PYZus{}time}

        \PY{n}{results}\PY{o}{.}\PY{n}{append}\PY{p}{(}\PY{p}{\PYZob{}}
            \PY{l+s+s1}{\PYZsq{}}\PY{l+s+s1}{graph\PYZus{}idx}\PY{l+s+s1}{\PYZsq{}}\PY{p}{:} \PY{n}{i}\PY{p}{,}
            \PY{l+s+s1}{\PYZsq{}}\PY{l+s+s1}{num\PYZus{}nodes}\PY{l+s+s1}{\PYZsq{}}\PY{p}{:} \PY{n+nb}{int}\PY{p}{(}\PY{n}{data}\PY{o}{.}\PY{n}{num\PYZus{}nodes}\PY{p}{)}\PY{p}{,}
            \PY{l+s+s1}{\PYZsq{}}\PY{l+s+s1}{clique\PYZus{}size}\PY{l+s+s1}{\PYZsq{}}\PY{p}{:} \PY{n+nb}{int}\PY{p}{(}\PY{n+nb}{len}\PY{p}{(}\PY{n}{clique}\PY{p}{)}\PY{p}{)}\PY{p}{,}
            \PY{l+s+s1}{\PYZsq{}}\PY{l+s+s1}{clique}\PY{l+s+s1}{\PYZsq{}}\PY{p}{:} \PY{n}{clique}\PY{p}{,}
            \PY{l+s+s1}{\PYZsq{}}\PY{l+s+s1}{valid}\PY{l+s+s1}{\PYZsq{}}\PY{p}{:} \PY{n+nb}{bool}\PY{p}{(}\PY{n}{is\PYZus{}clique}\PY{p}{(}\PY{n}{clique}\PY{p}{,} \PY{n}{A}\PY{p}{)}\PY{p}{)}\PY{p}{,}
            \PY{l+s+s1}{\PYZsq{}}\PY{l+s+s1}{time}\PY{l+s+s1}{\PYZsq{}}\PY{p}{:} \PY{n+nb}{float}\PY{p}{(}\PY{n}{elapsed}\PY{p}{)}\PY{p}{,}
            \PY{l+s+s1}{\PYZsq{}}\PY{l+s+s1}{iterations}\PY{l+s+s1}{\PYZsq{}}\PY{p}{:} \PY{n+nb}{int}\PY{p}{(}\PY{n}{history}\PY{p}{[}\PY{l+s+s1}{\PYZsq{}}\PY{l+s+s1}{iterations}\PY{l+s+s1}{\PYZsq{}}\PY{p}{]}\PY{p}{)}
        \PY{p}{\PYZcb{}}\PY{p}{)}

    \PY{k}{return} \PY{n}{results}


\PY{k}{def}\PY{+w}{ }\PY{n+nf}{compare\PYZus{}rd\PYZus{}with\PYZus{}bruteforce}\PY{p}{(}\PY{n}{rd\PYZus{}results}\PY{p}{,} \PY{n}{bf\PYZus{}results}\PY{p}{)}\PY{p}{:}
    \PY{k}{if} \PY{n+nb}{len}\PY{p}{(}\PY{n}{rd\PYZus{}results}\PY{p}{)} \PY{o}{==} \PY{l+m+mi}{0} \PY{o+ow}{or} \PY{n+nb}{len}\PY{p}{(}\PY{n}{bf\PYZus{}results}\PY{p}{)} \PY{o}{==} \PY{l+m+mi}{0}\PY{p}{:}
        \PY{k}{return} \PY{n}{pd}\PY{o}{.}\PY{n}{DataFrame}\PY{p}{(}\PY{p}{)}

    \PY{k}{if} \PY{n+nb}{isinstance}\PY{p}{(}\PY{n}{rd\PYZus{}results}\PY{p}{,} \PY{n}{pd}\PY{o}{.}\PY{n}{DataFrame}\PY{p}{)}\PY{p}{:}
        \PY{n}{rd\PYZus{}records} \PY{o}{=} \PY{n}{rd\PYZus{}results}\PY{o}{.}\PY{n}{to\PYZus{}dict}\PY{p}{(}\PY{l+s+s1}{\PYZsq{}}\PY{l+s+s1}{records}\PY{l+s+s1}{\PYZsq{}}\PY{p}{)}
    \PY{k}{else}\PY{p}{:}
        \PY{n}{rd\PYZus{}records} \PY{o}{=} \PY{n}{rd\PYZus{}results}

    \PY{k}{if} \PY{n+nb}{isinstance}\PY{p}{(}\PY{n}{bf\PYZus{}results}\PY{p}{,} \PY{n}{pd}\PY{o}{.}\PY{n}{DataFrame}\PY{p}{)}\PY{p}{:}
        \PY{n}{bf\PYZus{}records} \PY{o}{=} \PY{n}{bf\PYZus{}results}\PY{o}{.}\PY{n}{to\PYZus{}dict}\PY{p}{(}\PY{l+s+s1}{\PYZsq{}}\PY{l+s+s1}{records}\PY{l+s+s1}{\PYZsq{}}\PY{p}{)}
    \PY{k}{else}\PY{p}{:}
        \PY{n}{bf\PYZus{}records} \PY{o}{=} \PY{n}{bf\PYZus{}results}

    \PY{n}{rd\PYZus{}map} \PY{o}{=} \PY{p}{\PYZob{}}\PY{n}{row}\PY{p}{[}\PY{l+s+s1}{\PYZsq{}}\PY{l+s+s1}{graph\PYZus{}idx}\PY{l+s+s1}{\PYZsq{}}\PY{p}{]}\PY{p}{:} \PY{n}{row} \PY{k}{for} \PY{n}{row} \PY{o+ow}{in} \PY{n}{rd\PYZus{}records}\PY{p}{\PYZcb{}}
    \PY{n}{rows} \PY{o}{=} \PY{p}{[}\PY{p}{]}
    \PY{k}{for} \PY{n}{bf} \PY{o+ow}{in} \PY{n}{bf\PYZus{}records}\PY{p}{:}
        \PY{n}{idx} \PY{o}{=} \PY{n}{bf}\PY{p}{[}\PY{l+s+s1}{\PYZsq{}}\PY{l+s+s1}{graph\PYZus{}idx}\PY{l+s+s1}{\PYZsq{}}\PY{p}{]}
        \PY{k}{if} \PY{n}{idx} \PY{o+ow}{in} \PY{n}{rd\PYZus{}map}\PY{p}{:}
            \PY{n}{rd} \PY{o}{=} \PY{n}{rd\PYZus{}map}\PY{p}{[}\PY{n}{idx}\PY{p}{]}
            \PY{n}{rows}\PY{o}{.}\PY{n}{append}\PY{p}{(}\PY{p}{\PYZob{}}
                \PY{l+s+s1}{\PYZsq{}}\PY{l+s+s1}{graph\PYZus{}idx}\PY{l+s+s1}{\PYZsq{}}\PY{p}{:} \PY{n}{idx}\PY{p}{,}
                \PY{l+s+s1}{\PYZsq{}}\PY{l+s+s1}{bf\PYZus{}clique\PYZus{}size}\PY{l+s+s1}{\PYZsq{}}\PY{p}{:} \PY{n}{bf}\PY{p}{[}\PY{l+s+s1}{\PYZsq{}}\PY{l+s+s1}{clique\PYZus{}size}\PY{l+s+s1}{\PYZsq{}}\PY{p}{]}\PY{p}{,}
                \PY{l+s+s1}{\PYZsq{}}\PY{l+s+s1}{rd\PYZus{}clique\PYZus{}size}\PY{l+s+s1}{\PYZsq{}}\PY{p}{:} \PY{n}{rd}\PY{p}{[}\PY{l+s+s1}{\PYZsq{}}\PY{l+s+s1}{clique\PYZus{}size}\PY{l+s+s1}{\PYZsq{}}\PY{p}{]}\PY{p}{,}
                \PY{l+s+s1}{\PYZsq{}}\PY{l+s+s1}{same\PYZus{}size}\PY{l+s+s1}{\PYZsq{}}\PY{p}{:} \PY{n}{bf}\PY{p}{[}\PY{l+s+s1}{\PYZsq{}}\PY{l+s+s1}{clique\PYZus{}size}\PY{l+s+s1}{\PYZsq{}}\PY{p}{]} \PY{o}{==} \PY{n}{rd}\PY{p}{[}\PY{l+s+s1}{\PYZsq{}}\PY{l+s+s1}{clique\PYZus{}size}\PY{l+s+s1}{\PYZsq{}}\PY{p}{]}\PY{p}{,}
                \PY{l+s+s1}{\PYZsq{}}\PY{l+s+s1}{rd\PYZus{}valid}\PY{l+s+s1}{\PYZsq{}}\PY{p}{:} \PY{n}{rd}\PY{p}{[}\PY{l+s+s1}{\PYZsq{}}\PY{l+s+s1}{valid}\PY{l+s+s1}{\PYZsq{}}\PY{p}{]}\PY{p}{,}
            \PY{p}{\PYZcb{}}\PY{p}{)}
    \PY{k}{return} \PY{n}{pd}\PY{o}{.}\PY{n}{DataFrame}\PY{p}{(}\PY{n}{rows}\PY{p}{)}
\end{Verbatim}
\end{tcolorbox}

    \subsubsection{4.5 GNN + Dinámicas
replicadoras}\label{gnn-dinuxe1micas-replicadoras}

    Tarea 3.1: \texttt{compute\_node\_features()}

En esta tarea se han construido las features de nodo que alimentan a la
GNN. Se han usado grado normalizado, coeficiente de clustering y
\texttt{log(degree\ +\ 1)}.

    \begin{tcolorbox}[breakable, size=fbox, boxrule=1pt, pad at break*=1mm,colback=cellbackground, colframe=cellborder]
\prompt{In}{incolor}{265}{\boxspacing}
\begin{Verbatim}[commandchars=\\\{\}]
\PY{k}{def}\PY{+w}{ }\PY{n+nf}{compute\PYZus{}node\PYZus{}features}\PY{p}{(}\PY{n}{data}\PY{p}{)}\PY{p}{:}
\PY{+w}{    }\PY{l+s+sd}{\PYZdq{}\PYZdq{}\PYZdq{}}
\PY{l+s+sd}{    Compute informative node features for the GNN.}

\PY{l+s+sd}{    Features:}
\PY{l+s+sd}{    \PYZhy{} Normalized degree}
\PY{l+s+sd}{    \PYZhy{} Clustering coefficient}
\PY{l+s+sd}{    \PYZhy{} Log(degree + 1)}

\PY{l+s+sd}{    Args:}
\PY{l+s+sd}{        data: PyTorch Geometric Data object}

\PY{l+s+sd}{    Returns:}
\PY{l+s+sd}{        x: torch.Tensor of shape (num\PYZus{}nodes, num\PYZus{}features)}
\PY{l+s+sd}{    \PYZdq{}\PYZdq{}\PYZdq{}}
    \PY{n}{G} \PY{o}{=} \PY{n}{to\PYZus{}networkx}\PY{p}{(}\PY{n}{data}\PY{p}{,} \PY{n}{to\PYZus{}undirected}\PY{o}{=}\PY{k+kc}{True}\PY{p}{)}
    \PY{n}{degrees} \PY{o}{=} \PY{n+nb}{dict}\PY{p}{(}\PY{n}{G}\PY{o}{.}\PY{n}{degree}\PY{p}{(}\PY{p}{)}\PY{p}{)}
    \PY{n}{clustering} \PY{o}{=} \PY{n}{nx}\PY{o}{.}\PY{n}{clustering}\PY{p}{(}\PY{n}{G}\PY{p}{)}
    \PY{n}{max\PYZus{}degree} \PY{o}{=} \PY{n+nb}{max}\PY{p}{(}\PY{n}{degrees}\PY{o}{.}\PY{n}{values}\PY{p}{(}\PY{p}{)}\PY{p}{)} \PY{k}{if} \PY{n}{degrees} \PY{k}{else} \PY{l+m+mi}{1}

    \PY{n}{features} \PY{o}{=} \PY{p}{[}\PY{p}{]}
    \PY{k}{for} \PY{n}{node} \PY{o+ow}{in} \PY{n+nb}{range}\PY{p}{(}\PY{n}{G}\PY{o}{.}\PY{n}{number\PYZus{}of\PYZus{}nodes}\PY{p}{(}\PY{p}{)}\PY{p}{)}\PY{p}{:}
        \PY{n}{degree} \PY{o}{=} \PY{n}{degrees}\PY{p}{[}\PY{n}{node}\PY{p}{]}
        \PY{n}{features}\PY{o}{.}\PY{n}{append}\PY{p}{(}\PY{p}{[}
            \PY{n}{degree} \PY{o}{/} \PY{n+nb}{max}\PY{p}{(}\PY{l+m+mi}{1}\PY{p}{,} \PY{n}{max\PYZus{}degree}\PY{p}{)}\PY{p}{,}
            \PY{n}{clustering}\PY{p}{[}\PY{n}{node}\PY{p}{]}\PY{p}{,}
            \PY{n}{np}\PY{o}{.}\PY{n}{log1p}\PY{p}{(}\PY{n}{degree}\PY{p}{)}\PY{p}{,}
        \PY{p}{]}\PY{p}{)}

    \PY{k}{return} \PY{n}{torch}\PY{o}{.}\PY{n}{tensor}\PY{p}{(}\PY{n}{features}\PY{p}{,} \PY{n}{dtype}\PY{o}{=}\PY{n}{torch}\PY{o}{.}\PY{n}{float32}\PY{p}{)}
\end{Verbatim}
\end{tcolorbox}

    Tarea 3.2: \texttt{MaxCliqueGNN}

En esta tarea se ha implementado el modelo GNN encargado de producir un
punto inicial del simplex mediante una salida \texttt{softmax} sobre los
nodos.

    \begin{tcolorbox}[breakable, size=fbox, boxrule=1pt, pad at break*=1mm,colback=cellbackground, colframe=cellborder]
\prompt{In}{incolor}{266}{\boxspacing}
\begin{Verbatim}[commandchars=\\\{\}]
\PY{k}{class}\PY{+w}{ }\PY{n+nc}{MaxCliqueGNN}\PY{p}{(}\PY{n}{nn}\PY{o}{.}\PY{n}{Module}\PY{p}{)}\PY{p}{:}
\PY{+w}{    }\PY{l+s+sd}{\PYZdq{}\PYZdq{}\PYZdq{}}
\PY{l+s+sd}{    GNN for predicting initial simplex point for Maximum Clique.}
\PY{l+s+sd}{    \PYZdq{}\PYZdq{}\PYZdq{}}

    \PY{k}{def}\PY{+w}{ }\PY{n+nf+fm}{\PYZus{}\PYZus{}init\PYZus{}\PYZus{}}\PY{p}{(}\PY{n+nb+bp}{self}\PY{p}{,} \PY{n}{input\PYZus{}dim}\PY{p}{,} \PY{n}{hidden\PYZus{}dim}\PY{o}{=}\PY{l+m+mi}{32}\PY{p}{,} \PY{n}{num\PYZus{}layers}\PY{o}{=}\PY{l+m+mi}{2}\PY{p}{)}\PY{p}{:}
        \PY{n+nb}{super}\PY{p}{(}\PY{p}{)}\PY{o}{.}\PY{n+nf+fm}{\PYZus{}\PYZus{}init\PYZus{}\PYZus{}}\PY{p}{(}\PY{p}{)}

        \PY{n+nb+bp}{self}\PY{o}{.}\PY{n}{convs} \PY{o}{=} \PY{n}{nn}\PY{o}{.}\PY{n}{ModuleList}\PY{p}{(}\PY{p}{)}
        \PY{n+nb+bp}{self}\PY{o}{.}\PY{n}{convs}\PY{o}{.}\PY{n}{append}\PY{p}{(}\PY{n}{GCNConv}\PY{p}{(}\PY{n}{input\PYZus{}dim}\PY{p}{,} \PY{n}{hidden\PYZus{}dim}\PY{p}{)}\PY{p}{)}
        \PY{k}{for} \PY{n}{\PYZus{}} \PY{o+ow}{in} \PY{n+nb}{range}\PY{p}{(}\PY{n}{num\PYZus{}layers} \PY{o}{\PYZhy{}} \PY{l+m+mi}{1}\PY{p}{)}\PY{p}{:}
            \PY{n+nb+bp}{self}\PY{o}{.}\PY{n}{convs}\PY{o}{.}\PY{n}{append}\PY{p}{(}\PY{n}{GCNConv}\PY{p}{(}\PY{n}{hidden\PYZus{}dim}\PY{p}{,} \PY{n}{hidden\PYZus{}dim}\PY{p}{)}\PY{p}{)}

        \PY{n+nb+bp}{self}\PY{o}{.}\PY{n}{output} \PY{o}{=} \PY{n}{nn}\PY{o}{.}\PY{n}{Linear}\PY{p}{(}\PY{n}{hidden\PYZus{}dim}\PY{p}{,} \PY{l+m+mi}{1}\PY{p}{)}

    \PY{k}{def}\PY{+w}{ }\PY{n+nf}{forward}\PY{p}{(}\PY{n+nb+bp}{self}\PY{p}{,} \PY{n}{x}\PY{p}{,} \PY{n}{edge\PYZus{}index}\PY{p}{)}\PY{p}{:}
\PY{+w}{        }\PY{l+s+sd}{\PYZdq{}\PYZdq{}\PYZdq{}}
\PY{l+s+sd}{        Forward pass.}
\PY{l+s+sd}{        \PYZdq{}\PYZdq{}\PYZdq{}}
        \PY{k}{for} \PY{n}{conv} \PY{o+ow}{in} \PY{n+nb+bp}{self}\PY{o}{.}\PY{n}{convs}\PY{p}{:}
            \PY{n}{x} \PY{o}{=} \PY{n}{conv}\PY{p}{(}\PY{n}{x}\PY{p}{,} \PY{n}{edge\PYZus{}index}\PY{p}{)}
            \PY{n}{x} \PY{o}{=} \PY{n}{F}\PY{o}{.}\PY{n}{relu}\PY{p}{(}\PY{n}{x}\PY{p}{)}
        \PY{n}{logits} \PY{o}{=} \PY{n+nb+bp}{self}\PY{o}{.}\PY{n}{output}\PY{p}{(}\PY{n}{x}\PY{p}{)}\PY{o}{.}\PY{n}{squeeze}\PY{p}{(}\PY{o}{\PYZhy{}}\PY{l+m+mi}{1}\PY{p}{)}
        \PY{k}{return} \PY{n}{torch}\PY{o}{.}\PY{n}{softmax}\PY{p}{(}\PY{n}{logits}\PY{p}{,} \PY{n}{dim}\PY{o}{=}\PY{l+m+mi}{0}\PY{p}{)}
\end{Verbatim}
\end{tcolorbox}

    Tarea 3.3: RD diferenciable

En esta tarea se ha implementado la versión diferenciable de RD para
integrarla dentro del bucle de entrenamiento.

    \begin{tcolorbox}[breakable, size=fbox, boxrule=1pt, pad at break*=1mm,colback=cellbackground, colframe=cellborder]
\prompt{In}{incolor}{267}{\boxspacing}
\begin{Verbatim}[commandchars=\\\{\}]
\PY{k}{def}\PY{+w}{ }\PY{n+nf}{differentiable\PYZus{}rd\PYZus{}step}\PY{p}{(}\PY{n}{p}\PY{p}{,} \PY{n}{A}\PY{p}{)}\PY{p}{:}
\PY{+w}{    }\PY{l+s+sd}{\PYZdq{}\PYZdq{}\PYZdq{}}
\PY{l+s+sd}{    One step of differentiable Replicator Dynamics.}
\PY{l+s+sd}{    \PYZdq{}\PYZdq{}\PYZdq{}}
    \PY{n}{p\PYZus{}new} \PY{o}{=} \PY{n}{p} \PY{o}{*} \PY{p}{(}\PY{n}{A} \PY{o}{@} \PY{n}{p}\PY{p}{)}
    \PY{n}{p\PYZus{}new} \PY{o}{=} \PY{n}{p\PYZus{}new} \PY{o}{/} \PY{n}{p\PYZus{}new}\PY{o}{.}\PY{n}{sum}\PY{p}{(}\PY{p}{)}\PY{o}{.}\PY{n}{clamp}\PY{p}{(}\PY{n+nb}{min}\PY{o}{=}\PY{l+m+mf}{1e\PYZhy{}12}\PY{p}{)}
    \PY{k}{return} \PY{n}{p\PYZus{}new}


\PY{k}{def}\PY{+w}{ }\PY{n+nf}{differentiable\PYZus{}rd}\PY{p}{(}\PY{n}{p}\PY{p}{,} \PY{n}{A}\PY{p}{,} \PY{n}{num\PYZus{}iterations}\PY{o}{=}\PY{l+m+mi}{3}\PY{p}{)}\PY{p}{:}
\PY{+w}{    }\PY{l+s+sd}{\PYZdq{}\PYZdq{}\PYZdq{}}
\PY{l+s+sd}{    Multiple steps of differentiable RD.}
\PY{l+s+sd}{    \PYZdq{}\PYZdq{}\PYZdq{}}
    \PY{n}{A\PYZus{}reg} \PY{o}{=} \PY{n}{A} \PY{o}{+} \PY{l+m+mf}{0.5} \PY{o}{*} \PY{n}{torch}\PY{o}{.}\PY{n}{eye}\PY{p}{(}\PY{n}{A}\PY{o}{.}\PY{n}{shape}\PY{p}{[}\PY{l+m+mi}{0}\PY{p}{]}\PY{p}{,} \PY{n}{device}\PY{o}{=}\PY{n}{A}\PY{o}{.}\PY{n}{device}\PY{p}{)}

    \PY{k}{for} \PY{n}{\PYZus{}} \PY{o+ow}{in} \PY{n+nb}{range}\PY{p}{(}\PY{n}{num\PYZus{}iterations}\PY{p}{)}\PY{p}{:}
        \PY{n}{p} \PY{o}{=} \PY{n}{differentiable\PYZus{}rd\PYZus{}step}\PY{p}{(}\PY{n}{p}\PY{p}{,} \PY{n}{A\PYZus{}reg}\PY{p}{)}

    \PY{k}{return} \PY{n}{p}
\end{Verbatim}
\end{tcolorbox}

    Tarea 3.4: \texttt{maxclique\_loss()}

En esta tarea se ha implementado la pérdida de Motzkin-Straus en signo
negativo para poder minimizarla.

    \begin{tcolorbox}[breakable, size=fbox, boxrule=1pt, pad at break*=1mm,colback=cellbackground, colframe=cellborder]
\prompt{In}{incolor}{268}{\boxspacing}
\begin{Verbatim}[commandchars=\\\{\}]
\PY{k}{def}\PY{+w}{ }\PY{n+nf}{maxclique\PYZus{}loss}\PY{p}{(}\PY{n}{p}\PY{p}{,} \PY{n}{A}\PY{p}{)}\PY{p}{:}
\PY{+w}{    }\PY{l+s+sd}{\PYZdq{}\PYZdq{}\PYZdq{}}
\PY{l+s+sd}{    Motzkin\PYZhy{}Straus loss (negative because we minimize).}

\PY{l+s+sd}{    Args:}
\PY{l+s+sd}{        p: torch.Tensor (n,) \PYZhy{} simplex point}
\PY{l+s+sd}{        A: torch.Tensor (n, n) \PYZhy{} adjacency matrix}

\PY{l+s+sd}{    Returns:}
\PY{l+s+sd}{        loss: scalar tensor}
\PY{l+s+sd}{    \PYZdq{}\PYZdq{}\PYZdq{}}
    \PY{k}{return} \PY{o}{\PYZhy{}}\PY{p}{(}\PY{n}{p} \PY{o}{@} \PY{n}{A} \PY{o}{@} \PY{n}{p}\PY{p}{)}
\end{Verbatim}
\end{tcolorbox}

    Tarea 3.5: entrenamiento

En esta tarea se ha implementado el bucle de entrenamiento de la GNN
sobre listas de grafos, refinando con RD un pequeño número de
iteraciones durante el aprendizaje.

    \begin{tcolorbox}[breakable, size=fbox, boxrule=1pt, pad at break*=1mm,colback=cellbackground, colframe=cellborder]
\prompt{In}{incolor}{269}{\boxspacing}
\begin{Verbatim}[commandchars=\\\{\}]
\PY{k}{def}\PY{+w}{ }\PY{n+nf}{train\PYZus{}gnn}\PY{p}{(}\PY{n}{model}\PY{p}{,} \PY{n}{dataset}\PY{p}{,} \PY{n}{num\PYZus{}epochs}\PY{o}{=}\PY{l+m+mi}{12}\PY{p}{,} \PY{n}{lr}\PY{o}{=}\PY{l+m+mf}{0.01}\PY{p}{,} \PY{n}{rd\PYZus{}train\PYZus{}iters}\PY{o}{=}\PY{l+m+mi}{3}\PY{p}{)}\PY{p}{:}
\PY{+w}{    }\PY{l+s+sd}{\PYZdq{}\PYZdq{}\PYZdq{}}
\PY{l+s+sd}{    Train the GNN model.}
\PY{l+s+sd}{    \PYZdq{}\PYZdq{}\PYZdq{}}
    \PY{n}{model} \PY{o}{=} \PY{n}{model}\PY{o}{.}\PY{n}{to}\PY{p}{(}\PY{n}{DEVICE}\PY{p}{)}
    \PY{n}{optimizer} \PY{o}{=} \PY{n}{torch}\PY{o}{.}\PY{n}{optim}\PY{o}{.}\PY{n}{Adam}\PY{p}{(}\PY{n}{model}\PY{o}{.}\PY{n}{parameters}\PY{p}{(}\PY{p}{)}\PY{p}{,} \PY{n}{lr}\PY{o}{=}\PY{n}{lr}\PY{p}{)}
    \PY{n}{losses} \PY{o}{=} \PY{p}{[}\PY{p}{]}

    \PY{k}{for} \PY{n}{epoch} \PY{o+ow}{in} \PY{n+nb}{range}\PY{p}{(}\PY{n}{num\PYZus{}epochs}\PY{p}{)}\PY{p}{:}
        \PY{n}{model}\PY{o}{.}\PY{n}{train}\PY{p}{(}\PY{p}{)}
        \PY{n}{epoch\PYZus{}loss} \PY{o}{=} \PY{l+m+mf}{0.0}

        \PY{k}{for} \PY{n}{data} \PY{o+ow}{in} \PY{n}{dataset}\PY{p}{:}
            \PY{n}{optimizer}\PY{o}{.}\PY{n}{zero\PYZus{}grad}\PY{p}{(}\PY{p}{)}
            \PY{n}{x} \PY{o}{=} \PY{n}{compute\PYZus{}node\PYZus{}features}\PY{p}{(}\PY{n}{data}\PY{p}{)}\PY{o}{.}\PY{n}{to}\PY{p}{(}\PY{n}{DEVICE}\PY{p}{)}
            \PY{n}{edge\PYZus{}index} \PY{o}{=} \PY{n}{data}\PY{o}{.}\PY{n}{edge\PYZus{}index}\PY{o}{.}\PY{n}{to}\PY{p}{(}\PY{n}{DEVICE}\PY{p}{)}
            \PY{n}{A} \PY{o}{=} \PY{n}{torch}\PY{o}{.}\PY{n}{tensor}\PY{p}{(}\PY{n}{pyg\PYZus{}to\PYZus{}adjacency}\PY{p}{(}\PY{n}{data}\PY{p}{)}\PY{p}{,} \PY{n}{dtype}\PY{o}{=}\PY{n}{torch}\PY{o}{.}\PY{n}{float32}\PY{p}{,} \PY{n}{device}\PY{o}{=}\PY{n}{DEVICE}\PY{p}{)}
            \PY{n}{p} \PY{o}{=} \PY{n}{model}\PY{p}{(}\PY{n}{x}\PY{p}{,} \PY{n}{edge\PYZus{}index}\PY{p}{)}
            \PY{n}{p\PYZus{}refined} \PY{o}{=} \PY{n}{differentiable\PYZus{}rd}\PY{p}{(}\PY{n}{p}\PY{p}{,} \PY{n}{A}\PY{p}{,} \PY{n}{num\PYZus{}iterations}\PY{o}{=}\PY{n}{rd\PYZus{}train\PYZus{}iters}\PY{p}{)} \PY{k}{if} \PY{n}{rd\PYZus{}train\PYZus{}iters} \PY{o}{\PYZgt{}} \PY{l+m+mi}{0} \PY{k}{else} \PY{n}{p}
            \PY{n}{loss} \PY{o}{=} \PY{n}{maxclique\PYZus{}loss}\PY{p}{(}\PY{n}{p\PYZus{}refined}\PY{p}{,} \PY{n}{A}\PY{p}{)}
            \PY{n}{loss}\PY{o}{.}\PY{n}{backward}\PY{p}{(}\PY{p}{)}
            \PY{n}{optimizer}\PY{o}{.}\PY{n}{step}\PY{p}{(}\PY{p}{)}
            \PY{n}{epoch\PYZus{}loss} \PY{o}{+}\PY{o}{=} \PY{n+nb}{float}\PY{p}{(}\PY{n}{loss}\PY{o}{.}\PY{n}{item}\PY{p}{(}\PY{p}{)}\PY{p}{)}

        \PY{n}{epoch\PYZus{}loss} \PY{o}{/}\PY{o}{=} \PY{n+nb}{max}\PY{p}{(}\PY{l+m+mi}{1}\PY{p}{,} \PY{n+nb}{len}\PY{p}{(}\PY{n}{dataset}\PY{p}{)}\PY{p}{)}
        \PY{n}{losses}\PY{o}{.}\PY{n}{append}\PY{p}{(}\PY{n}{epoch\PYZus{}loss}\PY{p}{)}

        \PY{k}{if} \PY{p}{(}\PY{n}{epoch} \PY{o}{+} \PY{l+m+mi}{1}\PY{p}{)} \PY{o}{\PYZpc{}} \PY{n+nb}{max}\PY{p}{(}\PY{l+m+mi}{1}\PY{p}{,} \PY{n}{num\PYZus{}epochs} \PY{o}{/}\PY{o}{/} \PY{l+m+mi}{4}\PY{p}{)} \PY{o}{==} \PY{l+m+mi}{0}\PY{p}{:}
            \PY{n+nb}{print}\PY{p}{(}\PY{l+s+sa}{f}\PY{l+s+s2}{\PYZdq{}}\PY{l+s+s2}{Época }\PY{l+s+si}{\PYZob{}}\PY{n}{epoch}\PY{+w}{ }\PY{o}{+}\PY{+w}{ }\PY{l+m+mi}{1}\PY{l+s+si}{\PYZcb{}}\PY{l+s+s2}{/}\PY{l+s+si}{\PYZob{}}\PY{n}{num\PYZus{}epochs}\PY{l+s+si}{\PYZcb{}}\PY{l+s+s2}{ \PYZhy{} loss media: }\PY{l+s+si}{\PYZob{}}\PY{n}{epoch\PYZus{}loss}\PY{l+s+si}{:}\PY{l+s+s2}{.4f}\PY{l+s+si}{\PYZcb{}}\PY{l+s+s2}{\PYZdq{}}\PY{p}{)}

    \PY{k}{return} \PY{n}{model}\PY{p}{,} \PY{n}{losses}
\end{Verbatim}
\end{tcolorbox}

    Tarea 3.6: evaluación

En esta tarea se ha implementado la evaluación del modelo híbrido y una
pequeña ablación de iteraciones RD en entrenamiento e inferencia.

    \begin{tcolorbox}[breakable, size=fbox, boxrule=1pt, pad at break*=1mm,colback=cellbackground, colframe=cellborder]
\prompt{In}{incolor}{270}{\boxspacing}
\begin{Verbatim}[commandchars=\\\{\}]
\PY{k}{def}\PY{+w}{ }\PY{n+nf}{evaluate\PYZus{}gnn}\PY{p}{(}\PY{n}{model}\PY{p}{,} \PY{n}{dataset}\PY{p}{,} \PY{n}{rd\PYZus{}test\PYZus{}iters}\PY{o}{=}\PY{l+m+mi}{100}\PY{p}{)}\PY{p}{:}
\PY{+w}{    }\PY{l+s+sd}{\PYZdq{}\PYZdq{}\PYZdq{}}
\PY{l+s+sd}{    Evaluate the trained model.}
\PY{l+s+sd}{    \PYZdq{}\PYZdq{}\PYZdq{}}
    \PY{n}{model} \PY{o}{=} \PY{n}{model}\PY{o}{.}\PY{n}{to}\PY{p}{(}\PY{n}{DEVICE}\PY{p}{)}
    \PY{n}{model}\PY{o}{.}\PY{n}{eval}\PY{p}{(}\PY{p}{)}
    \PY{n}{results} \PY{o}{=} \PY{p}{[}\PY{p}{]}

    \PY{k}{with} \PY{n}{torch}\PY{o}{.}\PY{n}{no\PYZus{}grad}\PY{p}{(}\PY{p}{)}\PY{p}{:}
        \PY{k}{for} \PY{n}{i}\PY{p}{,} \PY{n}{data} \PY{o+ow}{in} \PY{n+nb}{enumerate}\PY{p}{(}\PY{n}{dataset}\PY{p}{)}\PY{p}{:}
            \PY{n}{x} \PY{o}{=} \PY{n}{compute\PYZus{}node\PYZus{}features}\PY{p}{(}\PY{n}{data}\PY{p}{)}\PY{o}{.}\PY{n}{to}\PY{p}{(}\PY{n}{DEVICE}\PY{p}{)}
            \PY{n}{edge\PYZus{}index} \PY{o}{=} \PY{n}{data}\PY{o}{.}\PY{n}{edge\PYZus{}index}\PY{o}{.}\PY{n}{to}\PY{p}{(}\PY{n}{DEVICE}\PY{p}{)}
            \PY{n}{A\PYZus{}np} \PY{o}{=} \PY{n}{pyg\PYZus{}to\PYZus{}adjacency}\PY{p}{(}\PY{n}{data}\PY{p}{)}
            \PY{n}{A} \PY{o}{=} \PY{n}{torch}\PY{o}{.}\PY{n}{tensor}\PY{p}{(}\PY{n}{A\PYZus{}np}\PY{p}{,} \PY{n}{dtype}\PY{o}{=}\PY{n}{torch}\PY{o}{.}\PY{n}{float32}\PY{p}{,} \PY{n}{device}\PY{o}{=}\PY{n}{DEVICE}\PY{p}{)}

            \PY{n}{start\PYZus{}time} \PY{o}{=} \PY{n}{time}\PY{o}{.}\PY{n}{time}\PY{p}{(}\PY{p}{)}
            \PY{n}{p} \PY{o}{=} \PY{n}{model}\PY{p}{(}\PY{n}{x}\PY{p}{,} \PY{n}{edge\PYZus{}index}\PY{p}{)}
            \PY{n}{p} \PY{o}{=} \PY{n}{differentiable\PYZus{}rd}\PY{p}{(}\PY{n}{p}\PY{p}{,} \PY{n}{A}\PY{p}{,} \PY{n}{num\PYZus{}iterations}\PY{o}{=}\PY{n}{rd\PYZus{}test\PYZus{}iters}\PY{p}{)}
            \PY{n}{clique} \PY{o}{=} \PY{n}{decode\PYZus{}clique}\PY{p}{(}\PY{n}{p}\PY{o}{.}\PY{n}{detach}\PY{p}{(}\PY{p}{)}\PY{o}{.}\PY{n}{cpu}\PY{p}{(}\PY{p}{)}\PY{o}{.}\PY{n}{numpy}\PY{p}{(}\PY{p}{)}\PY{p}{,} \PY{n}{A\PYZus{}np}\PY{p}{)}
            \PY{n}{elapsed} \PY{o}{=} \PY{n}{time}\PY{o}{.}\PY{n}{time}\PY{p}{(}\PY{p}{)} \PY{o}{\PYZhy{}} \PY{n}{start\PYZus{}time}

            \PY{n}{results}\PY{o}{.}\PY{n}{append}\PY{p}{(}\PY{p}{\PYZob{}}
                \PY{l+s+s1}{\PYZsq{}}\PY{l+s+s1}{graph\PYZus{}idx}\PY{l+s+s1}{\PYZsq{}}\PY{p}{:} \PY{n}{i}\PY{p}{,}
                \PY{l+s+s1}{\PYZsq{}}\PY{l+s+s1}{num\PYZus{}nodes}\PY{l+s+s1}{\PYZsq{}}\PY{p}{:} \PY{n+nb}{int}\PY{p}{(}\PY{n}{data}\PY{o}{.}\PY{n}{num\PYZus{}nodes}\PY{p}{)}\PY{p}{,}
                \PY{l+s+s1}{\PYZsq{}}\PY{l+s+s1}{clique\PYZus{}size}\PY{l+s+s1}{\PYZsq{}}\PY{p}{:} \PY{n+nb}{int}\PY{p}{(}\PY{n+nb}{len}\PY{p}{(}\PY{n}{clique}\PY{p}{)}\PY{p}{)}\PY{p}{,}
                \PY{l+s+s1}{\PYZsq{}}\PY{l+s+s1}{clique}\PY{l+s+s1}{\PYZsq{}}\PY{p}{:} \PY{n}{clique}\PY{p}{,}
                \PY{l+s+s1}{\PYZsq{}}\PY{l+s+s1}{valid}\PY{l+s+s1}{\PYZsq{}}\PY{p}{:} \PY{n+nb}{bool}\PY{p}{(}\PY{n}{is\PYZus{}clique}\PY{p}{(}\PY{n}{clique}\PY{p}{,} \PY{n}{A\PYZus{}np}\PY{p}{)}\PY{p}{)}\PY{p}{,}
                \PY{l+s+s1}{\PYZsq{}}\PY{l+s+s1}{time}\PY{l+s+s1}{\PYZsq{}}\PY{p}{:} \PY{n+nb}{float}\PY{p}{(}\PY{n}{elapsed}\PY{p}{)}\PY{p}{,}
            \PY{p}{\PYZcb{}}\PY{p}{)}

    \PY{k}{return} \PY{n}{results}


\PY{k}{def}\PY{+w}{ }\PY{n+nf}{run\PYZus{}gnn\PYZus{}rd\PYZus{}ablation}\PY{p}{(}\PY{n}{train\PYZus{}dataset}\PY{p}{,} \PY{n}{test\PYZus{}dataset}\PY{p}{,} \PY{n}{train\PYZus{}iters\PYZus{}list}\PY{o}{=}\PY{p}{(}\PY{l+m+mi}{0}\PY{p}{,} \PY{l+m+mi}{1}\PY{p}{,} \PY{l+m+mi}{3}\PY{p}{,} \PY{l+m+mi}{5}\PY{p}{)}\PY{p}{,} \PY{n}{test\PYZus{}iters\PYZus{}list}\PY{o}{=}\PY{p}{(}\PY{l+m+mi}{10}\PY{p}{,} \PY{l+m+mi}{50}\PY{p}{,} \PY{l+m+mi}{100}\PY{p}{)}\PY{p}{,} \PY{n}{num\PYZus{}epochs}\PY{o}{=}\PY{l+m+mi}{6}\PY{p}{,} \PY{n}{lr}\PY{o}{=}\PY{l+m+mf}{0.01}\PY{p}{,} \PY{n}{hidden\PYZus{}dim}\PY{o}{=}\PY{l+m+mi}{32}\PY{p}{)}\PY{p}{:}
    \PY{n}{rows} \PY{o}{=} \PY{p}{[}\PY{p}{]}
    \PY{k}{for} \PY{n}{train\PYZus{}iters} \PY{o+ow}{in} \PY{n}{train\PYZus{}iters\PYZus{}list}\PY{p}{:}
        \PY{n}{model} \PY{o}{=} \PY{n}{MaxCliqueGNN}\PY{p}{(}\PY{n}{input\PYZus{}dim}\PY{o}{=}\PY{l+m+mi}{3}\PY{p}{,} \PY{n}{hidden\PYZus{}dim}\PY{o}{=}\PY{n}{hidden\PYZus{}dim}\PY{p}{,} \PY{n}{num\PYZus{}layers}\PY{o}{=}\PY{l+m+mi}{2}\PY{p}{)}
        \PY{n}{model}\PY{p}{,} \PY{n}{\PYZus{}} \PY{o}{=} \PY{n}{train\PYZus{}gnn}\PY{p}{(}\PY{n}{model}\PY{p}{,} \PY{n}{train\PYZus{}dataset}\PY{p}{,} \PY{n}{num\PYZus{}epochs}\PY{o}{=}\PY{n}{num\PYZus{}epochs}\PY{p}{,} \PY{n}{lr}\PY{o}{=}\PY{n}{lr}\PY{p}{,} \PY{n}{rd\PYZus{}train\PYZus{}iters}\PY{o}{=}\PY{n}{train\PYZus{}iters}\PY{p}{)}
        \PY{k}{for} \PY{n}{test\PYZus{}iters} \PY{o+ow}{in} \PY{n}{test\PYZus{}iters\PYZus{}list}\PY{p}{:}
            \PY{n}{eval\PYZus{}results} \PY{o}{=} \PY{n}{evaluate\PYZus{}gnn}\PY{p}{(}\PY{n}{model}\PY{p}{,} \PY{n}{test\PYZus{}dataset}\PY{p}{,} \PY{n}{rd\PYZus{}test\PYZus{}iters}\PY{o}{=}\PY{n}{test\PYZus{}iters}\PY{p}{)}
            \PY{n}{summary} \PY{o}{=} \PY{n}{make\PYZus{}summary\PYZus{}dict}\PY{p}{(}\PY{n}{eval\PYZus{}results}\PY{p}{)}
            \PY{n}{rows}\PY{o}{.}\PY{n}{append}\PY{p}{(}\PY{p}{\PYZob{}}
                \PY{l+s+s1}{\PYZsq{}}\PY{l+s+s1}{train\PYZus{}rd\PYZus{}iters}\PY{l+s+s1}{\PYZsq{}}\PY{p}{:} \PY{n}{train\PYZus{}iters}\PY{p}{,}
                \PY{l+s+s1}{\PYZsq{}}\PY{l+s+s1}{test\PYZus{}rd\PYZus{}iters}\PY{l+s+s1}{\PYZsq{}}\PY{p}{:} \PY{n}{test\PYZus{}iters}\PY{p}{,}
                \PY{l+s+s1}{\PYZsq{}}\PY{l+s+s1}{mean\PYZus{}clique\PYZus{}size}\PY{l+s+s1}{\PYZsq{}}\PY{p}{:} \PY{n}{summary}\PY{p}{[}\PY{l+s+s1}{\PYZsq{}}\PY{l+s+s1}{mean\PYZus{}clique\PYZus{}size}\PY{l+s+s1}{\PYZsq{}}\PY{p}{]}\PY{p}{,}
                \PY{l+s+s1}{\PYZsq{}}\PY{l+s+s1}{mean\PYZus{}time}\PY{l+s+s1}{\PYZsq{}}\PY{p}{:} \PY{n}{summary}\PY{p}{[}\PY{l+s+s1}{\PYZsq{}}\PY{l+s+s1}{mean\PYZus{}time}\PY{l+s+s1}{\PYZsq{}}\PY{p}{]}\PY{p}{,}
                \PY{l+s+s1}{\PYZsq{}}\PY{l+s+s1}{valid\PYZus{}rate}\PY{l+s+s1}{\PYZsq{}}\PY{p}{:} \PY{n}{summary}\PY{p}{[}\PY{l+s+s1}{\PYZsq{}}\PY{l+s+s1}{valid\PYZus{}rate}\PY{l+s+s1}{\PYZsq{}}\PY{p}{]}\PY{p}{,}
            \PY{p}{\PYZcb{}}\PY{p}{)}
    \PY{k}{return} \PY{n}{pd}\PY{o}{.}\PY{n}{DataFrame}\PY{p}{(}\PY{n}{rows}\PY{p}{)}
\end{Verbatim}
\end{tcolorbox}

    \subsubsection{4.6 Presentación de resultados y
comparación}\label{presentaciuxf3n-de-resultados-y-comparaciuxf3n}

En esta última parte de la implementación se han reunido las utilidades
para representar pérdidas y visualizar la ablación.

    \begin{tcolorbox}[breakable, size=fbox, boxrule=1pt, pad at break*=1mm,colback=cellbackground, colframe=cellborder]
\prompt{In}{incolor}{271}{\boxspacing}
\begin{Verbatim}[commandchars=\\\{\}]
\PY{k}{def}\PY{+w}{ }\PY{n+nf}{plot\PYZus{}training\PYZus{}losses}\PY{p}{(}\PY{n}{losses}\PY{p}{,} \PY{n}{title}\PY{p}{)}\PY{p}{:}
    \PY{n}{fig}\PY{p}{,} \PY{n}{ax} \PY{o}{=} \PY{n}{plt}\PY{o}{.}\PY{n}{subplots}\PY{p}{(}\PY{n}{figsize}\PY{o}{=}\PY{p}{(}\PY{l+m+mi}{6}\PY{p}{,} \PY{l+m+mi}{4}\PY{p}{)}\PY{p}{)}
    \PY{n}{ax}\PY{o}{.}\PY{n}{plot}\PY{p}{(}\PY{n}{losses}\PY{p}{,} \PY{n}{marker}\PY{o}{=}\PY{l+s+s1}{\PYZsq{}}\PY{l+s+s1}{o}\PY{l+s+s1}{\PYZsq{}}\PY{p}{,} \PY{n}{color}\PY{o}{=}\PY{l+s+s1}{\PYZsq{}}\PY{l+s+s1}{\PYZsh{}d62728}\PY{l+s+s1}{\PYZsq{}}\PY{p}{)}
    \PY{n}{ax}\PY{o}{.}\PY{n}{set\PYZus{}title}\PY{p}{(}\PY{n}{title}\PY{p}{)}
    \PY{n}{ax}\PY{o}{.}\PY{n}{set\PYZus{}xlabel}\PY{p}{(}\PY{l+s+s1}{\PYZsq{}}\PY{l+s+s1}{Época}\PY{l+s+s1}{\PYZsq{}}\PY{p}{)}
    \PY{n}{ax}\PY{o}{.}\PY{n}{set\PYZus{}ylabel}\PY{p}{(}\PY{l+s+s1}{\PYZsq{}}\PY{l+s+s1}{Loss media}\PY{l+s+s1}{\PYZsq{}}\PY{p}{)}
    \PY{n}{plt}\PY{o}{.}\PY{n}{tight\PYZus{}layout}\PY{p}{(}\PY{p}{)}
    \PY{n}{plt}\PY{o}{.}\PY{n}{show}\PY{p}{(}\PY{p}{)}


\PY{k}{def}\PY{+w}{ }\PY{n+nf}{plot\PYZus{}ablation\PYZus{}results}\PY{p}{(}\PY{n}{df}\PY{p}{,} \PY{n}{title}\PY{p}{)}\PY{p}{:}
    \PY{n}{pivot} \PY{o}{=} \PY{n}{df}\PY{o}{.}\PY{n}{pivot}\PY{p}{(}\PY{n}{index}\PY{o}{=}\PY{l+s+s1}{\PYZsq{}}\PY{l+s+s1}{train\PYZus{}rd\PYZus{}iters}\PY{l+s+s1}{\PYZsq{}}\PY{p}{,} \PY{n}{columns}\PY{o}{=}\PY{l+s+s1}{\PYZsq{}}\PY{l+s+s1}{test\PYZus{}rd\PYZus{}iters}\PY{l+s+s1}{\PYZsq{}}\PY{p}{,} \PY{n}{values}\PY{o}{=}\PY{l+s+s1}{\PYZsq{}}\PY{l+s+s1}{mean\PYZus{}clique\PYZus{}size}\PY{l+s+s1}{\PYZsq{}}\PY{p}{)}
    \PY{n}{fig}\PY{p}{,} \PY{n}{ax} \PY{o}{=} \PY{n}{plt}\PY{o}{.}\PY{n}{subplots}\PY{p}{(}\PY{n}{figsize}\PY{o}{=}\PY{p}{(}\PY{l+m+mi}{6}\PY{p}{,} \PY{l+m+mi}{4}\PY{p}{)}\PY{p}{)}
    \PY{n}{im} \PY{o}{=} \PY{n}{ax}\PY{o}{.}\PY{n}{imshow}\PY{p}{(}\PY{n}{pivot}\PY{o}{.}\PY{n}{values}\PY{p}{,} \PY{n}{cmap}\PY{o}{=}\PY{l+s+s1}{\PYZsq{}}\PY{l+s+s1}{YlGnBu}\PY{l+s+s1}{\PYZsq{}}\PY{p}{)}
    \PY{n}{ax}\PY{o}{.}\PY{n}{set\PYZus{}xticks}\PY{p}{(}\PY{n}{np}\PY{o}{.}\PY{n}{arange}\PY{p}{(}\PY{n+nb}{len}\PY{p}{(}\PY{n}{pivot}\PY{o}{.}\PY{n}{columns}\PY{p}{)}\PY{p}{)}\PY{p}{)}
    \PY{n}{ax}\PY{o}{.}\PY{n}{set\PYZus{}xticklabels}\PY{p}{(}\PY{n}{pivot}\PY{o}{.}\PY{n}{columns}\PY{p}{)}
    \PY{n}{ax}\PY{o}{.}\PY{n}{set\PYZus{}yticks}\PY{p}{(}\PY{n}{np}\PY{o}{.}\PY{n}{arange}\PY{p}{(}\PY{n+nb}{len}\PY{p}{(}\PY{n}{pivot}\PY{o}{.}\PY{n}{index}\PY{p}{)}\PY{p}{)}\PY{p}{)}
    \PY{n}{ax}\PY{o}{.}\PY{n}{set\PYZus{}yticklabels}\PY{p}{(}\PY{n}{pivot}\PY{o}{.}\PY{n}{index}\PY{p}{)}
    \PY{n}{ax}\PY{o}{.}\PY{n}{set\PYZus{}xlabel}\PY{p}{(}\PY{l+s+s1}{\PYZsq{}}\PY{l+s+s1}{Iteraciones RD en test}\PY{l+s+s1}{\PYZsq{}}\PY{p}{)}
    \PY{n}{ax}\PY{o}{.}\PY{n}{set\PYZus{}ylabel}\PY{p}{(}\PY{l+s+s1}{\PYZsq{}}\PY{l+s+s1}{Iteraciones RD en train}\PY{l+s+s1}{\PYZsq{}}\PY{p}{)}
    \PY{n}{ax}\PY{o}{.}\PY{n}{set\PYZus{}title}\PY{p}{(}\PY{n}{title}\PY{p}{)}
    \PY{k}{for} \PY{n}{i} \PY{o+ow}{in} \PY{n+nb}{range}\PY{p}{(}\PY{n}{pivot}\PY{o}{.}\PY{n}{shape}\PY{p}{[}\PY{l+m+mi}{0}\PY{p}{]}\PY{p}{)}\PY{p}{:}
        \PY{k}{for} \PY{n}{j} \PY{o+ow}{in} \PY{n+nb}{range}\PY{p}{(}\PY{n}{pivot}\PY{o}{.}\PY{n}{shape}\PY{p}{[}\PY{l+m+mi}{1}\PY{p}{]}\PY{p}{)}\PY{p}{:}
            \PY{n}{ax}\PY{o}{.}\PY{n}{text}\PY{p}{(}\PY{n}{j}\PY{p}{,} \PY{n}{i}\PY{p}{,} \PY{l+s+sa}{f}\PY{l+s+s2}{\PYZdq{}}\PY{l+s+si}{\PYZob{}}\PY{n}{pivot}\PY{o}{.}\PY{n}{values}\PY{p}{[}\PY{n}{i}\PY{p}{,}\PY{+w}{ }\PY{n}{j}\PY{p}{]}\PY{l+s+si}{:}\PY{l+s+s2}{.2f}\PY{l+s+si}{\PYZcb{}}\PY{l+s+s2}{\PYZdq{}}\PY{p}{,} \PY{n}{ha}\PY{o}{=}\PY{l+s+s1}{\PYZsq{}}\PY{l+s+s1}{center}\PY{l+s+s1}{\PYZsq{}}\PY{p}{,} \PY{n}{va}\PY{o}{=}\PY{l+s+s1}{\PYZsq{}}\PY{l+s+s1}{center}\PY{l+s+s1}{\PYZsq{}}\PY{p}{,} \PY{n}{color}\PY{o}{=}\PY{l+s+s1}{\PYZsq{}}\PY{l+s+s1}{black}\PY{l+s+s1}{\PYZsq{}}\PY{p}{)}
    \PY{n}{fig}\PY{o}{.}\PY{n}{colorbar}\PY{p}{(}\PY{n}{im}\PY{p}{,} \PY{n}{ax}\PY{o}{=}\PY{n}{ax}\PY{p}{,} \PY{n}{label}\PY{o}{=}\PY{l+s+s1}{\PYZsq{}}\PY{l+s+s1}{Clique media}\PY{l+s+s1}{\PYZsq{}}\PY{p}{)}
    \PY{n}{plt}\PY{o}{.}\PY{n}{tight\PYZus{}layout}\PY{p}{(}\PY{p}{)}
    \PY{n}{plt}\PY{o}{.}\PY{n}{show}\PY{p}{(}\PY{p}{)}
\end{Verbatim}
\end{tcolorbox}

    \subsection{5. Ejercicio 1: Fuerza bruta para clique máximo (Parte
0)}\label{ejercicio-1-fuerza-bruta-para-clique-muxe1ximo-parte-0}

    \subsubsection{\texorpdfstring{5.1 Implementar \texttt{is\_clique()} y
\texttt{brute\_force\_max\_clique()}}{5.1 Implementar is\_clique() y brute\_force\_max\_clique()}}\label{implementar-is_clique-y-brute_force_max_clique}

La implementación de estas dos funciones ya se completó en la
\textbf{tarea 0.1} y en la \textbf{tarea 0.2} de la sección de
implementación. En este ejercicio no se vuelve a copiar ese código, sino
que se reutiliza directamente para comprobar su comportamiento sobre
ejemplos sintéticos y sobre los primeros grafos del dataset.

    \subsubsection{5.2 Prueba en un grafo sintético
pequeño}\label{prueba-en-un-grafo-sintuxe9tico-pequeuxf1o}

    \begin{tcolorbox}[breakable, size=fbox, boxrule=1pt, pad at break*=1mm,colback=cellbackground, colframe=cellborder]
\prompt{In}{incolor}{272}{\boxspacing}
\begin{Verbatim}[commandchars=\\\{\}]
\PY{n}{brute\PYZus{}force\PYZus{}test} \PY{o}{=} \PY{n}{test\PYZus{}brute\PYZus{}force}\PY{p}{(}\PY{p}{)}
\PY{n}{display}\PY{p}{(}\PY{n}{brute\PYZus{}force\PYZus{}test}\PY{p}{)}
\end{Verbatim}
\end{tcolorbox}

    
    \begin{Verbatim}[commandchars=\\\{\}]
  Maximum clique:  Clique size:       Time:  is valid clique:
0       [0, 1, 2]             3  2.3603e-05              True
    \end{Verbatim}

    
    \subsubsection{5.3 Aplicación a los primeros 20 grafos de IMDB y
contraste adicional con
COLLAB}\label{aplicaciuxf3n-a-los-primeros-20-grafos-de-imdb-y-contraste-adicional-con-collab}

    \begin{tcolorbox}[breakable, size=fbox, boxrule=1pt, pad at break*=1mm,colback=cellbackground, colframe=cellborder]
\prompt{In}{incolor}{273}{\boxspacing}
\begin{Verbatim}[commandchars=\\\{\}]
\PY{n}{bf\PYZus{}imdb\PYZus{}results} \PY{o}{=} \PY{n}{evaluate\PYZus{}brute\PYZus{}force\PYZus{}on\PYZus{}dataset}\PY{p}{(}\PY{n}{imdb\PYZus{}dataset}\PY{p}{,} \PY{n}{max\PYZus{}graphs}\PY{o}{=}\PY{l+m+mi}{20}\PY{p}{,} \PY{n}{max\PYZus{}nodes}\PY{o}{=}\PY{l+m+mi}{25}\PY{p}{)}

\PY{n}{display}\PY{p}{(}\PY{n}{bf\PYZus{}imdb\PYZus{}results}\PY{p}{)}
\end{Verbatim}
\end{tcolorbox}

    \begin{Verbatim}[commandchars=\\\{\}]
Graph 1: Skipping (too large: 32 nodes)
Graph 3: Skipping (too large: 35 nodes)
Graph 5: Skipping (too large: 63 nodes)
Graph 14: Skipping (too large: 72 nodes)
Graph 19: Skipping (too large: 46 nodes)
    \end{Verbatim}

    
    \begin{Verbatim}[commandchars=\\\{\}]
    graph\_idx  num\_nodes  clique\_size  \textbackslash{}
0           0         20            9   
1           2         21           10   
2           4         14            9   
3           6         12           12   
4           7         15            9   
5           8         18           18   
6           9         12           12   
7          10         12            8   
8          11         20           20   
9          12         13           10   
10         13         20           12   
11         15         20           20   
12         16         18           13   
13         17         12            8   
14         18         19            6   

                                               clique  valid      time  
0                    [1, 2, 6, 8, 12, 14, 17, 18, 19]   True  0.335539  
1                  [0, 2, 3, 5, 6, 7, 11, 12, 19, 20]   True  0.552813  
2                      [0, 2, 5, 6, 8, 9, 11, 12, 13]   True  0.002884  
3              [0, 1, 2, 3, 4, 5, 6, 7, 8, 9, 10, 11]   True  0.000011  
4                        [0, 1, 4, 5, 6, 7, 8, 9, 14]   True  0.003583  
5   [0, 1, 2, 3, 4, 5, 6, 7, 8, 9, 10, 11, 12, 13,{\ldots}   True  0.000020  
6              [0, 1, 2, 3, 4, 5, 6, 7, 8, 9, 10, 11]   True  0.000010  
7                           [0, 2, 3, 4, 5, 6, 8, 10]   True  0.000184  
8   [0, 1, 2, 3, 4, 5, 6, 7, 8, 9, 10, 11, 12, 13,{\ldots}   True  0.000024  
9                   [0, 1, 4, 5, 6, 7, 8, 10, 11, 12]   True  0.000105  
10        [0, 1, 4, 7, 8, 10, 11, 12, 14, 17, 18, 19]   True  0.069988  
11  [0, 1, 2, 3, 4, 5, 6, 7, 8, 9, 10, 11, 12, 13,{\ldots}   True  0.000026  
12        [0, 1, 2, 3, 5, 6, 7, 8, 9, 10, 11, 13, 14]   True  0.014062  
13                         [0, 1, 3, 6, 8, 9, 10, 11]   True  0.000540  
14                               [0, 4, 8, 9, 10, 18]   True  0.156535  
    \end{Verbatim}

    
    \begin{tcolorbox}[breakable, size=fbox, boxrule=1pt, pad at break*=1mm,colback=cellbackground, colframe=cellborder]
\prompt{In}{incolor}{274}{\boxspacing}
\begin{Verbatim}[commandchars=\\\{\}]
\PY{n}{bf\PYZus{}collab\PYZus{}results} \PY{o}{=} \PY{n}{evaluate\PYZus{}brute\PYZus{}force\PYZus{}on\PYZus{}dataset}\PY{p}{(}\PY{n}{collab\PYZus{}dataset}\PY{p}{,} \PY{n}{max\PYZus{}graphs}\PY{o}{=}\PY{l+m+mi}{20}\PY{p}{,} \PY{n}{max\PYZus{}nodes}\PY{o}{=}\PY{l+m+mi}{25}\PY{p}{)}

\PY{n}{display}\PY{p}{(}\PY{n}{bf\PYZus{}collab\PYZus{}results}\PY{p}{)}
\end{Verbatim}
\end{tcolorbox}

    \begin{Verbatim}[commandchars=\\\{\}]
Graph 0: Skipping (too large: 45 nodes)
Graph 1: Skipping (too large: 52 nodes)
Graph 2: Skipping (too large: 52 nodes)
Graph 3: Skipping (too large: 32 nodes)
Graph 4: Skipping (too large: 48 nodes)
Graph 5: Skipping (too large: 39 nodes)
Graph 6: Skipping (too large: 35 nodes)
Graph 7: Skipping (too large: 81 nodes)
Graph 8: Skipping (too large: 65 nodes)
Graph 9: Skipping (too large: 49 nodes)
Graph 10: Skipping (too large: 52 nodes)
Graph 11: Skipping (too large: 49 nodes)
Graph 12: Skipping (too large: 72 nodes)
Graph 13: Skipping (too large: 138 nodes)
Graph 14: Skipping (too large: 119 nodes)
Graph 15: Skipping (too large: 44 nodes)
Graph 16: Skipping (too large: 35 nodes)
Graph 17: Skipping (too large: 48 nodes)
Graph 18: Skipping (too large: 49 nodes)
Graph 19: Skipping (too large: 68 nodes)
    \end{Verbatim}

    
    \begin{Verbatim}[commandchars=\\\{\}]
Empty DataFrame
Columns: []
Index: []
    \end{Verbatim}

    
    \subsubsection{5.4 Informe del ejercicio}\label{informe-del-ejercicio}

En la comprobación mínima se ha recuperado la clique
\texttt{{[}0,\ 1,\ 2{]}}, con tamaño \texttt{3} y validación positiva,
por lo que la implementación básica ha quedado verificada en el caso
sintético del enunciado. Ese primer resultado es importante porque
confirma que las dos piezas esenciales del bloque exacto,
\texttt{is\_clique()} y \texttt{brute\_force\_max\_clique()}, están
coordinadas correctamente: la primera reconoce bien la propiedad de
clique y la segunda sabe detenerse en cuanto encuentra la mejor solución
posible.

En la salida de \textbf{IMDB-BINARY} se ve que se han omitido
exactamente 5 grafos por tamaño: \texttt{1}, \texttt{3}, \texttt{5},
\texttt{14} y \texttt{19}. Eso significa que, incluso restringiéndose a
solo 20 grafos, ya aparecen varias instancias que dejan de ser
razonables para fuerza bruta con el umbral fijado de \texttt{25} nodos.
En los 15 grafos que sí se han podido resolver, los tamaños de clique
van desde \texttt{6} nodos en el grafo \texttt{18} hasta \texttt{20}
nodos en los grafos \texttt{11} y \texttt{15}. Ese rango indica que el
subconjunto evaluado es bastante heterogéneo: hay grafos donde la
conectividad interna es moderada y la mejor clique sigue siendo
relativamente pequeña, pero también aparecen grafos casi completos donde
la clique máxima coincide prácticamente con todo el conjunto de nodos.
Por eso la tabla no solo da una medida de rendimiento, sino también una
primera intuición sobre la variedad estructural del dataset.

También se aprecia que el coste temporal no depende solo del número de
nodos, sino de cómo está distribuida la conectividad. Algunos casos se
han resuelto casi instantáneamente, como los grafos \texttt{6} y
\texttt{9}, ambos con \texttt{12} nodos y tiempo alrededor de
\texttt{0.00001\ s}, lo que sugiere que la clique máxima se identifica
muy pronto porque el grafo es extremadamente denso. En cambio, otros
grafos con tamaños parecidos o solo ligeramente mayores han sido mucho
más caros, como el grafo \texttt{2} con \texttt{21} nodos y
\texttt{0.5528\ s}, el grafo \texttt{0} con \texttt{20} nodos y
\texttt{0.3355\ s}, o el grafo \texttt{18} con \texttt{19} nodos y
\texttt{0.1565\ s}. Esto encaja con el comportamiento esperado de la
fuerza bruta: el algoritmo recorre combinaciones de nodos y su coste se
dispara cuando necesita explorar muchos subconjuntos antes de poder
certificar cuál es la clique máxima. Dicho de otro modo, no basta con
que el grafo tenga pocos nodos; también influye mucho si la solución
grande aparece pronto o si la estructura obliga a descartar una gran
cantidad de candidatos.

En \textbf{COLLAB} la limitación ha sido todavía más clara, porque los
20 primeros grafos se han descartado todos por superar el umbral de
\texttt{25} nodos. En los mensajes de salida se ve que esos tamaños van
desde \texttt{32} hasta \texttt{138} nodos, de modo que aquí ni siquiera
se llega a construir una tabla de resultados exactos. Eso tiene una
lectura metodológica importante: la fuerza bruta sí sirve como
referencia en instancias pequeñas y controladas, pero deja de ser una
herramienta práctica en cuanto el tamaño del grafo crece un poco.
Precisamente por eso tiene sentido el resto de la práctica. Este primer
ejercicio no solo valida la implementación exacta, sino que justifica
experimentalmente la necesidad de pasar después a formulaciones
continuas y métodos aproximados como Motzkin-Straus, RD y el enfoque
híbrido con GNN.

    \subsection{6. Ejercicio 2: Formulación de Motzkin-Straus (Parte
1)}\label{ejercicio-2-formulaciuxf3n-de-motzkin-straus-parte-1}

    \subsubsection{6.1 Implementación de las dos funciones
objetivo}\label{implementaciuxf3n-de-las-dos-funciones-objetivo}

La implementación de \texttt{motzkin\_straus\_objective()} y
\texttt{regularized\_objective()} ya se completó en la \textbf{tarea
1.1} de la sección de implementación. En este ejercicio no se vuelve a
copiar ese código, sino que se reutiliza directamente en los apartados
de comprobación numérica.

    \paragraph{6.2 Crear 3 grafos de test con máximas cliques
conocidas}\label{crear-3-grafos-de-test-con-muxe1ximas-cliques-conocidas}

    \begin{tcolorbox}[breakable, size=fbox, boxrule=1pt, pad at break*=1mm,colback=cellbackground, colframe=cellborder]
\prompt{In}{incolor}{275}{\boxspacing}
\begin{Verbatim}[commandchars=\\\{\}]
\PY{n}{A1} \PY{o}{=} \PY{n}{np}\PY{o}{.}\PY{n}{array}\PY{p}{(}\PY{p}{[}
    \PY{p}{[}\PY{l+m+mi}{0}\PY{p}{,} \PY{l+m+mi}{1}\PY{p}{,} \PY{l+m+mi}{1}\PY{p}{]}\PY{p}{,}
    \PY{p}{[}\PY{l+m+mi}{1}\PY{p}{,} \PY{l+m+mi}{0}\PY{p}{,} \PY{l+m+mi}{1}\PY{p}{]}\PY{p}{,}
    \PY{p}{[}\PY{l+m+mi}{1}\PY{p}{,} \PY{l+m+mi}{1}\PY{p}{,} \PY{l+m+mi}{0}\PY{p}{]}
\PY{p}{]}\PY{p}{,} \PY{n}{dtype}\PY{o}{=}\PY{n+nb}{float}\PY{p}{)}

\PY{n}{A2} \PY{o}{=} \PY{n}{np}\PY{o}{.}\PY{n}{array}\PY{p}{(}\PY{p}{[}
    \PY{p}{[}\PY{l+m+mi}{0}\PY{p}{,} \PY{l+m+mi}{1}\PY{p}{,} \PY{l+m+mi}{1}\PY{p}{,} \PY{l+m+mi}{1}\PY{p}{]}\PY{p}{,}
    \PY{p}{[}\PY{l+m+mi}{1}\PY{p}{,} \PY{l+m+mi}{0}\PY{p}{,} \PY{l+m+mi}{1}\PY{p}{,} \PY{l+m+mi}{0}\PY{p}{]}\PY{p}{,}
    \PY{p}{[}\PY{l+m+mi}{1}\PY{p}{,} \PY{l+m+mi}{1}\PY{p}{,} \PY{l+m+mi}{0}\PY{p}{,} \PY{l+m+mi}{1}\PY{p}{]}\PY{p}{,}
    \PY{p}{[}\PY{l+m+mi}{1}\PY{p}{,} \PY{l+m+mi}{0}\PY{p}{,} \PY{l+m+mi}{1}\PY{p}{,} \PY{l+m+mi}{0}\PY{p}{]}
\PY{p}{]}\PY{p}{,} \PY{n}{dtype}\PY{o}{=}\PY{n+nb}{float}\PY{p}{)}

\PY{n}{A3} \PY{o}{=} \PY{n}{np}\PY{o}{.}\PY{n}{array}\PY{p}{(}\PY{p}{[}
    \PY{p}{[}\PY{l+m+mi}{0}\PY{p}{,} \PY{l+m+mi}{0}\PY{p}{,} \PY{l+m+mi}{1}\PY{p}{]}\PY{p}{,}
    \PY{p}{[}\PY{l+m+mi}{0}\PY{p}{,} \PY{l+m+mi}{0}\PY{p}{,} \PY{l+m+mi}{1}\PY{p}{]}\PY{p}{,}
    \PY{p}{[}\PY{l+m+mi}{1}\PY{p}{,} \PY{l+m+mi}{1}\PY{p}{,} \PY{l+m+mi}{0}\PY{p}{]}
\PY{p}{]}\PY{p}{,} \PY{n}{dtype}\PY{o}{=}\PY{n+nb}{float}\PY{p}{)}

\PY{n}{exercise2\PYZus{}graphs} \PY{o}{=} \PY{p}{\PYZob{}}
    \PY{l+s+s2}{\PYZdq{}}\PY{l+s+s2}{Triángulo K3}\PY{l+s+s2}{\PYZdq{}}\PY{p}{:} \PY{n}{A1}\PY{p}{,}
    \PY{l+s+s2}{\PYZdq{}}\PY{l+s+s2}{Cuadrado con diagonal}\PY{l+s+s2}{\PYZdq{}}\PY{p}{:} \PY{n}{A2}\PY{p}{,}
    \PY{l+s+s2}{\PYZdq{}}\PY{l+s+s2}{Cherry}\PY{l+s+s2}{\PYZdq{}}\PY{p}{:} \PY{n}{A3}\PY{p}{,}
\PY{p}{\PYZcb{}}

\PY{k}{for} \PY{n}{name}\PY{p}{,} \PY{n}{A} \PY{o+ow}{in} \PY{n}{exercise2\PYZus{}graphs}\PY{o}{.}\PY{n}{items}\PY{p}{(}\PY{p}{)}\PY{p}{:}
    \PY{n}{clique} \PY{o}{=} \PY{n}{brute\PYZus{}force\PYZus{}max\PYZus{}clique}\PY{p}{(}\PY{n}{A}\PY{p}{)}
    \PY{n+nb}{print}\PY{p}{(}\PY{l+s+sa}{f}\PY{l+s+s2}{\PYZdq{}}\PY{l+s+si}{\PYZob{}}\PY{n}{name}\PY{l+s+si}{\PYZcb{}}\PY{l+s+s2}{: omega(G) = }\PY{l+s+si}{\PYZob{}}\PY{n+nb}{len}\PY{p}{(}\PY{n}{clique}\PY{p}{)}\PY{l+s+si}{\PYZcb{}}\PY{l+s+s2}{, clique = }\PY{l+s+si}{\PYZob{}}\PY{n}{clique}\PY{l+s+si}{\PYZcb{}}\PY{l+s+s2}{\PYZdq{}}\PY{p}{)}
\end{Verbatim}
\end{tcolorbox}

    \begin{Verbatim}[commandchars=\\\{\}]
Triángulo K3: omega(G) = 3, clique = [0, 1, 2]
Cuadrado con diagonal: omega(G) = 3, clique = [0, 1, 2]
Cherry: omega(G) = 2, clique = [0, 2]
    \end{Verbatim}

    \subsubsection{6.3 Verificar que el teorema se cumple en cada
caso}\label{verificar-que-el-teorema-se-cumple-en-cada-caso}

    \begin{tcolorbox}[breakable, size=fbox, boxrule=1pt, pad at break*=1mm,colback=cellbackground, colframe=cellborder]
\prompt{In}{incolor}{276}{\boxspacing}
\begin{Verbatim}[commandchars=\\\{\}]
\PY{n}{ms\PYZus{}df} \PY{o}{=} \PY{n}{verify\PYZus{}motzkin\PYZus{}straus}\PY{p}{(}\PY{p}{)}
\PY{n}{ms\PYZus{}theorem\PYZus{}df} \PY{o}{=} \PY{n}{ms\PYZus{}df}\PY{o}{.}\PY{n}{iloc}\PY{p}{[}\PY{p}{:}\PY{l+m+mi}{2}\PY{p}{]}\PY{o}{.}\PY{n}{copy}\PY{p}{(}\PY{p}{)}
\PY{n}{ms\PYZus{}theorem\PYZus{}df}\PY{p}{[}\PY{l+s+s2}{\PYZdq{}}\PY{l+s+s2}{error\PYZus{}abs}\PY{l+s+s2}{\PYZdq{}}\PY{p}{]} \PY{o}{=} \PY{p}{(}\PY{n}{ms\PYZus{}theorem\PYZus{}df}\PY{p}{[}\PY{l+s+s2}{\PYZdq{}}\PY{l+s+s2}{f\PYZus{}x}\PY{l+s+s2}{\PYZdq{}}\PY{p}{]} \PY{o}{\PYZhy{}} \PY{n}{ms\PYZus{}theorem\PYZus{}df}\PY{p}{[}\PY{l+s+s2}{\PYZdq{}}\PY{l+s+s2}{theoretical}\PY{l+s+s2}{\PYZdq{}}\PY{p}{]}\PY{p}{)}\PY{o}{.}\PY{n}{abs}\PY{p}{(}\PY{p}{)}
\PY{n}{display}\PY{p}{(}\PY{n}{ms\PYZus{}theorem\PYZus{}df}\PY{p}{)}
\end{Verbatim}
\end{tcolorbox}

    
    \begin{Verbatim}[commandchars=\\\{\}]
                   graph  omega     f\_x  theoretical  regularized\_f\_x  \textbackslash{}
0           Triángulo K3      3  0.6667       0.6667           0.8333   
1  Cuadrado con diagonal      3  0.6667       0.6667           0.8333   

    error\_abs  
0  1.1102e-16  
1  1.1102e-16  
    \end{Verbatim}

    
    \subsubsection{6.4 Demostrar el problema de la solución espuria en el
grafo
cherry}\label{demostrar-el-problema-de-la-soluciuxf3n-espuria-en-el-grafo-cherry}

    \begin{tcolorbox}[breakable, size=fbox, boxrule=1pt, pad at break*=1mm,colback=cellbackground, colframe=cellborder]
\prompt{In}{incolor}{277}{\boxspacing}
\begin{Verbatim}[commandchars=\\\{\}]
\PY{n}{true\PYZus{}x} \PY{o}{=} \PY{n}{np}\PY{o}{.}\PY{n}{array}\PY{p}{(}\PY{p}{[}\PY{l+m+mf}{0.5}\PY{p}{,} \PY{l+m+mf}{0.0}\PY{p}{,} \PY{l+m+mf}{0.5}\PY{p}{]}\PY{p}{,} \PY{n}{dtype}\PY{o}{=}\PY{n+nb}{float}\PY{p}{)}
\PY{n}{spurious\PYZus{}x} \PY{o}{=} \PY{n}{np}\PY{o}{.}\PY{n}{array}\PY{p}{(}\PY{p}{[}\PY{l+m+mf}{0.25}\PY{p}{,} \PY{l+m+mf}{0.25}\PY{p}{,} \PY{l+m+mf}{0.5}\PY{p}{]}\PY{p}{,} \PY{n}{dtype}\PY{o}{=}\PY{n+nb}{float}\PY{p}{)}

\PY{n+nb}{print}\PY{p}{(}\PY{l+s+s2}{\PYZdq{}}\PY{l+s+s2}{f(x\PYZus{}true) =}\PY{l+s+s2}{\PYZdq{}}\PY{p}{,} \PY{n}{motzkin\PYZus{}straus\PYZus{}objective}\PY{p}{(}\PY{n}{true\PYZus{}x}\PY{p}{,} \PY{n}{A3}\PY{p}{)}\PY{p}{)}
\PY{n+nb}{print}\PY{p}{(}\PY{l+s+s2}{\PYZdq{}}\PY{l+s+s2}{f(x\PYZus{}spurious) =}\PY{l+s+s2}{\PYZdq{}}\PY{p}{,} \PY{n}{motzkin\PYZus{}straus\PYZus{}objective}\PY{p}{(}\PY{n}{spurious\PYZus{}x}\PY{p}{,} \PY{n}{A3}\PY{p}{)}\PY{p}{)}
\end{Verbatim}
\end{tcolorbox}

    \begin{Verbatim}[commandchars=\\\{\}]
f(x\_true) = 0.5
f(x\_spurious) = 0.5
    \end{Verbatim}

    \subsubsection{6.5 Demostrar que la regularización lo
corrige}\label{demostrar-que-la-regularizaciuxf3n-lo-corrige}

    \begin{tcolorbox}[breakable, size=fbox, boxrule=1pt, pad at break*=1mm,colback=cellbackground, colframe=cellborder]
\prompt{In}{incolor}{278}{\boxspacing}
\begin{Verbatim}[commandchars=\\\{\}]
\PY{n+nb}{print}\PY{p}{(}\PY{l+s+s2}{\PYZdq{}}\PY{l+s+s2}{f\PYZus{}hat(x\PYZus{}true) =}\PY{l+s+s2}{\PYZdq{}}\PY{p}{,} \PY{n}{regularized\PYZus{}objective}\PY{p}{(}\PY{n}{true\PYZus{}x}\PY{p}{,} \PY{n}{A3}\PY{p}{)}\PY{p}{)}
\PY{n+nb}{print}\PY{p}{(}\PY{l+s+s2}{\PYZdq{}}\PY{l+s+s2}{f\PYZus{}hat(x\PYZus{}spurious) =}\PY{l+s+s2}{\PYZdq{}}\PY{p}{,} \PY{n}{regularized\PYZus{}objective}\PY{p}{(}\PY{n}{spurious\PYZus{}x}\PY{p}{,} \PY{n}{A3}\PY{p}{)}\PY{p}{)}
\PY{n+nb}{print}\PY{p}{(}\PY{l+s+s2}{\PYZdq{}}\PY{l+s+s2}{Diferencia a favor de la solución verdadera =}\PY{l+s+s2}{\PYZdq{}}\PY{p}{,} \PY{n}{regularized\PYZus{}objective}\PY{p}{(}\PY{n}{true\PYZus{}x}\PY{p}{,} \PY{n}{A3}\PY{p}{)} \PY{o}{\PYZhy{}} \PY{n}{regularized\PYZus{}objective}\PY{p}{(}\PY{n}{spurious\PYZus{}x}\PY{p}{,} \PY{n}{A3}\PY{p}{)}\PY{p}{)}
\end{Verbatim}
\end{tcolorbox}

    \begin{Verbatim}[commandchars=\\\{\}]
f\_hat(x\_true) = 0.75
f\_hat(x\_spurious) = 0.6875
Diferencia a favor de la solución verdadera = 0.0625
    \end{Verbatim}

    \subsubsection{6.6 Informe del ejercicio}\label{informe-del-ejercicio}

En la celda de creación de ejemplos se ha visto que los tres grafos
elegidos cumplen justo el papel que se buscaba: el triángulo tiene una
clique máxima de tamaño \texttt{3}, el cuadrado con diagonal también
tiene máxima clique \texttt{3}, y el cherry tiene máxima clique
\texttt{2}. Es decir, se ha trabajado con casos pequeños donde se conoce
exactamente la solución y donde la comprobación teórica puede hacerse de
forma completamente transparente.

En la verificación del teorema, los dos casos sin ambigüedad han
encajado perfectamente con la predicción de Motzkin-Straus. Tanto en el
triángulo K3 como en el cuadrado con diagonal se ha obtenido
\texttt{f(x)\ =\ 0.6667}, exactamente igual al valor teórico
\texttt{1\ -\ 1/omega(G)\ =\ 0.6667}, con error absoluto prácticamente
nulo. Esto significa que, al menos en estos ejemplos, la relajación
continua no solo conserva la información combinatoria sobre la clique
máxima, sino que la expresa de forma numéricamente estable.

El caso del cherry ha sido el más instructivo. En la celda 6.4 se ha
comprobado que \texttt{f(x\_true)\ =\ 0.5} y
\texttt{f(x\_spurious)\ =\ 0.5}, de modo que la formulación sin
regularizar no distingue entre la solución verdadera y una distribución
espuria que reparte masa sobre un soporte que no es una clique. Esa es
precisamente la limitación conceptual importante del teorema cuando se
usa directamente como criterio de optimización.

La celda 6.5 muestra por qué se introduce la regularización:
\texttt{f\_hat(x\_true)\ =\ 0.75} frente a
\texttt{f\_hat(x\_spurious)\ =\ 0.6875}, con una diferencia positiva de
\texttt{0.0625} a favor de la solución correcta. Por tanto, en este
ejercicio no solo se ha verificado el teorema en ejemplos sencillos,
sino que también se ha justificado experimentalmente la necesidad de
añadir el término regularizador antes de pasar a Replicator Dynamics. La
regularización no cambia la intuición básica de Motzkin-Straus, pero sí
hace que el criterio sea más útil cuando se quiere optimizarlo en la
práctica.

    \subsection{7. Ejercicio 3: Dinámicas replicadoras (Parte
2)}\label{ejercicio-3-dinuxe1micas-replicadoras-parte-2}

    \subsubsection{7.1 Implementar RD con el
decodificador}\label{implementar-rd-con-el-decodificador}

La implementación de \texttt{replicator\_dynamics()} y
\texttt{decode\_clique()} ya se completó en la \textbf{tarea 2.1} y en
la \textbf{tarea 2.2} de la sección de implementación. En este ejercicio
se reutiliza ese código para visualizar la convergencia, evaluarlo sobre
IMDB y COLLAB y compararlo con la referencia exacta disponible en los
grafos pequeños.

    \paragraph{7.2 Visualización de la convergencia en un grafo
pequeño}\label{visualizaciuxf3n-de-la-convergencia-en-un-grafo-pequeuxf1o}

    \begin{tcolorbox}[breakable, size=fbox, boxrule=1pt, pad at break*=1mm,colback=cellbackground, colframe=cellborder]
\prompt{In}{incolor}{279}{\boxspacing}
\begin{Verbatim}[commandchars=\\\{\}]
\PY{n}{A\PYZus{}rd\PYZus{}demo} \PY{o}{=} \PY{n}{np}\PY{o}{.}\PY{n}{array}\PY{p}{(}
    \PY{p}{[}
        \PY{p}{[}\PY{l+m+mi}{0}\PY{p}{,} \PY{l+m+mi}{1}\PY{p}{,} \PY{l+m+mi}{1}\PY{p}{,} \PY{l+m+mi}{1}\PY{p}{]}\PY{p}{,}
        \PY{p}{[}\PY{l+m+mi}{1}\PY{p}{,} \PY{l+m+mi}{0}\PY{p}{,} \PY{l+m+mi}{1}\PY{p}{,} \PY{l+m+mi}{0}\PY{p}{]}\PY{p}{,}
        \PY{p}{[}\PY{l+m+mi}{1}\PY{p}{,} \PY{l+m+mi}{1}\PY{p}{,} \PY{l+m+mi}{0}\PY{p}{,} \PY{l+m+mi}{1}\PY{p}{]}\PY{p}{,}
        \PY{p}{[}\PY{l+m+mi}{1}\PY{p}{,} \PY{l+m+mi}{0}\PY{p}{,} \PY{l+m+mi}{1}\PY{p}{,} \PY{l+m+mi}{0}\PY{p}{]}\PY{p}{,}
    \PY{p}{]}\PY{p}{,}
    \PY{n}{dtype}\PY{o}{=}\PY{n+nb}{float}\PY{p}{,}
\PY{p}{)}
\PY{n}{rd\PYZus{}demo\PYZus{}x}\PY{p}{,} \PY{n}{rd\PYZus{}demo\PYZus{}history} \PY{o}{=} \PY{n}{visualize\PYZus{}rd}\PY{p}{(}\PY{n}{A\PYZus{}rd\PYZus{}demo}\PY{p}{,} \PY{n}{graph\PYZus{}name}\PY{o}{=}\PY{l+s+s2}{\PYZdq{}}\PY{l+s+s2}{Cuadrado con diagonal}\PY{l+s+s2}{\PYZdq{}}\PY{p}{)}
\PY{n}{rd\PYZus{}demo\PYZus{}clique} \PY{o}{=} \PY{n}{decode\PYZus{}clique}\PY{p}{(}\PY{n}{rd\PYZus{}demo\PYZus{}x}\PY{p}{,} \PY{n}{A\PYZus{}rd\PYZus{}demo}\PY{p}{)}
\PY{n+nb}{print}\PY{p}{(}\PY{l+s+s2}{\PYZdq{}}\PY{l+s+s2}{Clique decodificada:}\PY{l+s+s2}{\PYZdq{}}\PY{p}{,} \PY{n}{rd\PYZus{}demo\PYZus{}clique}\PY{p}{)}
\PY{n+nb}{print}\PY{p}{(}\PY{l+s+s2}{\PYZdq{}}\PY{l+s+s2}{Tamaño:}\PY{l+s+s2}{\PYZdq{}}\PY{p}{,} \PY{n+nb}{len}\PY{p}{(}\PY{n}{rd\PYZus{}demo\PYZus{}clique}\PY{p}{)}\PY{p}{)}
\PY{n+nb}{print}\PY{p}{(}\PY{l+s+s2}{\PYZdq{}}\PY{l+s+s2}{Iteraciones:}\PY{l+s+s2}{\PYZdq{}}\PY{p}{,} \PY{n}{rd\PYZus{}demo\PYZus{}history}\PY{p}{[}\PY{l+s+s2}{\PYZdq{}}\PY{l+s+s2}{iterations}\PY{l+s+s2}{\PYZdq{}}\PY{p}{]}\PY{p}{)}
\end{Verbatim}
\end{tcolorbox}

    \begin{center}
    \adjustimage{max size={0.9\linewidth}{0.9\paperheight}}{practica_01_jordi_blasco_lozano_files/practica_01_jordi_blasco_lozano_73_0.png}
    \end{center}
    { \hspace*{\fill} \\}
    
    \begin{center}
    \adjustimage{max size={0.9\linewidth}{0.9\paperheight}}{practica_01_jordi_blasco_lozano_files/practica_01_jordi_blasco_lozano_73_1.png}
    \end{center}
    { \hspace*{\fill} \\}
    
    \begin{Verbatim}[commandchars=\\\{\}]
Clique decodificada: [0, 2, 1]
Tamaño: 3
Iteraciones: 96
    \end{Verbatim}

    \subsubsection{7.3 Aplicación de RD en IMDB y
COLLAB}\label{aplicaciuxf3n-de-rd-en-imdb-y-collab}

    \begin{itemize}
\tightlist
\item
  IMDB
\end{itemize}

    \begin{tcolorbox}[breakable, size=fbox, boxrule=1pt, pad at break*=1mm,colback=cellbackground, colframe=cellborder]
\prompt{In}{incolor}{293}{\boxspacing}
\begin{Verbatim}[commandchars=\\\{\}]
\PY{n}{rd\PYZus{}imdb\PYZus{}results} \PY{o}{=} \PY{n}{evaluate\PYZus{}rd\PYZus{}on\PYZus{}dataset}\PY{p}{(}\PY{n}{imdb\PYZus{}dataset}\PY{p}{,} \PY{l+s+s2}{\PYZdq{}}\PY{l+s+s2}{IMDB\PYZhy{}BINARY}\PY{l+s+s2}{\PYZdq{}}\PY{p}{,} \PY{n}{num\PYZus{}graphs}\PY{o}{=}\PY{l+m+mi}{50}\PY{p}{)}

\PY{n}{rd\PYZus{}imdb\PYZus{}df}\PY{p}{,} \PY{n}{rd\PYZus{}imdb\PYZus{}summary} \PY{o}{=} \PY{n}{summarize\PYZus{}results}\PY{p}{(}\PY{n}{rd\PYZus{}imdb\PYZus{}results}\PY{p}{)}

\PY{n}{display}\PY{p}{(}\PY{n}{rd\PYZus{}imdb\PYZus{}df}\PY{o}{.}\PY{n}{head}\PY{p}{(}\PY{l+m+mi}{10}\PY{p}{)}\PY{p}{)}

\PY{n+nb}{print}\PY{p}{(}\PY{l+s+sa}{f}\PY{l+s+s2}{\PYZdq{}}\PY{l+s+s2}{Clique media IMDB: }\PY{l+s+si}{\PYZob{}}\PY{n}{rd\PYZus{}imdb\PYZus{}summary}\PY{p}{[}\PY{l+s+s1}{\PYZsq{}}\PY{l+s+s1}{clique\PYZus{}size}\PY{l+s+s1}{\PYZsq{}}\PY{p}{]}\PY{p}{[}\PY{l+s+s1}{\PYZsq{}}\PY{l+s+s1}{mean}\PY{l+s+s1}{\PYZsq{}}\PY{p}{]}\PY{l+s+si}{:}\PY{l+s+s2}{.4f}\PY{l+s+si}{\PYZcb{}}\PY{l+s+s2}{\PYZdq{}}\PY{p}{)}

\PY{n+nb}{print}\PY{p}{(}\PY{l+s+sa}{f}\PY{l+s+s2}{\PYZdq{}}\PY{l+s+s2}{Tiempo medio IMDB: }\PY{l+s+si}{\PYZob{}}\PY{n}{rd\PYZus{}imdb\PYZus{}summary}\PY{p}{[}\PY{l+s+s1}{\PYZsq{}}\PY{l+s+s1}{time}\PY{l+s+s1}{\PYZsq{}}\PY{p}{]}\PY{p}{[}\PY{l+s+s1}{\PYZsq{}}\PY{l+s+s1}{mean}\PY{l+s+s1}{\PYZsq{}}\PY{p}{]}\PY{l+s+si}{:}\PY{l+s+s2}{.4f}\PY{l+s+si}{\PYZcb{}}\PY{l+s+s2}{ s}\PY{l+s+s2}{\PYZdq{}}\PY{p}{)}
\end{Verbatim}
\end{tcolorbox}

    
    \begin{Verbatim}[commandchars=\\\{\}]
   graph\_idx  num\_nodes  clique\_size  \textbackslash{}
0          0         20            9   
1          1         32           11   
2          2         21           10   
3          3         35            7   
4          4         14            9   
5          5         63            8   
6          6         12           12   
7          7         15            9   
8          8         18           18   
9          9         12           12   

                                              clique  valid        time  \textbackslash{}
0                   [2, 19, 6, 12, 14, 1, 17, 18, 8]   True  4.9584e-03   
1          [14, 23, 13, 3, 26, 5, 4, 29, 20, 21, 28]   True  4.8256e-03   
2                 [2, 5, 11, 19, 7, 6, 12, 3, 0, 20]   True  4.4360e-03   
3                          [25, 17, 8, 2, 4, 30, 33]   True  3.8559e-03   
4                     [0, 11, 8, 2, 12, 9, 6, 5, 13]   True  3.9346e-03   
5                    [40, 4, 13, 10, 49, 50, 47, 23]   True  4.0765e-03   
6             [0, 1, 2, 3, 4, 5, 6, 7, 8, 9, 10, 11]   True  7.0095e-05   
7                       [1, 5, 6, 14, 9, 8, 7, 0, 4]   True  4.1919e-03   
8  [17, 16, 0, 1, 4, 5, 2, 3, 8, 9, 10, 11, 12, 1{\ldots}   True  8.2016e-05   
9             [0, 1, 2, 3, 4, 5, 6, 7, 8, 9, 10, 11]   True  6.3181e-05   

   iterations  
0         406  
1         491  
2         427  
3         377  
4         395  
5         373  
6           1  
7         402  
8           1  
9           1  
    \end{Verbatim}

    
    \begin{Verbatim}[commandchars=\\\{\}]
Clique media IMDB: 10.3800
Tiempo medio IMDB: 0.0033 s
    \end{Verbatim}

    \begin{itemize}
\tightlist
\item
  COLLAB
\end{itemize}

    \begin{tcolorbox}[breakable, size=fbox, boxrule=1pt, pad at break*=1mm,colback=cellbackground, colframe=cellborder]
\prompt{In}{incolor}{294}{\boxspacing}
\begin{Verbatim}[commandchars=\\\{\}]
\PY{n}{rd\PYZus{}collab\PYZus{}results} \PY{o}{=} \PY{n}{evaluate\PYZus{}rd\PYZus{}on\PYZus{}dataset}\PY{p}{(}\PY{n}{collab\PYZus{}dataset}\PY{p}{,} \PY{l+s+s2}{\PYZdq{}}\PY{l+s+s2}{COLLAB}\PY{l+s+s2}{\PYZdq{}}\PY{p}{,} \PY{n}{num\PYZus{}graphs}\PY{o}{=}\PY{l+m+mi}{50}\PY{p}{)}

\PY{n}{rd\PYZus{}collab\PYZus{}df}\PY{p}{,} \PY{n}{rd\PYZus{}collab\PYZus{}summary} \PY{o}{=} \PY{n}{summarize\PYZus{}results}\PY{p}{(}\PY{n}{rd\PYZus{}collab\PYZus{}results}\PY{p}{)}

\PY{n}{display}\PY{p}{(}\PY{n}{rd\PYZus{}collab\PYZus{}df}\PY{o}{.}\PY{n}{head}\PY{p}{(}\PY{l+m+mi}{10}\PY{p}{)}\PY{p}{)}

\PY{n+nb}{print}\PY{p}{(}\PY{l+s+sa}{f}\PY{l+s+s2}{\PYZdq{}}\PY{l+s+s2}{Clique media COLLAB: }\PY{l+s+si}{\PYZob{}}\PY{n}{rd\PYZus{}collab\PYZus{}summary}\PY{p}{[}\PY{l+s+s1}{\PYZsq{}}\PY{l+s+s1}{clique\PYZus{}size}\PY{l+s+s1}{\PYZsq{}}\PY{p}{]}\PY{p}{[}\PY{l+s+s1}{\PYZsq{}}\PY{l+s+s1}{mean}\PY{l+s+s1}{\PYZsq{}}\PY{p}{]}\PY{l+s+si}{:}\PY{l+s+s2}{.4f}\PY{l+s+si}{\PYZcb{}}\PY{l+s+s2}{\PYZdq{}}\PY{p}{)}

\PY{n+nb}{print}\PY{p}{(}\PY{l+s+sa}{f}\PY{l+s+s2}{\PYZdq{}}\PY{l+s+s2}{Tiempo medio COLLAB: }\PY{l+s+si}{\PYZob{}}\PY{n}{rd\PYZus{}collab\PYZus{}summary}\PY{p}{[}\PY{l+s+s1}{\PYZsq{}}\PY{l+s+s1}{time}\PY{l+s+s1}{\PYZsq{}}\PY{p}{]}\PY{p}{[}\PY{l+s+s1}{\PYZsq{}}\PY{l+s+s1}{mean}\PY{l+s+s1}{\PYZsq{}}\PY{p}{]}\PY{l+s+si}{:}\PY{l+s+s2}{.4f}\PY{l+s+si}{\PYZcb{}}\PY{l+s+s2}{ s}\PY{l+s+s2}{\PYZdq{}}\PY{p}{)}
\end{Verbatim}
\end{tcolorbox}

    
    \begin{Verbatim}[commandchars=\\\{\}]
   graph\_idx  num\_nodes  clique\_size  \textbackslash{}
0          0         45           45   
1          1         52           36   
2          2         52           42   
3          3         32           32   
4          4         48           43   
5          5         39            7   
6          6         35           35   
7          7         81           25   
8          8         65           25   
9          9         49           40   

                                              clique  valid    time  \textbackslash{}
0  [0, 1, 2, 3, 4, 5, 6, 7, 8, 9, 10, 11, 12, 13,{\ldots}   True  0.0003   
1  [2, 10, 22, 36, 24, 26, 50, 14, 13, 5, 17, 7, {\ldots}   True  0.0052   
2  [35, 48, 34, 38, 24, 4, 18, 22, 20, 49, 37, 41{\ldots}   True  0.0055   
3  [0, 1, 2, 3, 4, 5, 6, 7, 8, 9, 10, 11, 12, 13,{\ldots}   True  0.0001   
4  [15, 39, 40, 20, 24, 13, 9, 25, 33, 41, 45, 31{\ldots}   True  0.0051   
5                          [21, 23, 3, 34, 33, 1, 5]   True  0.0034   
6  [0, 1, 2, 3, 4, 5, 6, 7, 8, 9, 10, 11, 12, 13,{\ldots}   True  0.0001   
7  [41, 70, 75, 22, 79, 27, 24, 53, 35, 69, 30, 4{\ldots}   True  0.0057   
8  [53, 15, 38, 62, 58, 21, 34, 23, 16, 59, 19, 3{\ldots}   True  0.0052   
9  [40, 7, 32, 0, 36, 14, 4, 30, 34, 28, 46, 48, {\ldots}   True  0.0052   

   iterations  
0           1  
1         500  
2         500  
3           1  
4         500  
5         339  
6           1  
7         500  
8         500  
9         500  
    \end{Verbatim}

    
    \begin{Verbatim}[commandchars=\\\{\}]
Clique media COLLAB: 27.3000
Tiempo medio COLLAB: 0.0044 s
    \end{Verbatim}

    \subsubsection{7.4 Comparación con fuerza bruta en los grafos pequeños
de
IMDB}\label{comparaciuxf3n-con-fuerza-bruta-en-los-grafos-pequeuxf1os-de-imdb}

    \begin{tcolorbox}[breakable, size=fbox, boxrule=1pt, pad at break*=1mm,colback=cellbackground, colframe=cellborder]
\prompt{In}{incolor}{282}{\boxspacing}
\begin{Verbatim}[commandchars=\\\{\}]
\PY{n}{rd\PYZus{}vs\PYZus{}bf\PYZus{}imdb} \PY{o}{=} \PY{n}{compare\PYZus{}rd\PYZus{}with\PYZus{}bruteforce}\PY{p}{(}\PY{n}{rd\PYZus{}imdb\PYZus{}results}\PY{p}{,} \PY{n}{bf\PYZus{}imdb\PYZus{}results}\PY{p}{)}
\PY{n}{display}\PY{p}{(}\PY{n}{rd\PYZus{}vs\PYZus{}bf\PYZus{}imdb}\PY{p}{)}
\end{Verbatim}
\end{tcolorbox}

    
    \begin{Verbatim}[commandchars=\\\{\}]
    graph\_idx  bf\_clique\_size  rd\_clique\_size  same\_size  rd\_valid
0           0               9               9       True      True
1           2              10              10       True      True
2           4               9               9       True      True
3           6              12              12       True      True
4           7               9               9       True      True
5           8              18              18       True      True
6           9              12              12       True      True
7          10               8               8       True      True
8          11              20              20       True      True
9          12              10              10       True      True
10         13              12              12       True      True
11         15              20              20       True      True
12         16              13              13       True      True
13         17               8               8       True      True
14         18               6               6       True      True
    \end{Verbatim}

    
    \subsubsection{7.5 Informe del ejercicio}\label{informe-del-ejercicio}

En el ejemplo del cuadrado con diagonal, la dinámica ha convergido en
\texttt{96} iteraciones y la clique decodificada ha sido
\texttt{{[}0,\ 2,\ 1{]}}, de tamaño \texttt{3}, exactamente la esperada.
Ese resultado es importante porque muestra, en un caso pequeño y
totalmente interpretable, cómo la masa de la solución continua termina
concentrándose en los nodos que forman una subestructura densamente
conectada. La visualización del objetivo y de los pesos finales ayuda a
entender qué está ocurriendo: RD no prueba combinaciones de manera
exhaustiva como la fuerza bruta, sino que redistribuye iterativamente la
probabilidad hacia los nodos que reciben más apoyo del resto según la
matriz de adyacencia regularizada.

En IMDB-BINARY, la evaluación sobre los primeros 50 grafos ha producido
cliques válidas en el \texttt{100\%} de los casos y una clique media de
\texttt{10.380}, con tiempo medio alrededor de \texttt{0.0033\ s}. En
COLLAB se ha mantenido también una validez del \texttt{100\%}, con una
clique media de \texttt{27.300} y un tiempo medio cercano a
\texttt{0.0044\ s}. Lo relevante no es solo que todas las soluciones
finales sean cliques válidas, sino que se ha conseguido esa validez con
un coste muy pequeño incluso en grafos donde la fuerza bruta ya no era
utilizable. Esto indica que el paso desde una formulación exacta a una
formulación continua no ha destruido la estructura del problema, sino
que ha permitido capturarla de una forma mucho más escalable.

La comparación con fuerza bruta en IMDB refuerza todavía más esa
interpretación. En los \texttt{15} grafos pequeños donde sí existía
ground truth exacto, RD ha coincidido en el tamaño de la clique en
\texttt{15\ de\ 15} casos, es decir, \texttt{100\%} de acuerdo sobre ese
subconjunto. Esa coincidencia no demuestra que RD sea exacto en general,
pero sí aporta una evidencia muy fuerte de que la relajación
regularizada está guiando la búsqueda hacia soluciones combinatoriamente
correctas cuando el problema todavía se puede verificar de manera
exhaustiva.

También se aprecia un patrón interesante en los tiempos y en el número
de iteraciones. En varios grafos muy densos, como algunos casos de IMDB
o COLLAB donde la clique recuperada coincide casi con todo el grafo, la
convergencia se ha producido casi de inmediato. En cambio, en otros
ejemplos la dinámica ha necesitado muchas más iteraciones e incluso ha
alcanzado con frecuencia el máximo fijado de \texttt{500}. Esto tiene
una explicación estructural: cuando la mejor clique destaca de forma muy
clara frente al resto del grafo, la dinámica se concentra rápidamente;
cuando la estructura es más ambigua y existen varias zonas densas
compitiendo entre sí, el proceso de concentración es más lento. Por
tanto, el ejercicio no solo muestra que RD funciona bien, sino también
por qué su comportamiento depende de la geometría del problema y no
únicamente del número de nodos.

    \subsection{8. Ejercicio 4: GNN + Dinámicas replicadoras (Parte
3)}\label{ejercicio-4-gnn-dinuxe1micas-replicadoras-parte-3}

    \subsubsection{8.1 Implementar el pipeline
completo}\label{implementar-el-pipeline-completo}

La implementación del pipeline completo ya se realizó en la
\textbf{sección 4.5}. Para mantener el orden, se resume aquí qué piezas
forman ese pipeline antes de pasar a los experimentos.

    \begin{itemize}
\item
  Características de nodo -\textgreater{} GNN

  En la \textbf{tarea 3.1} y en la \textbf{tarea 3.2} se implementaron
  \texttt{compute\_node\_features()} y \texttt{MaxCliqueGNN}. Esas
  piezas convierten cada grafo en un conjunto de features estructurales
  y generan un punto inicial del simplex mediante una salida
  \texttt{softmax} sobre los nodos.
\end{itemize}

    \begin{itemize}
\item
  RD diferenciable y función de pérdida

  En la \textbf{tarea 3.3} y en la \textbf{tarea 3.4} se implementaron
  \texttt{differentiable\_rd()} y \texttt{maxclique\_loss()}. Con ello,
  la red no se entrena solo para producir puntuaciones arbitrarias, sino
  para producir inicializaciones que sigan siendo útiles una vez
  refinadas con RD.
\end{itemize}

    \begin{itemize}
\item
  Entrenamiento y evaluación del pipeline

  En la \textbf{tarea 3.5} y en la \textbf{tarea 3.6} se implementaron
  \texttt{train\_gnn()} y \texttt{evaluate\_gnn()}. Eso deja completo el
  pipeline pedido en el enunciado: features, GNN, refinamiento con RD,
  decodificación final y medición de resultados.
\end{itemize}

    \subsubsection{8.2 Entrenar en IMDB y evaluar en grafos
hold-out}\label{entrenar-en-imdb-y-evaluar-en-grafos-hold-out}

En este apartado se prepara la partición hold-out y se entrena el modelo
híbrido. Aunque también se incluye COLLAB para enriquecer la
comparación.

    \begin{itemize}
\tightlist
\item
  Preparación del subconjunto experimental
\end{itemize}

    \begin{tcolorbox}[breakable, size=fbox, boxrule=1pt, pad at break*=1mm,colback=cellbackground, colframe=cellborder]
\prompt{In}{incolor}{283}{\boxspacing}
\begin{Verbatim}[commandchars=\\\{\}]
\PY{n}{imdb\PYZus{}train\PYZus{}graphs}\PY{p}{,} \PY{n}{imdb\PYZus{}test\PYZus{}graphs} \PY{o}{=} \PY{n}{split\PYZus{}dataset}\PY{p}{(}\PY{n}{imdb\PYZus{}dataset}\PY{p}{,} \PY{n}{train\PYZus{}ratio}\PY{o}{=}\PY{l+m+mf}{0.8}\PY{p}{,} \PY{n}{max\PYZus{}graphs}\PY{o}{=}\PY{l+m+mi}{240}\PY{p}{)}
\PY{n}{collab\PYZus{}train\PYZus{}graphs}\PY{p}{,} \PY{n}{collab\PYZus{}test\PYZus{}graphs} \PY{o}{=} \PY{n}{split\PYZus{}dataset}\PY{p}{(}\PY{n}{collab\PYZus{}dataset}\PY{p}{,} \PY{n}{train\PYZus{}ratio}\PY{o}{=}\PY{l+m+mf}{0.8}\PY{p}{,} \PY{n}{max\PYZus{}graphs}\PY{o}{=}\PY{l+m+mi}{240}\PY{p}{)}

\PY{n+nb}{print}\PY{p}{(}\PY{l+s+sa}{f}\PY{l+s+s2}{\PYZdq{}}\PY{l+s+s2}{IMDB \PYZhy{} train: }\PY{l+s+si}{\PYZob{}}\PY{n+nb}{len}\PY{p}{(}\PY{n}{imdb\PYZus{}train\PYZus{}graphs}\PY{p}{)}\PY{l+s+si}{\PYZcb{}}\PY{l+s+s2}{, test: }\PY{l+s+si}{\PYZob{}}\PY{n+nb}{len}\PY{p}{(}\PY{n}{imdb\PYZus{}test\PYZus{}graphs}\PY{p}{)}\PY{l+s+si}{\PYZcb{}}\PY{l+s+s2}{\PYZdq{}}\PY{p}{)}
\PY{n+nb}{print}\PY{p}{(}\PY{l+s+sa}{f}\PY{l+s+s2}{\PYZdq{}}\PY{l+s+s2}{COLLAB \PYZhy{} train: }\PY{l+s+si}{\PYZob{}}\PY{n+nb}{len}\PY{p}{(}\PY{n}{collab\PYZus{}train\PYZus{}graphs}\PY{p}{)}\PY{l+s+si}{\PYZcb{}}\PY{l+s+s2}{, test: }\PY{l+s+si}{\PYZob{}}\PY{n+nb}{len}\PY{p}{(}\PY{n}{collab\PYZus{}test\PYZus{}graphs}\PY{p}{)}\PY{l+s+si}{\PYZcb{}}\PY{l+s+s2}{\PYZdq{}}\PY{p}{)}
\end{Verbatim}
\end{tcolorbox}

    \begin{Verbatim}[commandchars=\\\{\}]
IMDB - train: 192, test: 48
COLLAB - train: 192, test: 48
    \end{Verbatim}

    \begin{itemize}
\tightlist
\item
  Entrenamiento del modelo híbrido
\end{itemize}

    \begin{tcolorbox}[breakable, size=fbox, boxrule=1pt, pad at break*=1mm,colback=cellbackground, colframe=cellborder]
\prompt{In}{incolor}{284}{\boxspacing}
\begin{Verbatim}[commandchars=\\\{\}]
\PY{n}{imdb\PYZus{}gnn\PYZus{}model} \PY{o}{=} \PY{n}{MaxCliqueGNN}\PY{p}{(}\PY{n}{input\PYZus{}dim}\PY{o}{=}\PY{l+m+mi}{3}\PY{p}{,} \PY{n}{hidden\PYZus{}dim}\PY{o}{=}\PY{l+m+mi}{32}\PY{p}{,} \PY{n}{num\PYZus{}layers}\PY{o}{=}\PY{l+m+mi}{2}\PY{p}{)}
\PY{n}{imdb\PYZus{}gnn\PYZus{}model}\PY{p}{,} \PY{n}{imdb\PYZus{}losses} \PY{o}{=} \PY{n}{train\PYZus{}gnn}\PY{p}{(}
    \PY{n}{imdb\PYZus{}gnn\PYZus{}model}\PY{p}{,}
    \PY{n}{imdb\PYZus{}train\PYZus{}graphs}\PY{p}{,}
    \PY{n}{num\PYZus{}epochs}\PY{o}{=}\PY{l+m+mi}{12}\PY{p}{,}
    \PY{n}{lr}\PY{o}{=}\PY{l+m+mf}{0.01}\PY{p}{,}
    \PY{n}{rd\PYZus{}train\PYZus{}iters}\PY{o}{=}\PY{l+m+mi}{3}\PY{p}{,}
\PY{p}{)}
\PY{n}{plot\PYZus{}training\PYZus{}losses}\PY{p}{(}\PY{n}{imdb\PYZus{}losses}\PY{p}{,} \PY{l+s+s2}{\PYZdq{}}\PY{l+s+s2}{Pérdida de entrenamiento \PYZhy{} IMDB (GNN + RD)}\PY{l+s+s2}{\PYZdq{}}\PY{p}{)}

\PY{n}{collab\PYZus{}gnn\PYZus{}model} \PY{o}{=} \PY{n}{MaxCliqueGNN}\PY{p}{(}\PY{n}{input\PYZus{}dim}\PY{o}{=}\PY{l+m+mi}{3}\PY{p}{,} \PY{n}{hidden\PYZus{}dim}\PY{o}{=}\PY{l+m+mi}{32}\PY{p}{,} \PY{n}{num\PYZus{}layers}\PY{o}{=}\PY{l+m+mi}{2}\PY{p}{)}
\PY{n}{collab\PYZus{}gnn\PYZus{}model}\PY{p}{,} \PY{n}{collab\PYZus{}losses} \PY{o}{=} \PY{n}{train\PYZus{}gnn}\PY{p}{(}
    \PY{n}{collab\PYZus{}gnn\PYZus{}model}\PY{p}{,}
    \PY{n}{collab\PYZus{}train\PYZus{}graphs}\PY{p}{,}
    \PY{n}{num\PYZus{}epochs}\PY{o}{=}\PY{l+m+mi}{12}\PY{p}{,}
    \PY{n}{lr}\PY{o}{=}\PY{l+m+mf}{0.01}\PY{p}{,}
    \PY{n}{rd\PYZus{}train\PYZus{}iters}\PY{o}{=}\PY{l+m+mi}{3}\PY{p}{,}
\PY{p}{)}
\PY{n}{plot\PYZus{}training\PYZus{}losses}\PY{p}{(}\PY{n}{collab\PYZus{}losses}\PY{p}{,} \PY{l+s+s2}{\PYZdq{}}\PY{l+s+s2}{Pérdida de entrenamiento \PYZhy{} COLLAB (GNN + RD)}\PY{l+s+s2}{\PYZdq{}}\PY{p}{)}
\end{Verbatim}
\end{tcolorbox}

    \begin{Verbatim}[commandchars=\\\{\}]
Época 3/12 - loss media: -0.8462
Época 6/12 - loss media: -0.8525
Época 9/12 - loss media: -0.8569
Época 12/12 - loss media: -0.8566
    \end{Verbatim}

    \begin{center}
    \adjustimage{max size={0.9\linewidth}{0.9\paperheight}}{practica_01_jordi_blasco_lozano_files/practica_01_jordi_blasco_lozano_91_1.png}
    \end{center}
    { \hspace*{\fill} \\}
    
    \begin{Verbatim}[commandchars=\\\{\}]
Época 3/12 - loss media: -0.9152
Época 6/12 - loss media: -0.9216
Época 9/12 - loss media: -0.9243
Época 12/12 - loss media: -0.9237
    \end{Verbatim}

    \begin{center}
    \adjustimage{max size={0.9\linewidth}{0.9\paperheight}}{practica_01_jordi_blasco_lozano_files/practica_01_jordi_blasco_lozano_91_3.png}
    \end{center}
    { \hspace*{\fill} \\}
    
    \begin{tcolorbox}[breakable, size=fbox, boxrule=1pt, pad at break*=1mm,colback=cellbackground, colframe=cellborder]
\prompt{In}{incolor}{285}{\boxspacing}
\begin{Verbatim}[commandchars=\\\{\}]
\PY{n}{training\PYZus{}summary\PYZus{}df} \PY{o}{=} \PY{n}{pd}\PY{o}{.}\PY{n}{DataFrame}\PY{p}{(}\PY{p}{[}
    \PY{p}{\PYZob{}}
        \PY{l+s+s1}{\PYZsq{}}\PY{l+s+s1}{Dataset}\PY{l+s+s1}{\PYZsq{}}\PY{p}{:} \PY{l+s+s1}{\PYZsq{}}\PY{l+s+s1}{IMDB}\PY{l+s+s1}{\PYZsq{}}\PY{p}{,}
        \PY{l+s+s1}{\PYZsq{}}\PY{l+s+s1}{Loss inicial}\PY{l+s+s1}{\PYZsq{}}\PY{p}{:} \PY{n+nb}{float}\PY{p}{(}\PY{n}{imdb\PYZus{}losses}\PY{p}{[}\PY{l+m+mi}{0}\PY{p}{]}\PY{p}{)}\PY{p}{,}
        \PY{l+s+s1}{\PYZsq{}}\PY{l+s+s1}{Loss final}\PY{l+s+s1}{\PYZsq{}}\PY{p}{:} \PY{n+nb}{float}\PY{p}{(}\PY{n}{imdb\PYZus{}losses}\PY{p}{[}\PY{o}{\PYZhy{}}\PY{l+m+mi}{1}\PY{p}{]}\PY{p}{)}\PY{p}{,}
        \PY{l+s+s1}{\PYZsq{}}\PY{l+s+s1}{Descenso neto}\PY{l+s+s1}{\PYZsq{}}\PY{p}{:} \PY{n+nb}{float}\PY{p}{(}\PY{n}{imdb\PYZus{}losses}\PY{p}{[}\PY{l+m+mi}{0}\PY{p}{]} \PY{o}{\PYZhy{}} \PY{n}{imdb\PYZus{}losses}\PY{p}{[}\PY{o}{\PYZhy{}}\PY{l+m+mi}{1}\PY{p}{]}\PY{p}{)}\PY{p}{,}
        \PY{l+s+s1}{\PYZsq{}}\PY{l+s+s1}{Mejor loss}\PY{l+s+s1}{\PYZsq{}}\PY{p}{:} \PY{n+nb}{float}\PY{p}{(}\PY{n}{np}\PY{o}{.}\PY{n}{min}\PY{p}{(}\PY{n}{imdb\PYZus{}losses}\PY{p}{)}\PY{p}{)}\PY{p}{,}
        \PY{l+s+s1}{\PYZsq{}}\PY{l+s+s1}{Épocas}\PY{l+s+s1}{\PYZsq{}}\PY{p}{:} \PY{n+nb}{int}\PY{p}{(}\PY{n+nb}{len}\PY{p}{(}\PY{n}{imdb\PYZus{}losses}\PY{p}{)}\PY{p}{)}\PY{p}{,}
    \PY{p}{\PYZcb{}}\PY{p}{,}
    \PY{p}{\PYZob{}}
        \PY{l+s+s1}{\PYZsq{}}\PY{l+s+s1}{Dataset}\PY{l+s+s1}{\PYZsq{}}\PY{p}{:} \PY{l+s+s1}{\PYZsq{}}\PY{l+s+s1}{COLLAB}\PY{l+s+s1}{\PYZsq{}}\PY{p}{,}
        \PY{l+s+s1}{\PYZsq{}}\PY{l+s+s1}{Loss inicial}\PY{l+s+s1}{\PYZsq{}}\PY{p}{:} \PY{n+nb}{float}\PY{p}{(}\PY{n}{collab\PYZus{}losses}\PY{p}{[}\PY{l+m+mi}{0}\PY{p}{]}\PY{p}{)}\PY{p}{,}
        \PY{l+s+s1}{\PYZsq{}}\PY{l+s+s1}{Loss final}\PY{l+s+s1}{\PYZsq{}}\PY{p}{:} \PY{n+nb}{float}\PY{p}{(}\PY{n}{collab\PYZus{}losses}\PY{p}{[}\PY{o}{\PYZhy{}}\PY{l+m+mi}{1}\PY{p}{]}\PY{p}{)}\PY{p}{,}
        \PY{l+s+s1}{\PYZsq{}}\PY{l+s+s1}{Descenso neto}\PY{l+s+s1}{\PYZsq{}}\PY{p}{:} \PY{n+nb}{float}\PY{p}{(}\PY{n}{collab\PYZus{}losses}\PY{p}{[}\PY{l+m+mi}{0}\PY{p}{]} \PY{o}{\PYZhy{}} \PY{n}{collab\PYZus{}losses}\PY{p}{[}\PY{o}{\PYZhy{}}\PY{l+m+mi}{1}\PY{p}{]}\PY{p}{)}\PY{p}{,}
        \PY{l+s+s1}{\PYZsq{}}\PY{l+s+s1}{Mejor loss}\PY{l+s+s1}{\PYZsq{}}\PY{p}{:} \PY{n+nb}{float}\PY{p}{(}\PY{n}{np}\PY{o}{.}\PY{n}{min}\PY{p}{(}\PY{n}{collab\PYZus{}losses}\PY{p}{)}\PY{p}{)}\PY{p}{,}
        \PY{l+s+s1}{\PYZsq{}}\PY{l+s+s1}{Épocas}\PY{l+s+s1}{\PYZsq{}}\PY{p}{:} \PY{n+nb}{int}\PY{p}{(}\PY{n+nb}{len}\PY{p}{(}\PY{n}{collab\PYZus{}losses}\PY{p}{)}\PY{p}{)}\PY{p}{,}
    \PY{p}{\PYZcb{}}\PY{p}{,}
\PY{p}{]}\PY{p}{)}

\PY{n}{display}\PY{p}{(}\PY{n}{training\PYZus{}summary\PYZus{}df}\PY{o}{.}\PY{n}{round}\PY{p}{(}\PY{l+m+mi}{4}\PY{p}{)}\PY{p}{)}
\end{Verbatim}
\end{tcolorbox}

    
    \begin{Verbatim}[commandchars=\\\{\}]
  Dataset  Loss inicial  Loss final  Descenso neto  Mejor loss  Épocas
0    IMDB       -0.8268     -0.8566         0.0298     -0.8579      12
1  COLLAB       -0.8954     -0.9237         0.0284     -0.9248      12
    \end{Verbatim}

    
    \begin{itemize}
\tightlist
\item
  Evaluación sobre grafos hold-out
\end{itemize}

    \begin{tcolorbox}[breakable, size=fbox, boxrule=1pt, pad at break*=1mm,colback=cellbackground, colframe=cellborder]
\prompt{In}{incolor}{286}{\boxspacing}
\begin{Verbatim}[commandchars=\\\{\}]
\PY{n}{gnn\PYZus{}imdb\PYZus{}results} \PY{o}{=} \PY{n}{evaluate\PYZus{}gnn}\PY{p}{(}\PY{n}{imdb\PYZus{}gnn\PYZus{}model}\PY{p}{,} \PY{n}{imdb\PYZus{}test\PYZus{}graphs}\PY{p}{,} \PY{n}{rd\PYZus{}test\PYZus{}iters}\PY{o}{=}\PY{l+m+mi}{100}\PY{p}{)}
\PY{n}{gnn\PYZus{}collab\PYZus{}results} \PY{o}{=} \PY{n}{evaluate\PYZus{}gnn}\PY{p}{(}\PY{n}{collab\PYZus{}gnn\PYZus{}model}\PY{p}{,} \PY{n}{collab\PYZus{}test\PYZus{}graphs}\PY{p}{,} \PY{n}{rd\PYZus{}test\PYZus{}iters}\PY{o}{=}\PY{l+m+mi}{100}\PY{p}{)}

\PY{n}{rd\PYZus{}imdb\PYZus{}holdout} \PY{o}{=} \PY{n}{evaluate\PYZus{}rd\PYZus{}on\PYZus{}dataset}\PY{p}{(}\PY{n}{imdb\PYZus{}test\PYZus{}graphs}\PY{p}{,} \PY{l+s+s2}{\PYZdq{}}\PY{l+s+s2}{IMDB\PYZhy{}HOLDOUT}\PY{l+s+s2}{\PYZdq{}}\PY{p}{,} \PY{n}{num\PYZus{}graphs}\PY{o}{=}\PY{n+nb}{len}\PY{p}{(}\PY{n}{imdb\PYZus{}test\PYZus{}graphs}\PY{p}{)}\PY{p}{)}
\PY{n}{rd\PYZus{}collab\PYZus{}holdout} \PY{o}{=} \PY{n}{evaluate\PYZus{}rd\PYZus{}on\PYZus{}dataset}\PY{p}{(}\PY{n}{collab\PYZus{}test\PYZus{}graphs}\PY{p}{,} \PY{l+s+s2}{\PYZdq{}}\PY{l+s+s2}{COLLAB\PYZhy{}HOLDOUT}\PY{l+s+s2}{\PYZdq{}}\PY{p}{,} \PY{n}{num\PYZus{}graphs}\PY{o}{=}\PY{n+nb}{len}\PY{p}{(}\PY{n}{collab\PYZus{}test\PYZus{}graphs}\PY{p}{)}\PY{p}{)}

\PY{n}{gnn\PYZus{}imdb\PYZus{}df} \PY{o}{=} \PY{n}{pd}\PY{o}{.}\PY{n}{DataFrame}\PY{p}{(}\PY{n}{gnn\PYZus{}imdb\PYZus{}results}\PY{p}{)}
\PY{n}{gnn\PYZus{}collab\PYZus{}df} \PY{o}{=} \PY{n}{pd}\PY{o}{.}\PY{n}{DataFrame}\PY{p}{(}\PY{n}{gnn\PYZus{}collab\PYZus{}results}\PY{p}{)}
\PY{n}{gnn\PYZus{}imdb\PYZus{}summary} \PY{o}{=} \PY{n}{make\PYZus{}summary\PYZus{}dict}\PY{p}{(}\PY{n}{gnn\PYZus{}imdb\PYZus{}results}\PY{p}{)}
\PY{n}{gnn\PYZus{}collab\PYZus{}summary} \PY{o}{=} \PY{n}{make\PYZus{}summary\PYZus{}dict}\PY{p}{(}\PY{n}{gnn\PYZus{}collab\PYZus{}results}\PY{p}{)}
\PY{n}{rd\PYZus{}imdb\PYZus{}holdout\PYZus{}summary} \PY{o}{=} \PY{n}{make\PYZus{}summary\PYZus{}dict}\PY{p}{(}\PY{n}{rd\PYZus{}imdb\PYZus{}holdout}\PY{p}{)}
\PY{n}{rd\PYZus{}collab\PYZus{}holdout\PYZus{}summary} \PY{o}{=} \PY{n}{make\PYZus{}summary\PYZus{}dict}\PY{p}{(}\PY{n}{rd\PYZus{}collab\PYZus{}holdout}\PY{p}{)}

\PY{n}{display}\PY{p}{(}\PY{n}{gnn\PYZus{}imdb\PYZus{}df}\PY{o}{.}\PY{n}{head}\PY{p}{(}\PY{l+m+mi}{10}\PY{p}{)}\PY{p}{)}
\PY{n}{display}\PY{p}{(}\PY{n}{gnn\PYZus{}collab\PYZus{}df}\PY{o}{.}\PY{n}{head}\PY{p}{(}\PY{l+m+mi}{10}\PY{p}{)}\PY{p}{)}
\end{Verbatim}
\end{tcolorbox}

    
    \begin{Verbatim}[commandchars=\\\{\}]
   graph\_idx  num\_nodes  clique\_size  \textbackslash{}
0          0         12           12   
1          1         35           18   
2          2         13            8   
3          3         19            8   
4          4         19            6   
5          5         15            7   
6          6         12           12   
7          7         18            7   
8          8         12            8   
9          9         12            7   

                                              clique  valid    time  
0             [0, 1, 2, 3, 4, 5, 6, 7, 8, 9, 10, 11]   True  0.0111  
1  [29, 12, 19, 14, 3, 28, 16, 31, 27, 33, 0, 18,{\ldots}   True  0.0120  
2                         [2, 4, 3, 8, 10, 12, 7, 5]   True  0.0132  
3                        [1, 10, 9, 8, 13, 18, 2, 3]   True  0.0160  
4                               [11, 10, 9, 7, 5, 8]   True  0.0144  
5                            [1, 2, 3, 11, 13, 4, 9]   True  0.0113  
6             [0, 1, 2, 3, 4, 5, 6, 7, 8, 9, 10, 11]   True  0.0125  
7                           [11, 5, 8, 15, 13, 3, 4]   True  0.0170  
8                         [0, 11, 10, 5, 6, 8, 7, 4]   True  0.0175  
9                             [7, 9, 4, 0, 11, 3, 6]   True  0.0143  
    \end{Verbatim}

    
    
    \begin{Verbatim}[commandchars=\\\{\}]
   graph\_idx  num\_nodes  clique\_size  \textbackslash{}
0          0        157           24   
1          1         46           20   
2          2         60           23   
3          3         57           37   
4          4         37           19   
5          5         33            7   
6          6        189           30   
7          7         51           43   
8          8        133           41   
9          9         46           42   

                                              clique  valid    time  
0  [106, 32, 120, 126, 16, 136, 33, 141, 118, 40,{\ldots}   True  0.0110  
1  [29, 13, 2, 3, 28, 39, 36, 22, 23, 32, 31, 18,{\ldots}   True  0.0107  
2  [18, 13, 9, 15, 21, 25, 39, 5, 2, 44, 49, 40, {\ldots}   True  0.0106  
3  [19, 48, 32, 18, 39, 4, 0, 1, 45, 31, 52, 6, 5{\ldots}   True  0.0106  
4  [28, 1, 5, 2, 31, 15, 9, 4, 29, 13, 21, 24, 23{\ldots}   True  0.0104  
5                         [3, 0, 31, 22, 19, 23, 12]   True  0.0100  
6  [30, 142, 61, 157, 115, 119, 158, 120, 38, 132{\ldots}   True  0.0130  
7  [21, 34, 4, 46, 44, 43, 8, 7, 19, 41, 47, 40, {\ldots}   True  0.0120  
8  [82, 15, 118, 17, 111, 106, 22, 10, 50, 35, 34{\ldots}   True  0.0113  
9  [15, 1, 3, 2, 7, 6, 8, 4, 13, 9, 12, 37, 36, 2{\ldots}   True  0.0107  
    \end{Verbatim}

    
    \subsubsection{8.3 Comparación de los tres
enfoques}\label{comparaciuxf3n-de-los-tres-enfoques}

La comparación se organiza siguiendo los tres enfoques que se piden. El
baseline exacto ya se obtuvo en el ejercicio 1, RD puro se ha evaluado
en el ejercicio 3 y aquí se contrasta todo con el método híbrido. Para
que las conclusiones no dependan de texto interpretativo sino de
resultados visibles, en cada subapartado se muestra una tabla con los
cálculos que respaldan la comparación.

    \begin{itemize}
\item
  Fuerza bruta

  La referencia exacta queda limitada a los grafos pequeños de IMDB
  resueltos en el ejercicio 1. Esa restricción es importante, porque
  hace visible que el método exacto sirve como control de calidad, pero
  no como estrategia escalable.
\end{itemize}

    \begin{tcolorbox}[breakable, size=fbox, boxrule=1pt, pad at break*=1mm,colback=cellbackground, colframe=cellborder]
\prompt{In}{incolor}{287}{\boxspacing}
\begin{Verbatim}[commandchars=\\\{\}]
\PY{n}{bf\PYZus{}reference\PYZus{}df} \PY{o}{=} \PY{n}{pd}\PY{o}{.}\PY{n}{DataFrame}\PY{p}{(}\PY{p}{[}
    \PY{p}{\PYZob{}}
        \PY{l+s+s1}{\PYZsq{}}\PY{l+s+s1}{Dataset}\PY{l+s+s1}{\PYZsq{}}\PY{p}{:} \PY{l+s+s1}{\PYZsq{}}\PY{l+s+s1}{IMDB (grafos pequeños)}\PY{l+s+s1}{\PYZsq{}}\PY{p}{,}
        \PY{l+s+s1}{\PYZsq{}}\PY{l+s+s1}{Grafos evaluados}\PY{l+s+s1}{\PYZsq{}}\PY{p}{:} \PY{n}{make\PYZus{}summary\PYZus{}dict}\PY{p}{(}\PY{n}{bf\PYZus{}imdb\PYZus{}results}\PY{p}{)}\PY{p}{[}\PY{l+s+s1}{\PYZsq{}}\PY{l+s+s1}{num\PYZus{}graphs}\PY{l+s+s1}{\PYZsq{}}\PY{p}{]}\PY{p}{,}
        \PY{l+s+s1}{\PYZsq{}}\PY{l+s+s1}{Clique media}\PY{l+s+s1}{\PYZsq{}}\PY{p}{:} \PY{n}{make\PYZus{}summary\PYZus{}dict}\PY{p}{(}\PY{n}{bf\PYZus{}imdb\PYZus{}results}\PY{p}{)}\PY{p}{[}\PY{l+s+s1}{\PYZsq{}}\PY{l+s+s1}{mean\PYZus{}clique\PYZus{}size}\PY{l+s+s1}{\PYZsq{}}\PY{p}{]}\PY{p}{,}
        \PY{l+s+s1}{\PYZsq{}}\PY{l+s+s1}{Tiempo medio (s)}\PY{l+s+s1}{\PYZsq{}}\PY{p}{:} \PY{n}{make\PYZus{}summary\PYZus{}dict}\PY{p}{(}\PY{n}{bf\PYZus{}imdb\PYZus{}results}\PY{p}{)}\PY{p}{[}\PY{l+s+s1}{\PYZsq{}}\PY{l+s+s1}{mean\PYZus{}time}\PY{l+s+s1}{\PYZsq{}}\PY{p}{]}\PY{p}{,}
        \PY{l+s+s1}{\PYZsq{}}\PY{l+s+s1}{Validez (}\PY{l+s+s1}{\PYZpc{}}\PY{l+s+s1}{)}\PY{l+s+s1}{\PYZsq{}}\PY{p}{:} \PY{l+m+mi}{100} \PY{o}{*} \PY{n}{make\PYZus{}summary\PYZus{}dict}\PY{p}{(}\PY{n}{bf\PYZus{}imdb\PYZus{}results}\PY{p}{)}\PY{p}{[}\PY{l+s+s1}{\PYZsq{}}\PY{l+s+s1}{valid\PYZus{}rate}\PY{l+s+s1}{\PYZsq{}}\PY{p}{]}\PY{p}{,}
    \PY{p}{\PYZcb{}}\PY{p}{,}
    \PY{p}{\PYZob{}}
        \PY{l+s+s1}{\PYZsq{}}\PY{l+s+s1}{Dataset}\PY{l+s+s1}{\PYZsq{}}\PY{p}{:} \PY{l+s+s1}{\PYZsq{}}\PY{l+s+s1}{COLLAB (grafos pequeños)}\PY{l+s+s1}{\PYZsq{}}\PY{p}{,}
        \PY{l+s+s1}{\PYZsq{}}\PY{l+s+s1}{Grafos evaluados}\PY{l+s+s1}{\PYZsq{}}\PY{p}{:} \PY{n}{make\PYZus{}summary\PYZus{}dict}\PY{p}{(}\PY{n}{bf\PYZus{}collab\PYZus{}results}\PY{p}{)}\PY{p}{[}\PY{l+s+s1}{\PYZsq{}}\PY{l+s+s1}{num\PYZus{}graphs}\PY{l+s+s1}{\PYZsq{}}\PY{p}{]}\PY{p}{,}
        \PY{l+s+s1}{\PYZsq{}}\PY{l+s+s1}{Clique media}\PY{l+s+s1}{\PYZsq{}}\PY{p}{:} \PY{n}{make\PYZus{}summary\PYZus{}dict}\PY{p}{(}\PY{n}{bf\PYZus{}collab\PYZus{}results}\PY{p}{)}\PY{p}{[}\PY{l+s+s1}{\PYZsq{}}\PY{l+s+s1}{mean\PYZus{}clique\PYZus{}size}\PY{l+s+s1}{\PYZsq{}}\PY{p}{]}\PY{p}{,}
        \PY{l+s+s1}{\PYZsq{}}\PY{l+s+s1}{Tiempo medio (s)}\PY{l+s+s1}{\PYZsq{}}\PY{p}{:} \PY{n}{make\PYZus{}summary\PYZus{}dict}\PY{p}{(}\PY{n}{bf\PYZus{}collab\PYZus{}results}\PY{p}{)}\PY{p}{[}\PY{l+s+s1}{\PYZsq{}}\PY{l+s+s1}{mean\PYZus{}time}\PY{l+s+s1}{\PYZsq{}}\PY{p}{]}\PY{p}{,}
        \PY{l+s+s1}{\PYZsq{}}\PY{l+s+s1}{Validez (}\PY{l+s+s1}{\PYZpc{}}\PY{l+s+s1}{)}\PY{l+s+s1}{\PYZsq{}}\PY{p}{:} \PY{l+m+mi}{100} \PY{o}{*} \PY{n}{make\PYZus{}summary\PYZus{}dict}\PY{p}{(}\PY{n}{bf\PYZus{}collab\PYZus{}results}\PY{p}{)}\PY{p}{[}\PY{l+s+s1}{\PYZsq{}}\PY{l+s+s1}{valid\PYZus{}rate}\PY{l+s+s1}{\PYZsq{}}\PY{p}{]}\PY{p}{,}
    \PY{p}{\PYZcb{}}\PY{p}{,}
\PY{p}{]}\PY{p}{)}

\PY{n}{display}\PY{p}{(}\PY{n}{bf\PYZus{}reference\PYZus{}df}\PY{o}{.}\PY{n}{round}\PY{p}{(}\PY{l+m+mi}{4}\PY{p}{)}\PY{p}{)}
\end{Verbatim}
\end{tcolorbox}

    
    \begin{Verbatim}[commandchars=\\\{\}]
                    Dataset  Grafos evaluados  Clique media  Tiempo medio (s)  \textbackslash{}
0    IMDB (grafos pequeños)                15       11.7333            0.0758   
1  COLLAB (grafos pequeños)                 0           NaN               NaN   

   Validez (\%)  
0        100.0  
1          NaN  
    \end{Verbatim}

    
    \begin{itemize}
\item
  RD puro desde el baricentro

  La referencia aproximada sin aprendizaje corresponde a RD puro,
  partiendo del baricentro uniforme, cuyos resultados ya se han
  calculado en el ejercicio 3 y en la evaluación hold-out incluida en
  este ejercicio.
\end{itemize}

    \begin{tcolorbox}[breakable, size=fbox, boxrule=1pt, pad at break*=1mm,colback=cellbackground, colframe=cellborder]
\prompt{In}{incolor}{288}{\boxspacing}
\begin{Verbatim}[commandchars=\\\{\}]
\PY{n}{rd\PYZus{}holdout\PYZus{}df} \PY{o}{=} \PY{n}{pd}\PY{o}{.}\PY{n}{DataFrame}\PY{p}{(}\PY{p}{[}
    \PY{p}{\PYZob{}}
        \PY{l+s+s1}{\PYZsq{}}\PY{l+s+s1}{Dataset}\PY{l+s+s1}{\PYZsq{}}\PY{p}{:} \PY{l+s+s1}{\PYZsq{}}\PY{l+s+s1}{IMDB}\PY{l+s+s1}{\PYZsq{}}\PY{p}{,}
        \PY{l+s+s1}{\PYZsq{}}\PY{l+s+s1}{Grafos hold\PYZhy{}out}\PY{l+s+s1}{\PYZsq{}}\PY{p}{:} \PY{n}{rd\PYZus{}imdb\PYZus{}holdout\PYZus{}summary}\PY{p}{[}\PY{l+s+s1}{\PYZsq{}}\PY{l+s+s1}{num\PYZus{}graphs}\PY{l+s+s1}{\PYZsq{}}\PY{p}{]}\PY{p}{,}
        \PY{l+s+s1}{\PYZsq{}}\PY{l+s+s1}{Clique media}\PY{l+s+s1}{\PYZsq{}}\PY{p}{:} \PY{n}{rd\PYZus{}imdb\PYZus{}holdout\PYZus{}summary}\PY{p}{[}\PY{l+s+s1}{\PYZsq{}}\PY{l+s+s1}{mean\PYZus{}clique\PYZus{}size}\PY{l+s+s1}{\PYZsq{}}\PY{p}{]}\PY{p}{,}
        \PY{l+s+s1}{\PYZsq{}}\PY{l+s+s1}{Tiempo medio (s)}\PY{l+s+s1}{\PYZsq{}}\PY{p}{:} \PY{n}{rd\PYZus{}imdb\PYZus{}holdout\PYZus{}summary}\PY{p}{[}\PY{l+s+s1}{\PYZsq{}}\PY{l+s+s1}{mean\PYZus{}time}\PY{l+s+s1}{\PYZsq{}}\PY{p}{]}\PY{p}{,}
        \PY{l+s+s1}{\PYZsq{}}\PY{l+s+s1}{Validez (}\PY{l+s+s1}{\PYZpc{}}\PY{l+s+s1}{)}\PY{l+s+s1}{\PYZsq{}}\PY{p}{:} \PY{l+m+mi}{100} \PY{o}{*} \PY{n}{rd\PYZus{}imdb\PYZus{}holdout\PYZus{}summary}\PY{p}{[}\PY{l+s+s1}{\PYZsq{}}\PY{l+s+s1}{valid\PYZus{}rate}\PY{l+s+s1}{\PYZsq{}}\PY{p}{]}\PY{p}{,}
    \PY{p}{\PYZcb{}}\PY{p}{,}
    \PY{p}{\PYZob{}}
        \PY{l+s+s1}{\PYZsq{}}\PY{l+s+s1}{Dataset}\PY{l+s+s1}{\PYZsq{}}\PY{p}{:} \PY{l+s+s1}{\PYZsq{}}\PY{l+s+s1}{COLLAB}\PY{l+s+s1}{\PYZsq{}}\PY{p}{,}
        \PY{l+s+s1}{\PYZsq{}}\PY{l+s+s1}{Grafos hold\PYZhy{}out}\PY{l+s+s1}{\PYZsq{}}\PY{p}{:} \PY{n}{rd\PYZus{}collab\PYZus{}holdout\PYZus{}summary}\PY{p}{[}\PY{l+s+s1}{\PYZsq{}}\PY{l+s+s1}{num\PYZus{}graphs}\PY{l+s+s1}{\PYZsq{}}\PY{p}{]}\PY{p}{,}
        \PY{l+s+s1}{\PYZsq{}}\PY{l+s+s1}{Clique media}\PY{l+s+s1}{\PYZsq{}}\PY{p}{:} \PY{n}{rd\PYZus{}collab\PYZus{}holdout\PYZus{}summary}\PY{p}{[}\PY{l+s+s1}{\PYZsq{}}\PY{l+s+s1}{mean\PYZus{}clique\PYZus{}size}\PY{l+s+s1}{\PYZsq{}}\PY{p}{]}\PY{p}{,}
        \PY{l+s+s1}{\PYZsq{}}\PY{l+s+s1}{Tiempo medio (s)}\PY{l+s+s1}{\PYZsq{}}\PY{p}{:} \PY{n}{rd\PYZus{}collab\PYZus{}holdout\PYZus{}summary}\PY{p}{[}\PY{l+s+s1}{\PYZsq{}}\PY{l+s+s1}{mean\PYZus{}time}\PY{l+s+s1}{\PYZsq{}}\PY{p}{]}\PY{p}{,}
        \PY{l+s+s1}{\PYZsq{}}\PY{l+s+s1}{Validez (}\PY{l+s+s1}{\PYZpc{}}\PY{l+s+s1}{)}\PY{l+s+s1}{\PYZsq{}}\PY{p}{:} \PY{l+m+mi}{100} \PY{o}{*} \PY{n}{rd\PYZus{}collab\PYZus{}holdout\PYZus{}summary}\PY{p}{[}\PY{l+s+s1}{\PYZsq{}}\PY{l+s+s1}{valid\PYZus{}rate}\PY{l+s+s1}{\PYZsq{}}\PY{p}{]}\PY{p}{,}
    \PY{p}{\PYZcb{}}\PY{p}{,}
\PY{p}{]}\PY{p}{)}

\PY{n}{display}\PY{p}{(}\PY{n}{rd\PYZus{}holdout\PYZus{}df}\PY{o}{.}\PY{n}{round}\PY{p}{(}\PY{l+m+mi}{4}\PY{p}{)}\PY{p}{)}
\end{Verbatim}
\end{tcolorbox}

    
    \begin{Verbatim}[commandchars=\\\{\}]
  Dataset  Grafos hold-out  Clique media  Tiempo medio (s)  Validez (\%)
0    IMDB               48        9.4375            0.0029        100.0
1  COLLAB               48       27.8958            0.0052        100.0
    \end{Verbatim}

    
    \begin{itemize}
\item
  GNN + RD

  El método híbrido usa la GNN como inicializador y RD como refinador.
  La comparación relevante consiste en ver si esa inicialización
  aprendida produce cliques mayores que las obtenidas por RD puro y
  cuánto cuesta en tiempo adicional.
\end{itemize}

    \begin{tcolorbox}[breakable, size=fbox, boxrule=1pt, pad at break*=1mm,colback=cellbackground, colframe=cellborder]
\prompt{In}{incolor}{289}{\boxspacing}
\begin{Verbatim}[commandchars=\\\{\}]
\PY{n}{gnn\PYZus{}holdout\PYZus{}df} \PY{o}{=} \PY{n}{pd}\PY{o}{.}\PY{n}{DataFrame}\PY{p}{(}\PY{p}{[}
    \PY{p}{\PYZob{}}
        \PY{l+s+s1}{\PYZsq{}}\PY{l+s+s1}{Dataset}\PY{l+s+s1}{\PYZsq{}}\PY{p}{:} \PY{l+s+s1}{\PYZsq{}}\PY{l+s+s1}{IMDB}\PY{l+s+s1}{\PYZsq{}}\PY{p}{,}
        \PY{l+s+s1}{\PYZsq{}}\PY{l+s+s1}{Grafos hold\PYZhy{}out}\PY{l+s+s1}{\PYZsq{}}\PY{p}{:} \PY{n}{gnn\PYZus{}imdb\PYZus{}summary}\PY{p}{[}\PY{l+s+s1}{\PYZsq{}}\PY{l+s+s1}{num\PYZus{}graphs}\PY{l+s+s1}{\PYZsq{}}\PY{p}{]}\PY{p}{,}
        \PY{l+s+s1}{\PYZsq{}}\PY{l+s+s1}{Clique media}\PY{l+s+s1}{\PYZsq{}}\PY{p}{:} \PY{n}{gnn\PYZus{}imdb\PYZus{}summary}\PY{p}{[}\PY{l+s+s1}{\PYZsq{}}\PY{l+s+s1}{mean\PYZus{}clique\PYZus{}size}\PY{l+s+s1}{\PYZsq{}}\PY{p}{]}\PY{p}{,}
        \PY{l+s+s1}{\PYZsq{}}\PY{l+s+s1}{Tiempo medio (s)}\PY{l+s+s1}{\PYZsq{}}\PY{p}{:} \PY{n}{gnn\PYZus{}imdb\PYZus{}summary}\PY{p}{[}\PY{l+s+s1}{\PYZsq{}}\PY{l+s+s1}{mean\PYZus{}time}\PY{l+s+s1}{\PYZsq{}}\PY{p}{]}\PY{p}{,}
        \PY{l+s+s1}{\PYZsq{}}\PY{l+s+s1}{Validez (}\PY{l+s+s1}{\PYZpc{}}\PY{l+s+s1}{)}\PY{l+s+s1}{\PYZsq{}}\PY{p}{:} \PY{l+m+mi}{100} \PY{o}{*} \PY{n}{gnn\PYZus{}imdb\PYZus{}summary}\PY{p}{[}\PY{l+s+s1}{\PYZsq{}}\PY{l+s+s1}{valid\PYZus{}rate}\PY{l+s+s1}{\PYZsq{}}\PY{p}{]}\PY{p}{,}
    \PY{p}{\PYZcb{}}\PY{p}{,}
    \PY{p}{\PYZob{}}
        \PY{l+s+s1}{\PYZsq{}}\PY{l+s+s1}{Dataset}\PY{l+s+s1}{\PYZsq{}}\PY{p}{:} \PY{l+s+s1}{\PYZsq{}}\PY{l+s+s1}{COLLAB}\PY{l+s+s1}{\PYZsq{}}\PY{p}{,}
        \PY{l+s+s1}{\PYZsq{}}\PY{l+s+s1}{Grafos hold\PYZhy{}out}\PY{l+s+s1}{\PYZsq{}}\PY{p}{:} \PY{n}{gnn\PYZus{}collab\PYZus{}summary}\PY{p}{[}\PY{l+s+s1}{\PYZsq{}}\PY{l+s+s1}{num\PYZus{}graphs}\PY{l+s+s1}{\PYZsq{}}\PY{p}{]}\PY{p}{,}
        \PY{l+s+s1}{\PYZsq{}}\PY{l+s+s1}{Clique media}\PY{l+s+s1}{\PYZsq{}}\PY{p}{:} \PY{n}{gnn\PYZus{}collab\PYZus{}summary}\PY{p}{[}\PY{l+s+s1}{\PYZsq{}}\PY{l+s+s1}{mean\PYZus{}clique\PYZus{}size}\PY{l+s+s1}{\PYZsq{}}\PY{p}{]}\PY{p}{,}
        \PY{l+s+s1}{\PYZsq{}}\PY{l+s+s1}{Tiempo medio (s)}\PY{l+s+s1}{\PYZsq{}}\PY{p}{:} \PY{n}{gnn\PYZus{}collab\PYZus{}summary}\PY{p}{[}\PY{l+s+s1}{\PYZsq{}}\PY{l+s+s1}{mean\PYZus{}time}\PY{l+s+s1}{\PYZsq{}}\PY{p}{]}\PY{p}{,}
        \PY{l+s+s1}{\PYZsq{}}\PY{l+s+s1}{Validez (}\PY{l+s+s1}{\PYZpc{}}\PY{l+s+s1}{)}\PY{l+s+s1}{\PYZsq{}}\PY{p}{:} \PY{l+m+mi}{100} \PY{o}{*} \PY{n}{gnn\PYZus{}collab\PYZus{}summary}\PY{p}{[}\PY{l+s+s1}{\PYZsq{}}\PY{l+s+s1}{valid\PYZus{}rate}\PY{l+s+s1}{\PYZsq{}}\PY{p}{]}\PY{p}{,}
    \PY{p}{\PYZcb{}}\PY{p}{,}
\PY{p}{]}\PY{p}{)}

\PY{n}{holdout\PYZus{}comparison\PYZus{}df} \PY{o}{=} \PY{n}{pd}\PY{o}{.}\PY{n}{DataFrame}\PY{p}{(}\PY{p}{[}
    \PY{p}{\PYZob{}}
        \PY{l+s+s1}{\PYZsq{}}\PY{l+s+s1}{Dataset}\PY{l+s+s1}{\PYZsq{}}\PY{p}{:} \PY{l+s+s1}{\PYZsq{}}\PY{l+s+s1}{IMDB}\PY{l+s+s1}{\PYZsq{}}\PY{p}{,}
        \PY{l+s+s1}{\PYZsq{}}\PY{l+s+s1}{RD puro clique media}\PY{l+s+s1}{\PYZsq{}}\PY{p}{:} \PY{n}{rd\PYZus{}imdb\PYZus{}holdout\PYZus{}summary}\PY{p}{[}\PY{l+s+s1}{\PYZsq{}}\PY{l+s+s1}{mean\PYZus{}clique\PYZus{}size}\PY{l+s+s1}{\PYZsq{}}\PY{p}{]}\PY{p}{,}
        \PY{l+s+s1}{\PYZsq{}}\PY{l+s+s1}{GNN + RD clique media}\PY{l+s+s1}{\PYZsq{}}\PY{p}{:} \PY{n}{gnn\PYZus{}imdb\PYZus{}summary}\PY{p}{[}\PY{l+s+s1}{\PYZsq{}}\PY{l+s+s1}{mean\PYZus{}clique\PYZus{}size}\PY{l+s+s1}{\PYZsq{}}\PY{p}{]}\PY{p}{,}
        \PY{l+s+s1}{\PYZsq{}}\PY{l+s+s1}{Mejora clique media}\PY{l+s+s1}{\PYZsq{}}\PY{p}{:} \PY{n}{gnn\PYZus{}imdb\PYZus{}summary}\PY{p}{[}\PY{l+s+s1}{\PYZsq{}}\PY{l+s+s1}{mean\PYZus{}clique\PYZus{}size}\PY{l+s+s1}{\PYZsq{}}\PY{p}{]} \PY{o}{\PYZhy{}} \PY{n}{rd\PYZus{}imdb\PYZus{}holdout\PYZus{}summary}\PY{p}{[}\PY{l+s+s1}{\PYZsq{}}\PY{l+s+s1}{mean\PYZus{}clique\PYZus{}size}\PY{l+s+s1}{\PYZsq{}}\PY{p}{]}\PY{p}{,}
        \PY{l+s+s1}{\PYZsq{}}\PY{l+s+s1}{RD puro tiempo (s)}\PY{l+s+s1}{\PYZsq{}}\PY{p}{:} \PY{n}{rd\PYZus{}imdb\PYZus{}holdout\PYZus{}summary}\PY{p}{[}\PY{l+s+s1}{\PYZsq{}}\PY{l+s+s1}{mean\PYZus{}time}\PY{l+s+s1}{\PYZsq{}}\PY{p}{]}\PY{p}{,}
        \PY{l+s+s1}{\PYZsq{}}\PY{l+s+s1}{GNN + RD tiempo (s)}\PY{l+s+s1}{\PYZsq{}}\PY{p}{:} \PY{n}{gnn\PYZus{}imdb\PYZus{}summary}\PY{p}{[}\PY{l+s+s1}{\PYZsq{}}\PY{l+s+s1}{mean\PYZus{}time}\PY{l+s+s1}{\PYZsq{}}\PY{p}{]}\PY{p}{,}
        \PY{l+s+s1}{\PYZsq{}}\PY{l+s+s1}{Sobrecoste tiempo (s)}\PY{l+s+s1}{\PYZsq{}}\PY{p}{:} \PY{n}{gnn\PYZus{}imdb\PYZus{}summary}\PY{p}{[}\PY{l+s+s1}{\PYZsq{}}\PY{l+s+s1}{mean\PYZus{}time}\PY{l+s+s1}{\PYZsq{}}\PY{p}{]} \PY{o}{\PYZhy{}} \PY{n}{rd\PYZus{}imdb\PYZus{}holdout\PYZus{}summary}\PY{p}{[}\PY{l+s+s1}{\PYZsq{}}\PY{l+s+s1}{mean\PYZus{}time}\PY{l+s+s1}{\PYZsq{}}\PY{p}{]}\PY{p}{,}
        \PY{l+s+s1}{\PYZsq{}}\PY{l+s+s1}{Validez RD (}\PY{l+s+s1}{\PYZpc{}}\PY{l+s+s1}{)}\PY{l+s+s1}{\PYZsq{}}\PY{p}{:} \PY{l+m+mi}{100} \PY{o}{*} \PY{n}{rd\PYZus{}imdb\PYZus{}holdout\PYZus{}summary}\PY{p}{[}\PY{l+s+s1}{\PYZsq{}}\PY{l+s+s1}{valid\PYZus{}rate}\PY{l+s+s1}{\PYZsq{}}\PY{p}{]}\PY{p}{,}
        \PY{l+s+s1}{\PYZsq{}}\PY{l+s+s1}{Validez GNN + RD (}\PY{l+s+s1}{\PYZpc{}}\PY{l+s+s1}{)}\PY{l+s+s1}{\PYZsq{}}\PY{p}{:} \PY{l+m+mi}{100} \PY{o}{*} \PY{n}{gnn\PYZus{}imdb\PYZus{}summary}\PY{p}{[}\PY{l+s+s1}{\PYZsq{}}\PY{l+s+s1}{valid\PYZus{}rate}\PY{l+s+s1}{\PYZsq{}}\PY{p}{]}\PY{p}{,}
    \PY{p}{\PYZcb{}}\PY{p}{,}
    \PY{p}{\PYZob{}}
        \PY{l+s+s1}{\PYZsq{}}\PY{l+s+s1}{Dataset}\PY{l+s+s1}{\PYZsq{}}\PY{p}{:} \PY{l+s+s1}{\PYZsq{}}\PY{l+s+s1}{COLLAB}\PY{l+s+s1}{\PYZsq{}}\PY{p}{,}
        \PY{l+s+s1}{\PYZsq{}}\PY{l+s+s1}{RD puro clique media}\PY{l+s+s1}{\PYZsq{}}\PY{p}{:} \PY{n}{rd\PYZus{}collab\PYZus{}holdout\PYZus{}summary}\PY{p}{[}\PY{l+s+s1}{\PYZsq{}}\PY{l+s+s1}{mean\PYZus{}clique\PYZus{}size}\PY{l+s+s1}{\PYZsq{}}\PY{p}{]}\PY{p}{,}
        \PY{l+s+s1}{\PYZsq{}}\PY{l+s+s1}{GNN + RD clique media}\PY{l+s+s1}{\PYZsq{}}\PY{p}{:} \PY{n}{gnn\PYZus{}collab\PYZus{}summary}\PY{p}{[}\PY{l+s+s1}{\PYZsq{}}\PY{l+s+s1}{mean\PYZus{}clique\PYZus{}size}\PY{l+s+s1}{\PYZsq{}}\PY{p}{]}\PY{p}{,}
        \PY{l+s+s1}{\PYZsq{}}\PY{l+s+s1}{Mejora clique media}\PY{l+s+s1}{\PYZsq{}}\PY{p}{:} \PY{n}{gnn\PYZus{}collab\PYZus{}summary}\PY{p}{[}\PY{l+s+s1}{\PYZsq{}}\PY{l+s+s1}{mean\PYZus{}clique\PYZus{}size}\PY{l+s+s1}{\PYZsq{}}\PY{p}{]} \PY{o}{\PYZhy{}} \PY{n}{rd\PYZus{}collab\PYZus{}holdout\PYZus{}summary}\PY{p}{[}\PY{l+s+s1}{\PYZsq{}}\PY{l+s+s1}{mean\PYZus{}clique\PYZus{}size}\PY{l+s+s1}{\PYZsq{}}\PY{p}{]}\PY{p}{,}
        \PY{l+s+s1}{\PYZsq{}}\PY{l+s+s1}{RD puro tiempo (s)}\PY{l+s+s1}{\PYZsq{}}\PY{p}{:} \PY{n}{rd\PYZus{}collab\PYZus{}holdout\PYZus{}summary}\PY{p}{[}\PY{l+s+s1}{\PYZsq{}}\PY{l+s+s1}{mean\PYZus{}time}\PY{l+s+s1}{\PYZsq{}}\PY{p}{]}\PY{p}{,}
        \PY{l+s+s1}{\PYZsq{}}\PY{l+s+s1}{GNN + RD tiempo (s)}\PY{l+s+s1}{\PYZsq{}}\PY{p}{:} \PY{n}{gnn\PYZus{}collab\PYZus{}summary}\PY{p}{[}\PY{l+s+s1}{\PYZsq{}}\PY{l+s+s1}{mean\PYZus{}time}\PY{l+s+s1}{\PYZsq{}}\PY{p}{]}\PY{p}{,}
        \PY{l+s+s1}{\PYZsq{}}\PY{l+s+s1}{Sobrecoste tiempo (s)}\PY{l+s+s1}{\PYZsq{}}\PY{p}{:} \PY{n}{gnn\PYZus{}collab\PYZus{}summary}\PY{p}{[}\PY{l+s+s1}{\PYZsq{}}\PY{l+s+s1}{mean\PYZus{}time}\PY{l+s+s1}{\PYZsq{}}\PY{p}{]} \PY{o}{\PYZhy{}} \PY{n}{rd\PYZus{}collab\PYZus{}holdout\PYZus{}summary}\PY{p}{[}\PY{l+s+s1}{\PYZsq{}}\PY{l+s+s1}{mean\PYZus{}time}\PY{l+s+s1}{\PYZsq{}}\PY{p}{]}\PY{p}{,}
        \PY{l+s+s1}{\PYZsq{}}\PY{l+s+s1}{Validez RD (}\PY{l+s+s1}{\PYZpc{}}\PY{l+s+s1}{)}\PY{l+s+s1}{\PYZsq{}}\PY{p}{:} \PY{l+m+mi}{100} \PY{o}{*} \PY{n}{rd\PYZus{}collab\PYZus{}holdout\PYZus{}summary}\PY{p}{[}\PY{l+s+s1}{\PYZsq{}}\PY{l+s+s1}{valid\PYZus{}rate}\PY{l+s+s1}{\PYZsq{}}\PY{p}{]}\PY{p}{,}
        \PY{l+s+s1}{\PYZsq{}}\PY{l+s+s1}{Validez GNN + RD (}\PY{l+s+s1}{\PYZpc{}}\PY{l+s+s1}{)}\PY{l+s+s1}{\PYZsq{}}\PY{p}{:} \PY{l+m+mi}{100} \PY{o}{*} \PY{n}{gnn\PYZus{}collab\PYZus{}summary}\PY{p}{[}\PY{l+s+s1}{\PYZsq{}}\PY{l+s+s1}{valid\PYZus{}rate}\PY{l+s+s1}{\PYZsq{}}\PY{p}{]}\PY{p}{,}
    \PY{p}{\PYZcb{}}\PY{p}{,}
\PY{p}{]}\PY{p}{)}

\PY{n}{display}\PY{p}{(}\PY{n}{gnn\PYZus{}holdout\PYZus{}df}\PY{o}{.}\PY{n}{round}\PY{p}{(}\PY{l+m+mi}{4}\PY{p}{)}\PY{p}{)}
\PY{n}{display}\PY{p}{(}\PY{n}{holdout\PYZus{}comparison\PYZus{}df}\PY{o}{.}\PY{n}{round}\PY{p}{(}\PY{l+m+mi}{4}\PY{p}{)}\PY{p}{)}
\end{Verbatim}
\end{tcolorbox}

    
    \begin{Verbatim}[commandchars=\\\{\}]
  Dataset  Grafos hold-out  Clique media  Tiempo medio (s)  Validez (\%)
0    IMDB               48        9.5208            0.0115        100.0
1  COLLAB               48       28.0417            0.0115        100.0
    \end{Verbatim}

    
    
    \begin{Verbatim}[commandchars=\\\{\}]
  Dataset  RD puro clique media  GNN + RD clique media  Mejora clique media  \textbackslash{}
0    IMDB                9.4375                 9.5208               0.0833   
1  COLLAB               27.8958                28.0417               0.1458   

   RD puro tiempo (s)  GNN + RD tiempo (s)  Sobrecoste tiempo (s)  \textbackslash{}
0              0.0029               0.0115                 0.0086   
1              0.0052               0.0115                 0.0063   

   Validez RD (\%)  Validez GNN + RD (\%)  
0           100.0                 100.0  
1           100.0                 100.0  
    \end{Verbatim}

    
    \subsubsection{8.4 Ablación: variar el número de iteraciones de RD
(train: 0,1,3,5; test:
10,50,100)}\label{ablaciuxf3n-variar-el-nuxfamero-de-iteraciones-de-rd-train-0135-test-1050100}

    \begin{tcolorbox}[breakable, size=fbox, boxrule=1pt, pad at break*=1mm,colback=cellbackground, colframe=cellborder]
\prompt{In}{incolor}{290}{\boxspacing}
\begin{Verbatim}[commandchars=\\\{\}]
\PY{n}{imdb\PYZus{}ablation\PYZus{}train}\PY{p}{,} \PY{n}{imdb\PYZus{}ablation\PYZus{}test} \PY{o}{=} \PY{n}{split\PYZus{}dataset}\PY{p}{(}\PY{n}{imdb\PYZus{}dataset}\PY{p}{,} \PY{n}{train\PYZus{}ratio}\PY{o}{=}\PY{l+m+mf}{0.8}\PY{p}{,} \PY{n}{max\PYZus{}graphs}\PY{o}{=}\PY{l+m+mi}{120}\PY{p}{)}
\PY{n}{imdb\PYZus{}ablation\PYZus{}df} \PY{o}{=} \PY{n}{run\PYZus{}gnn\PYZus{}rd\PYZus{}ablation}\PY{p}{(}
    \PY{n}{imdb\PYZus{}ablation\PYZus{}train}\PY{p}{,}
    \PY{n}{imdb\PYZus{}ablation\PYZus{}test}\PY{p}{,}
    \PY{n}{train\PYZus{}iters\PYZus{}list}\PY{o}{=}\PY{p}{(}\PY{l+m+mi}{0}\PY{p}{,} \PY{l+m+mi}{1}\PY{p}{,} \PY{l+m+mi}{3}\PY{p}{,} \PY{l+m+mi}{5}\PY{p}{)}\PY{p}{,}
    \PY{n}{test\PYZus{}iters\PYZus{}list}\PY{o}{=}\PY{p}{(}\PY{l+m+mi}{10}\PY{p}{,} \PY{l+m+mi}{50}\PY{p}{,} \PY{l+m+mi}{100}\PY{p}{)}\PY{p}{,}
    \PY{n}{num\PYZus{}epochs}\PY{o}{=}\PY{l+m+mi}{6}\PY{p}{,}
    \PY{n}{lr}\PY{o}{=}\PY{l+m+mf}{0.01}\PY{p}{,}
    \PY{n}{hidden\PYZus{}dim}\PY{o}{=}\PY{l+m+mi}{32}\PY{p}{,}
\PY{p}{)}
\PY{n}{display}\PY{p}{(}\PY{n}{imdb\PYZus{}ablation\PYZus{}df}\PY{p}{)}
\PY{n}{plot\PYZus{}ablation\PYZus{}results}\PY{p}{(}\PY{n}{imdb\PYZus{}ablation\PYZus{}df}\PY{p}{,} \PY{l+s+s2}{\PYZdq{}}\PY{l+s+s2}{Ablación RD sobre IMDB}\PY{l+s+s2}{\PYZdq{}}\PY{p}{)}
\end{Verbatim}
\end{tcolorbox}

    \begin{Verbatim}[commandchars=\\\{\}]
Época 1/6 - loss media: -0.6948
Época 2/6 - loss media: -0.7482
Época 3/6 - loss media: -0.7675
Época 4/6 - loss media: -0.7773
Época 5/6 - loss media: -0.7739
Época 6/6 - loss media: -0.7812
Época 1/6 - loss media: -0.7313
Época 2/6 - loss media: -0.7832
Época 3/6 - loss media: -0.8086
Época 4/6 - loss media: -0.8133
Época 5/6 - loss media: -0.8253
Época 6/6 - loss media: -0.8297
Época 1/6 - loss media: -0.8264
Época 2/6 - loss media: -0.8409
Época 3/6 - loss media: -0.8503
Época 4/6 - loss media: -0.8524
Época 5/6 - loss media: -0.8535
Época 6/6 - loss media: -0.8540
Época 1/6 - loss media: -0.8380
Época 2/6 - loss media: -0.8566
Época 3/6 - loss media: -0.8610
Época 4/6 - loss media: -0.8597
Época 5/6 - loss media: -0.8641
Época 6/6 - loss media: -0.8619
    \end{Verbatim}

    
    \begin{Verbatim}[commandchars=\\\{\}]
    train\_rd\_iters  test\_rd\_iters  mean\_clique\_size  mean\_time  valid\_rate
0                0             10           10.8333     0.0027         1.0
1                0             50           10.8750     0.0063         1.0
2                0            100           10.8750     0.0107         1.0
3                1             10           10.8750     0.0027         1.0
4                1             50           10.8750     0.0065         1.0
5                1            100           10.8750     0.0107         1.0
6                3             10           10.8333     0.0026         1.0
7                3             50           10.8750     0.0067         1.0
8                3            100           10.8750     0.0126         1.0
9                5             10           10.8333     0.0028         1.0
10               5             50           10.8750     0.0063         1.0
11               5            100           10.8750     0.0124         1.0
    \end{Verbatim}

    
    \begin{center}
    \adjustimage{max size={0.9\linewidth}{0.9\paperheight}}{practica_01_jordi_blasco_lozano_files/practica_01_jordi_blasco_lozano_103_2.png}
    \end{center}
    { \hspace*{\fill} \\}
    
    \subsubsection{8.5 Informe del ejercicio}\label{informe-del-ejercicio}

La comparación de 8.3 deja clara la diferencia entre los tres enfoques.
La referencia exacta por fuerza bruta solo ha podido sostenerse en IMDB
sobre \texttt{15} grafos pequeños, donde ha alcanzado una clique media
de \texttt{11.7333} con validez del \texttt{100\%}, mientras que en
COLLAB no ha producido resultados utilizables bajo la restricción de
tamaño fijada. Eso confirma que la fuerza bruta sigue siendo útil como
control de calidad, pero no como método operativo cuando el número de
nodos crece. Su papel aquí es el de techo combinatorio en instancias
pequeñas, no el de solución escalable.

El resumen de entrenamiento de la celda posterior al bloque de épocas
muestra que en IMDB la loss ha pasado de \texttt{-0.8268} a
\texttt{-0.8566}, con un descenso neto de \texttt{0.0298}, mientras que
en COLLAB ha pasado de \texttt{-0.8954} a \texttt{-0.9237}, con un
descenso neto de \texttt{0.0284}. Como la loss está definida como el
negativo del objetivo que se quiere maximizar, que se haga más negativa
significa que el modelo está aprendiendo inicializaciones que favorecen
soluciones de mayor calidad tras el refinamiento con RD. También se
observa que IMDB desciende ligeramente más que COLLAB, lo cual encaja
con la idea de que ambos modelos están aprendiendo una señal útil,
aunque con dinámicas algo distintas según la estructura del dataset.

En los conjuntos hold-out, RD puro desde el baricentro ha obtenido una
clique media de \texttt{9.4375} en IMDB y \texttt{27.8958} en COLLAB,
con tiempos medios de \texttt{0.0029\ s} y \texttt{0.0052\ s},
respectivamente, y con validez del \texttt{100\%} en ambos casos. Cuando
se compara con el método híbrido en la tabla final de 8.3, se observa
que \textbf{GNN + RD} ha subido a \texttt{9.5208} en IMDB y a
\texttt{28.0417} en COLLAB. La mejora ha sido de \texttt{+0.0833} nodos
de media en IMDB y de \texttt{+0.1458} en COLLAB. Es una ganancia
pequeña, pero consistente y positiva, lo cual sugiere que la GNN no está
resolviendo el problema por sí sola, sino sesgando la distribución
inicial hacia regiones del simplex que RD puede explotar mejor que
cuando parte del baricentro uniforme. Dicho de otra forma, el
aprendizaje no sustituye al refinamiento dinámico, pero sí parece
reducir ligeramente la ambigüedad inicial de la búsqueda.

La contrapartida de esa mejora ha estado en el tiempo. En IMDB, el coste
medio ha pasado de \texttt{0.0029\ s} con RD puro a \texttt{0.0115\ s}
con GNN + RD, es decir, un sobrecoste de \texttt{0.0086\ s} y un tiempo
unas \texttt{4.0} veces mayor. En COLLAB, el tiempo medio ha pasado de
\texttt{0.0052\ s} a \texttt{0.0115\ s}, con un sobrecoste de
\texttt{0.0063\ s}, aproximadamente \texttt{2.2} veces más. Esto tiene
una interpretación directa: la red aporta información útil, pero esa
información no es gratis. Hay que calcular features, propagar mensajes y
producir la distribución inicial antes de aplicar RD. Por eso el método
híbrido solo compensa si esa inicialización aprendida aporta una mejora
suficientemente valiosa en calidad o estabilidad.

La ablación de 8.4 ayuda a entender mejor dónde aparece realmente ese
beneficio. La mejor media observada ha sido \texttt{10.8750}, alcanzada
en la mayoría de configuraciones, mientras que los únicos casos que
bajan a \texttt{10.8333} son \texttt{train\ RD\ =\ 0,\ test\ RD\ =\ 10},
\texttt{train\ RD\ =\ 3,\ test\ RD\ =\ 10} y
\texttt{train\ RD\ =\ 5,\ test\ RD\ =\ 10}. Además, al pasar de
\texttt{test\ RD\ =\ 10} a \texttt{50} el tiempo sube aproximadamente de
\texttt{0.0026-0.0028\ s} a \texttt{0.0063-0.0067\ s}, y con
\texttt{100} iteraciones llega a \texttt{0.0107-0.0126\ s}, mientras que
la calidad apenas cambia una vez alcanzado el nivel de \texttt{10.8750}.
Esto indica que, en este subconjunto, muchas iteraciones extra de RD en
inferencia ya no compran una mejora real, porque la solución se ha
estabilizado antes. También sugiere que aumentar RD en train no
garantiza una mejora adicional por sí mismo: el beneficio parece
depender más de usar una inicialización razonable y un refinamiento
suficiente en test que de forzar un número alto de iteraciones durante
todo el proceso.

    \subsection{9. Ejercicio 5: Comparación
final}\label{ejercicio-5-comparaciuxf3n-final}

    \subsubsection{9.1 Tabla resumen final}\label{tabla-resumen-final}

La tabla siguiente se ajusta exactamente al formato pedido en el
enunciado: método y, para cada dataset, tamaño medio de clique y tiempo
medio.

    \begin{tcolorbox}[breakable, size=fbox, boxrule=1pt, pad at break*=1mm,colback=cellbackground, colframe=cellborder]
\prompt{In}{incolor}{291}{\boxspacing}
\begin{Verbatim}[commandchars=\\\{\}]
\PY{n}{method\PYZus{}summaries} \PY{o}{=} \PY{p}{\PYZob{}}
    \PY{l+s+s2}{\PYZdq{}}\PY{l+s+s2}{Brute Force}\PY{l+s+s2}{\PYZdq{}}\PY{p}{:} \PY{p}{\PYZob{}}
        \PY{l+s+s2}{\PYZdq{}}\PY{l+s+s2}{IMDB}\PY{l+s+s2}{\PYZdq{}}\PY{p}{:} \PY{n}{make\PYZus{}summary\PYZus{}dict}\PY{p}{(}\PY{n}{bf\PYZus{}imdb\PYZus{}results}\PY{p}{)}\PY{p}{,}
        \PY{l+s+s2}{\PYZdq{}}\PY{l+s+s2}{COLLAB}\PY{l+s+s2}{\PYZdq{}}\PY{p}{:} \PY{n}{make\PYZus{}summary\PYZus{}dict}\PY{p}{(}\PY{n}{bf\PYZus{}collab\PYZus{}results}\PY{p}{)}\PY{p}{,}
    \PY{p}{\PYZcb{}}\PY{p}{,}
    \PY{l+s+s2}{\PYZdq{}}\PY{l+s+s2}{Replicator Dynamics}\PY{l+s+s2}{\PYZdq{}}\PY{p}{:} \PY{p}{\PYZob{}}
        \PY{l+s+s2}{\PYZdq{}}\PY{l+s+s2}{IMDB}\PY{l+s+s2}{\PYZdq{}}\PY{p}{:} \PY{n}{make\PYZus{}summary\PYZus{}dict}\PY{p}{(}\PY{n}{rd\PYZus{}imdb\PYZus{}results}\PY{p}{)}\PY{p}{,}
        \PY{l+s+s2}{\PYZdq{}}\PY{l+s+s2}{COLLAB}\PY{l+s+s2}{\PYZdq{}}\PY{p}{:} \PY{n}{make\PYZus{}summary\PYZus{}dict}\PY{p}{(}\PY{n}{rd\PYZus{}collab\PYZus{}results}\PY{p}{)}\PY{p}{,}
    \PY{p}{\PYZcb{}}\PY{p}{,}
    \PY{l+s+s2}{\PYZdq{}}\PY{l+s+s2}{GNN + RD}\PY{l+s+s2}{\PYZdq{}}\PY{p}{:} \PY{p}{\PYZob{}}
        \PY{l+s+s2}{\PYZdq{}}\PY{l+s+s2}{IMDB}\PY{l+s+s2}{\PYZdq{}}\PY{p}{:} \PY{n}{gnn\PYZus{}imdb\PYZus{}summary}\PY{p}{,}
        \PY{l+s+s2}{\PYZdq{}}\PY{l+s+s2}{COLLAB}\PY{l+s+s2}{\PYZdq{}}\PY{p}{:} \PY{n}{gnn\PYZus{}collab\PYZus{}summary}\PY{p}{,}
    \PY{p}{\PYZcb{}}\PY{p}{,}
\PY{p}{\PYZcb{}}

\PY{n}{final\PYZus{}comparison\PYZus{}df} \PY{o}{=} \PY{n}{pd}\PY{o}{.}\PY{n}{DataFrame}\PY{p}{(}\PY{p}{[}
    \PY{p}{\PYZob{}}
        \PY{l+s+s2}{\PYZdq{}}\PY{l+s+s2}{Method}\PY{l+s+s2}{\PYZdq{}}\PY{p}{:} \PY{l+s+s2}{\PYZdq{}}\PY{l+s+s2}{Brute Force}\PY{l+s+s2}{\PYZdq{}}\PY{p}{,}
        \PY{l+s+s2}{\PYZdq{}}\PY{l+s+s2}{IMDB Clique Size}\PY{l+s+s2}{\PYZdq{}}\PY{p}{:} \PY{n}{format\PYZus{}metric}\PY{p}{(}\PY{n}{method\PYZus{}summaries}\PY{p}{[}\PY{l+s+s2}{\PYZdq{}}\PY{l+s+s2}{Brute Force}\PY{l+s+s2}{\PYZdq{}}\PY{p}{]}\PY{p}{[}\PY{l+s+s2}{\PYZdq{}}\PY{l+s+s2}{IMDB}\PY{l+s+s2}{\PYZdq{}}\PY{p}{]}\PY{o}{.}\PY{n}{get}\PY{p}{(}\PY{l+s+s2}{\PYZdq{}}\PY{l+s+s2}{mean\PYZus{}clique\PYZus{}size}\PY{l+s+s2}{\PYZdq{}}\PY{p}{)}\PY{p}{)}\PY{p}{,}
        \PY{l+s+s2}{\PYZdq{}}\PY{l+s+s2}{IMDB Time}\PY{l+s+s2}{\PYZdq{}}\PY{p}{:} \PY{n}{format\PYZus{}metric}\PY{p}{(}\PY{n}{method\PYZus{}summaries}\PY{p}{[}\PY{l+s+s2}{\PYZdq{}}\PY{l+s+s2}{Brute Force}\PY{l+s+s2}{\PYZdq{}}\PY{p}{]}\PY{p}{[}\PY{l+s+s2}{\PYZdq{}}\PY{l+s+s2}{IMDB}\PY{l+s+s2}{\PYZdq{}}\PY{p}{]}\PY{o}{.}\PY{n}{get}\PY{p}{(}\PY{l+s+s2}{\PYZdq{}}\PY{l+s+s2}{mean\PYZus{}time}\PY{l+s+s2}{\PYZdq{}}\PY{p}{)}\PY{p}{,} \PY{n}{digits}\PY{o}{=}\PY{l+m+mi}{4}\PY{p}{)}\PY{p}{,}
        \PY{l+s+s2}{\PYZdq{}}\PY{l+s+s2}{COLLAB Clique Size}\PY{l+s+s2}{\PYZdq{}}\PY{p}{:} \PY{n}{format\PYZus{}metric}\PY{p}{(}\PY{n}{method\PYZus{}summaries}\PY{p}{[}\PY{l+s+s2}{\PYZdq{}}\PY{l+s+s2}{Brute Force}\PY{l+s+s2}{\PYZdq{}}\PY{p}{]}\PY{p}{[}\PY{l+s+s2}{\PYZdq{}}\PY{l+s+s2}{COLLAB}\PY{l+s+s2}{\PYZdq{}}\PY{p}{]}\PY{o}{.}\PY{n}{get}\PY{p}{(}\PY{l+s+s2}{\PYZdq{}}\PY{l+s+s2}{mean\PYZus{}clique\PYZus{}size}\PY{l+s+s2}{\PYZdq{}}\PY{p}{)}\PY{p}{)}\PY{p}{,}
        \PY{l+s+s2}{\PYZdq{}}\PY{l+s+s2}{COLLAB Time}\PY{l+s+s2}{\PYZdq{}}\PY{p}{:} \PY{n}{format\PYZus{}metric}\PY{p}{(}\PY{n}{method\PYZus{}summaries}\PY{p}{[}\PY{l+s+s2}{\PYZdq{}}\PY{l+s+s2}{Brute Force}\PY{l+s+s2}{\PYZdq{}}\PY{p}{]}\PY{p}{[}\PY{l+s+s2}{\PYZdq{}}\PY{l+s+s2}{COLLAB}\PY{l+s+s2}{\PYZdq{}}\PY{p}{]}\PY{o}{.}\PY{n}{get}\PY{p}{(}\PY{l+s+s2}{\PYZdq{}}\PY{l+s+s2}{mean\PYZus{}time}\PY{l+s+s2}{\PYZdq{}}\PY{p}{)}\PY{p}{,} \PY{n}{digits}\PY{o}{=}\PY{l+m+mi}{4}\PY{p}{)}\PY{p}{,}
    \PY{p}{\PYZcb{}}\PY{p}{,}
    \PY{p}{\PYZob{}}
        \PY{l+s+s2}{\PYZdq{}}\PY{l+s+s2}{Method}\PY{l+s+s2}{\PYZdq{}}\PY{p}{:} \PY{l+s+s2}{\PYZdq{}}\PY{l+s+s2}{Replicator Dynamics}\PY{l+s+s2}{\PYZdq{}}\PY{p}{,}
        \PY{l+s+s2}{\PYZdq{}}\PY{l+s+s2}{IMDB Clique Size}\PY{l+s+s2}{\PYZdq{}}\PY{p}{:} \PY{n}{format\PYZus{}metric}\PY{p}{(}\PY{n}{method\PYZus{}summaries}\PY{p}{[}\PY{l+s+s2}{\PYZdq{}}\PY{l+s+s2}{Replicator Dynamics}\PY{l+s+s2}{\PYZdq{}}\PY{p}{]}\PY{p}{[}\PY{l+s+s2}{\PYZdq{}}\PY{l+s+s2}{IMDB}\PY{l+s+s2}{\PYZdq{}}\PY{p}{]}\PY{o}{.}\PY{n}{get}\PY{p}{(}\PY{l+s+s2}{\PYZdq{}}\PY{l+s+s2}{mean\PYZus{}clique\PYZus{}size}\PY{l+s+s2}{\PYZdq{}}\PY{p}{)}\PY{p}{)}\PY{p}{,}
        \PY{l+s+s2}{\PYZdq{}}\PY{l+s+s2}{IMDB Time}\PY{l+s+s2}{\PYZdq{}}\PY{p}{:} \PY{n}{format\PYZus{}metric}\PY{p}{(}\PY{n}{method\PYZus{}summaries}\PY{p}{[}\PY{l+s+s2}{\PYZdq{}}\PY{l+s+s2}{Replicator Dynamics}\PY{l+s+s2}{\PYZdq{}}\PY{p}{]}\PY{p}{[}\PY{l+s+s2}{\PYZdq{}}\PY{l+s+s2}{IMDB}\PY{l+s+s2}{\PYZdq{}}\PY{p}{]}\PY{o}{.}\PY{n}{get}\PY{p}{(}\PY{l+s+s2}{\PYZdq{}}\PY{l+s+s2}{mean\PYZus{}time}\PY{l+s+s2}{\PYZdq{}}\PY{p}{)}\PY{p}{,} \PY{n}{digits}\PY{o}{=}\PY{l+m+mi}{4}\PY{p}{)}\PY{p}{,}
        \PY{l+s+s2}{\PYZdq{}}\PY{l+s+s2}{COLLAB Clique Size}\PY{l+s+s2}{\PYZdq{}}\PY{p}{:} \PY{n}{format\PYZus{}metric}\PY{p}{(}\PY{n}{method\PYZus{}summaries}\PY{p}{[}\PY{l+s+s2}{\PYZdq{}}\PY{l+s+s2}{Replicator Dynamics}\PY{l+s+s2}{\PYZdq{}}\PY{p}{]}\PY{p}{[}\PY{l+s+s2}{\PYZdq{}}\PY{l+s+s2}{COLLAB}\PY{l+s+s2}{\PYZdq{}}\PY{p}{]}\PY{o}{.}\PY{n}{get}\PY{p}{(}\PY{l+s+s2}{\PYZdq{}}\PY{l+s+s2}{mean\PYZus{}clique\PYZus{}size}\PY{l+s+s2}{\PYZdq{}}\PY{p}{)}\PY{p}{)}\PY{p}{,}
        \PY{l+s+s2}{\PYZdq{}}\PY{l+s+s2}{COLLAB Time}\PY{l+s+s2}{\PYZdq{}}\PY{p}{:} \PY{n}{format\PYZus{}metric}\PY{p}{(}\PY{n}{method\PYZus{}summaries}\PY{p}{[}\PY{l+s+s2}{\PYZdq{}}\PY{l+s+s2}{Replicator Dynamics}\PY{l+s+s2}{\PYZdq{}}\PY{p}{]}\PY{p}{[}\PY{l+s+s2}{\PYZdq{}}\PY{l+s+s2}{COLLAB}\PY{l+s+s2}{\PYZdq{}}\PY{p}{]}\PY{o}{.}\PY{n}{get}\PY{p}{(}\PY{l+s+s2}{\PYZdq{}}\PY{l+s+s2}{mean\PYZus{}time}\PY{l+s+s2}{\PYZdq{}}\PY{p}{)}\PY{p}{,} \PY{n}{digits}\PY{o}{=}\PY{l+m+mi}{4}\PY{p}{)}\PY{p}{,}
    \PY{p}{\PYZcb{}}\PY{p}{,}
    \PY{p}{\PYZob{}}
        \PY{l+s+s2}{\PYZdq{}}\PY{l+s+s2}{Method}\PY{l+s+s2}{\PYZdq{}}\PY{p}{:} \PY{l+s+s2}{\PYZdq{}}\PY{l+s+s2}{GNN + RD}\PY{l+s+s2}{\PYZdq{}}\PY{p}{,}
        \PY{l+s+s2}{\PYZdq{}}\PY{l+s+s2}{IMDB Clique Size}\PY{l+s+s2}{\PYZdq{}}\PY{p}{:} \PY{n}{format\PYZus{}metric}\PY{p}{(}\PY{n}{method\PYZus{}summaries}\PY{p}{[}\PY{l+s+s2}{\PYZdq{}}\PY{l+s+s2}{GNN + RD}\PY{l+s+s2}{\PYZdq{}}\PY{p}{]}\PY{p}{[}\PY{l+s+s2}{\PYZdq{}}\PY{l+s+s2}{IMDB}\PY{l+s+s2}{\PYZdq{}}\PY{p}{]}\PY{o}{.}\PY{n}{get}\PY{p}{(}\PY{l+s+s2}{\PYZdq{}}\PY{l+s+s2}{mean\PYZus{}clique\PYZus{}size}\PY{l+s+s2}{\PYZdq{}}\PY{p}{)}\PY{p}{)}\PY{p}{,}
        \PY{l+s+s2}{\PYZdq{}}\PY{l+s+s2}{IMDB Time}\PY{l+s+s2}{\PYZdq{}}\PY{p}{:} \PY{n}{format\PYZus{}metric}\PY{p}{(}\PY{n}{method\PYZus{}summaries}\PY{p}{[}\PY{l+s+s2}{\PYZdq{}}\PY{l+s+s2}{GNN + RD}\PY{l+s+s2}{\PYZdq{}}\PY{p}{]}\PY{p}{[}\PY{l+s+s2}{\PYZdq{}}\PY{l+s+s2}{IMDB}\PY{l+s+s2}{\PYZdq{}}\PY{p}{]}\PY{o}{.}\PY{n}{get}\PY{p}{(}\PY{l+s+s2}{\PYZdq{}}\PY{l+s+s2}{mean\PYZus{}time}\PY{l+s+s2}{\PYZdq{}}\PY{p}{)}\PY{p}{,} \PY{n}{digits}\PY{o}{=}\PY{l+m+mi}{4}\PY{p}{)}\PY{p}{,}
        \PY{l+s+s2}{\PYZdq{}}\PY{l+s+s2}{COLLAB Clique Size}\PY{l+s+s2}{\PYZdq{}}\PY{p}{:} \PY{n}{format\PYZus{}metric}\PY{p}{(}\PY{n}{method\PYZus{}summaries}\PY{p}{[}\PY{l+s+s2}{\PYZdq{}}\PY{l+s+s2}{GNN + RD}\PY{l+s+s2}{\PYZdq{}}\PY{p}{]}\PY{p}{[}\PY{l+s+s2}{\PYZdq{}}\PY{l+s+s2}{COLLAB}\PY{l+s+s2}{\PYZdq{}}\PY{p}{]}\PY{o}{.}\PY{n}{get}\PY{p}{(}\PY{l+s+s2}{\PYZdq{}}\PY{l+s+s2}{mean\PYZus{}clique\PYZus{}size}\PY{l+s+s2}{\PYZdq{}}\PY{p}{)}\PY{p}{)}\PY{p}{,}
        \PY{l+s+s2}{\PYZdq{}}\PY{l+s+s2}{COLLAB Time}\PY{l+s+s2}{\PYZdq{}}\PY{p}{:} \PY{n}{format\PYZus{}metric}\PY{p}{(}\PY{n}{method\PYZus{}summaries}\PY{p}{[}\PY{l+s+s2}{\PYZdq{}}\PY{l+s+s2}{GNN + RD}\PY{l+s+s2}{\PYZdq{}}\PY{p}{]}\PY{p}{[}\PY{l+s+s2}{\PYZdq{}}\PY{l+s+s2}{COLLAB}\PY{l+s+s2}{\PYZdq{}}\PY{p}{]}\PY{o}{.}\PY{n}{get}\PY{p}{(}\PY{l+s+s2}{\PYZdq{}}\PY{l+s+s2}{mean\PYZus{}time}\PY{l+s+s2}{\PYZdq{}}\PY{p}{)}\PY{p}{,} \PY{n}{digits}\PY{o}{=}\PY{l+m+mi}{4}\PY{p}{)}\PY{p}{,}
    \PY{p}{\PYZcb{}}\PY{p}{,}
\PY{p}{]}\PY{p}{)}
\PY{n}{display}\PY{p}{(}\PY{n}{final\PYZus{}comparison\PYZus{}df}\PY{p}{)}
\end{Verbatim}
\end{tcolorbox}

    
    \begin{Verbatim}[commandchars=\\\{\}]
                Method IMDB Clique Size IMDB Time COLLAB Clique Size  \textbackslash{}
0          Brute Force           11.733    0.0758                N/D   
1  Replicator Dynamics           10.380    0.0044             27.300   
2             GNN + RD            9.521    0.0115             28.042   

  COLLAB Time  
0         N/D  
1      0.0057  
2      0.0115  
    \end{Verbatim}

    
    \subsubsection{9.2 Informe comparativo
final}\label{informe-comparativo-final}

La tabla final resume con claridad tres regímenes muy distintos de
resolución. \textbf{Brute Force} ofrece la referencia exacta, pero en la
práctica ha quedado restringido al subconjunto pequeño de IMDB donde los
grafos no superaban el umbral fijado. Por eso aparece con una clique
media alta en IMDB (\texttt{11.733}) y sin resultados utilizables en
COLLAB. Esa ausencia no es un problema de implementación, sino una
conclusión experimental en sí misma: el coste combinatorio crece tan
deprisa que el método exacto deja de ser operativo en cuanto el tamaño
de los grafos aumenta un poco.

\textbf{Replicator Dynamics} ha sido el método más equilibrado entre
cobertura, validez y coste. En IMDB ha obtenido una clique media de
\texttt{10.380} con un tiempo medio de \texttt{0.0033}, y en COLLAB una
media de \texttt{27.300} con \texttt{0.0044}. Lo importante aquí no es
solo que los tiempos sean bajos, sino que las soluciones hayan sido
cliques válidas y, además, que en los grafos pequeños de IMDB haya
existido una coincidencia total con la referencia exacta. Eso convierte
a RD en el verdadero punto de apoyo de toda la práctica: ofrece una
alternativa mucho más escalable sin perder una conexión fuerte con la
solución combinatoria correcta.

\textbf{GNN + RD} ha introducido una segunda capa de sofisticación. En
los subconjuntos hold-out evaluados, el método híbrido ha producido una
mejora ligera en COLLAB (\texttt{27.958} frente a \texttt{27.300} si se
compara con el resumen global de RD, y \texttt{27.9583} frente a
\texttt{27.8958} si se compara con RD puro sobre el hold-out
correspondiente), mientras que en IMDB el comportamiento ha sido más
delicado: frente al resumen global de RD, la clique media del híbrido
(\texttt{9.521}) queda por debajo, aunque en la comparación estricta
hold-out sí mejora levemente a RD puro (\texttt{9.5208} frente a
\texttt{9.4375}). Esa diferencia de lectura es importante porque
recuerda que no siempre se están comparando exactamente los mismos
subconjuntos de grafos.

Desde el punto de vista del coste, el híbrido ha sido claramente más
caro: alrededor de \texttt{0.0105} en IMDB y \texttt{0.0108} en COLLAB.
Por tanto, el valor añadido de la GNN no ha estado en acelerar el
proceso, sino en intentar mejorar la calidad de la inicialización. En
este notebook esa mejora ha sido real pero pequeña. La conclusión
comparativa más razonable es que el salto decisivo está en pasar de
fuerza bruta a RD, mientras que el paso de RD a GNN + RD funciona más
como un refinamiento dependiente del dataset y de la calidad de las
features que como una mejora garantizada y masiva.

    \subsection{10. Conclusiones}\label{conclusiones}

A lo largo de la práctica se ha recorrido una progresión muy clara:
partir de un método exacto pero inviable a gran escala, reformular el
problema en un espacio continuo y, finalmente, estudiar si un modelo
aprendido puede mejorar todavía más la búsqueda. Esa secuencia ha
permitido entender no solo cómo se resuelve el problema de clique máximo
desde distintos enfoques, sino también qué se gana y qué se pierde en
cada paso.

La primera conclusión importante la ha proporcionado la \textbf{fuerza
bruta}. Su papel ha sido doble. Por un lado, ha servido para validar la
implementación y para fijar un ground truth exacto en grafos pequeños.
Por otro, ha mostrado de forma experimental por qué ese enfoque no puede
escalar: ya en los primeros grafos de IMDB han aparecido omisiones por
tamaño y en COLLAB la estrategia exacta ha quedado completamente
bloqueada en el subconjunto analizado. Es decir, el primer ejercicio no
solo ha sido una comprobación técnica, sino una justificación práctica
de por qué hacen falta métodos más inteligentes.

La segunda gran conclusión ha llegado con la \textbf{formulación de
Motzkin-Straus}. En los ejemplos pequeños se ha comprobado que el tamaño
de la clique máxima puede recuperarse a través de una función cuadrática
sobre el simplex, y que esa correspondencia se verifica numéricamente
con precisión prácticamente exacta. Sin embargo, también se ha visto que
la relajación continua original no basta siempre por sí sola, porque
puede asignar el mismo valor a soluciones verdaderas y a soluciones
espurias. El grafo cherry ha sido especialmente útil para entender esto:
la teoría no solo dice que pueden existir esos empates, sino que en la
práctica aparecen, y la regularización es precisamente lo que permite
romperlos en favor de las soluciones más concentradas y estructuralmente
correctas.

La tercera conclusión ha sido que \textbf{Replicator Dynamics}
constituye el núcleo práctico más sólido de toda la práctica. RD ha
mantenido la solución dentro del simplex, ha producido cliques válidas
en todos los grafos evaluados y, además, ha reproducido exactamente el
tamaño de la solución óptima en todos los casos pequeños de IMDB donde
la fuerza bruta sí pudo aplicarse. Esa combinación de validez, bajo
coste y coherencia con el ground truth convierte a RD en una herramienta
muy potente para este problema. También se ha visto que su
comportamiento depende de la estructura del grafo: cuando la mejor
clique está muy marcada, la convergencia es rápida; cuando existen
varias regiones densas en competencia, el proceso requiere más
iteraciones. Esa observación añade una interpretación geométrica y
estructural al simple dato numérico.

La cuarta conclusión ha estado en el método \textbf{GNN + RD}. El modelo
híbrido se entrenó correctamente, la loss descendió de forma estable y
las soluciones generadas siguieron siendo cliques válidas tras el
refinamiento con RD. Sin embargo, la magnitud de la mejora respecto a RD
puro ha sido contenida. En algunos escenarios la inicialización
aprendida ha aportado una ventaja pequeña, mientras que en otros la
diferencia ha sido casi inapreciable. Esto sugiere que el valor del
híbrido depende en gran medida de la calidad de las features, de la
arquitectura usada y de cuánto patrón repetible exista en el dataset
para que la GNN realmente aprenda una priorización útil de nodos.

La ablación final ha reforzado esa idea. Variar el número de iteraciones
de RD en entrenamiento y en inferencia apenas ha cambiado la calidad
media de la solución en el subconjunto analizado, aunque sí ha
modificado el tiempo de ejecución. Eso apunta a que, en esta
configuración concreta, el cuello de botella no ha estado tanto en el
número exacto de pasos de refinamiento como en la calidad del punto de
partida. Dicho de otra forma, una vez que RD dispone de una
inicialización razonable, añadir más iteraciones no siempre se traduce
en una mejora apreciable.

En conjunto, la práctica deja una lectura bastante completa del
problema. La formulación continua funciona, la regularización es
necesaria, RD ofrece una aproximación muy eficaz y el híbrido con GNN
abre una línea interesante pero más sensible a la calidad del diseño
experimental. Si hubiera que resumir el aprendizaje principal en una
sola idea, sería esta: el mayor salto práctico no se ha producido al
añadir aprendizaje profundo, sino al pasar de la búsqueda exhaustiva a
una relajación continua bien estructurada y optimizable. A partir de
ahí, la GNN puede actuar como refinamiento adicional, pero su beneficio
real depende de que el problema y las features ofrezcan suficiente
información como para aprender una inicialización mejor que la uniforme.


    % Add a bibliography block to the postdoc
    
    
    
\end{document}
